\documentclass[11pt]{article}

\usepackage[T1]{fontenc}
\usepackage[utf8]{inputenc}
\usepackage{lmodern}
\usepackage{microtype}
\usepackage[letterpaper,margin=1in]{geometry}
\usepackage{amsmath,amssymb}
\usepackage{array,longtable,booktabs}
\usepackage[dvipsnames]{xcolor}
\usepackage{titlesec}
\usepackage{xurl}
\usepackage[hidelinks,bookmarksnumbered]{hyperref}

\setlength{\parindent}{0pt}
\setlength{\parskip}{0.6em}
\definecolor{heading}{HTML}{1F3864}
\titleformat{\section}{\Large\bfseries\color{heading}}{}{0pt}{}
\titleformat{\subsection}{\large\bfseries\color{heading}}{}{0pt}{}
\titleformat{\subsubsection}{\normalsize\bfseries}{}{0pt}{}
\titlespacing*{\subsubsection}{0pt}{1.1em}{0.2em}

% Reference entry with a 0.5in hanging indent (APA style).
\newenvironment{hangref}{\par\hangindent=0.5in\hangafter=1\noindent}{\par}
% Question heading inside a paper or survey.
\newcommand{\question}[1]{\subsubsection*{#1}}
% Unnumbered heading that still appears in the PDF bookmarks.
\newcommand{\bookmarked}[3]{#1*{#3}\phantomsection\addcontentsline{toc}{#2}{#3}}
\newcolumntype{P}[1]{>{\raggedright\arraybackslash}p{#1}}

\title{\color{heading}Literature Review: Offline Reinforcement Learning\\[0.3em]\large Parts A and B}
\author{[Student name]\\{}[Course]}
\date{[Date]}

\begin{document}
\maketitle


\bookmarked{\section}{section}{Research Problem}

Offline reinforcement learning (RL) trains a decision-making policy entirely from a fixed dataset of previously logged interactions, with no further access to the environment. Because the learned policy can prefer actions that rarely or never appear in the data, the critic's value estimates for those actions are unreliable and tend to be overly optimistic, which pulls the policy toward poor behavior. Most offline RL methods counter this with a conservatism or behavior-regularization term whose strength is controlled by a single global coefficient that must be tuned for each dataset.

My research problem is how to set the strength of that regularization in a calibrated, per-sample way rather than with one hand-tuned constant. The approach I am developing, Bayesian Conformal Aggregation (BCA), estimates the uncertainty of the critic's Bellman residuals, calibrates it with weighted Bayesian conformal prediction, and uses the resulting width to adjust the regularization term inside existing algorithms.

The fifteen works reviewed below fall into three groups. The first group covers offline RL algorithms and their evaluation: CQL, TD3+BC, IQL and ReBRAC are the four host algorithms that BCA modifies; Unifloral supplies the unified implementation and evaluation protocol used to compare them; and A-Crab, Quantile Q-Learning, UNIQ and Bayesian Flow Networks offer other ways to stay conservative or to plan from offline data. The second group concerns calibrated uncertainty: conformal prediction for policy evaluation, Bellman calibration for value learning, and weighted Bayesian conformal prediction, the procedure BCA uses to set its threshold. The third group, reviewed in its own section after the research papers, consists of three surveys of offline RL that place these papers in the wider field.

\clearpage

\bookmarked{\section}{section}{Part A: Research Papers}

This section answers the eleven research-paper questions for each of the twelve research papers.

\bookmarked{\subsection}{subsection}{Paper 1: Conservative Q-Learning for Offline Reinforcement Learning (CQL)}

\question{1. What is the citation?}

\begin{hangref}Kumar, A., Zhou, A., Tucker, G., \& Levine, S. (2020). Conservative Q-learning for offline reinforcement learning. \emph{Advances in Neural Information Processing Systems}, \emph{33}, 1179--1191. \url{https://proceedings.neurips.cc/paper/2020/hash/0d2b2061826a5df3221116a5085a6052-Abstract.html}\end{hangref}
\begin{hangref}Version reviewed: Kumar, A., Zhou, A., Tucker, G., \& Levine, S. (2020). \emph{Conservative Q-learning for offline reinforcement learning} (arXiv:2006.04779v3) [Preprint]. arXiv. \url{https://doi.org/10.48550/arXiv.2006.04779}\end{hangref}

\question{2. What is the research area?}

The paper belongs to machine learning and, within deep RL, to offline (batch) RL, where a policy is learned from a fixed dataset of logged transitions without further interaction. CQL is model-free and can be instantiated as an actor-critic method (on soft actor-critic, for continuous control) or a value-based method (on QR-DQN, for discrete Atari control); its conservative evaluation step can also serve alone for off-policy evaluation. The problem family is overestimation of Q-values at actions missing from the dataset. Among prior approaches, CQL represents the value-regularization, or pessimistic critic, family. BEAR and both BRAC variants penalize divergence from an explicitly estimated behavior policy (BRAC-v also inside the critic target), and other methods use ensemble or uncertainty penalties; CQL instead penalizes the Q-values themselves and needs no behavior model.

\question{3. What is the research question?}

Can an offline RL agent avoid the overestimation caused by distribution shift by learning a Q-function whose expected value under the learned policy provably lower-bounds the policy's true value, and does optimizing a policy against such a conservative critic produce better policies than existing policy-constraint methods?

The authors break this into several sub-questions. Theoretically, they ask whether pushing Q-values down under a chosen action distribution while pushing them up on dataset actions yields a bound tighter than a pointwise one, whether that bound survives sampling error and function approximation, and whether the resulting policy improves on the behavior policy with high probability. Empirically, they test whether the learned values really fall below true returns, whether the method beats baselines on data that mix several policies, come from human demonstrations or consist of image-based Atari replay, and which design choices (variant, maximization term, automatic tuning of \ensuremath{\alpha}) matter.

\question{4. What is the significance of the question?}

The question matters in practice because the domains where RL would be most useful, such as robotics, recommendation and medicine, make online trial and error expensive or hazardous, while logged data are often abundant. Reliable learning from such logs would let RL reuse large datasets much as supervised learning does. The authors stress that this promise had not been realized, because standard off-policy algorithms applied to static data usually fail.

Technically, the failure comes from bootstrapping. The Bellman target evaluates the critic at actions proposed by the learned policy, which may lie outside the data; errors there are never corrected by new experience, and because the policy maximizes the critic, it is drawn toward exactly the inflated values. Earlier remedies either constrained the policy toward an estimated behavior model, which is hard when data come from several sources, or used uncertainty estimates that proved too unreliable offline. When the paper appeared, the D4RL benchmark had just shown that existing methods often could not beat behavior cloning on multi-modal or human-collected data, so a simpler and more robust method with formal guarantees addressed a clear gap in the field.

\question{5. What are the contributions reported in this paper?}

The central contribution is the CQL framework, a family of critic objectives indexed by a regularizer R(\ensuremath{\mu}) on the distribution whose Q-values are pushed down. The critic is trained with the usual Bellman error plus a penalty that lowers Q-values under \ensuremath{\mu}, chosen to favor high-valued actions, and raises them on dataset actions. Two members are derived in closed form: CQL(H), whose penalty is a log-sum-exp (soft maximum) of Q over actions, and CQL(\ensuremath{\rho}), which replaces the soft maximum with an exponentially weighted average of Q-values over actions sampled from the previous policy. An appendix adds a variance-penalizing variant, CQL(var), linked to distributionally robust optimization.

The theoretical contributions are a chain of results: the simpler penalty yields a pointwise lower bound on Q; the full penalty yields a tighter lower bound on the policy's expected value; both extend to linear and neural-tangent-kernel approximation; the iterates stay conservative when the policy changes slowly; the update widens the value gap between in-data and unseen actions; and the induced policy objective is a penalized empirical return with a high-probability safe-improvement guarantee over the behavior policy. A further appendix result shows the maximization term must use the behavior distribution for the bound to hold.

On the practical side, the paper provides an algorithm that adds about twenty lines to SAC or QR-DQN, needs no behavior-policy model, and can tune \ensuremath{\alpha} automatically through a Lagrangian threshold. Empirically, it contributes a broad evaluation on D4RL and offline Atari, direct measurements of value underestimation and gap expansion, and ablations of the main design choices.

\question{6. What is the research method?}

The research method combines algorithm design, theoretical analysis (mostly tabular, with extensions to linear and NTK approximation), and empirical benchmarking supported by diagnostic measurements and ablations.

The approach modifies only the critic objective. For the CQL(H) variant used in most experiments, the critic minimizes \begin{equation*}\alpha\,\mathbb{E}_{s\sim\mathcal{D}}\Big[\log\sum_{a}\exp Q(s,a)-\mathbb{E}_{a\sim\hat{\pi}_{\beta}}Q(s,a)\Big]+\frac{1}{2}\,\mathbb{E}_{(s,a,s')\sim\mathcal{D}}\Big[\big(Q(s,a)-\hat{\mathcal{B}}^{\pi}\hat{Q}^{k}(s,a)\big)^{2}\Big].\end{equation*} The log-sum-exp term pushes down whatever actions the critic currently rates highly, while the second term props up actions that appear in the data, so the net effect falls on unseen actions. In continuous action spaces the log-sum-exp is estimated by importance sampling with ten uniform and ten policy actions per state; in discrete domains it is computed exactly. The actor is a SAC-style entropy-regularized policy maximizing $\mathbb{E}[Q(s,a)-\log\pi(a\mid s)]$, and no behavior policy is fitted.

The trade-off coefficient \ensuremath{\alpha} is the key hyperparameter. For continuous control, dual gradient descent raises \ensuremath{\alpha} when the gap between soft-maximum and dataset Q-values exceeds a threshold \ensuremath{\tau} and lets it fall toward zero when the gap is below \ensuremath{\tau} (\ensuremath{\tau} = 10 for Gym-MuJoCo, \ensuremath{\tau} = 5 for Adroit and Kitchen); for Atari, \ensuremath{\alpha} is fixed at 0.5, 1.0 or 4.0 depending on the dataset. Guided by Theorem 3.3, the policy learning rate is reduced to 3e-5 while the critic keeps 3e-4; other settings, including twin critics and soft target updates, follow SAC, and D4RL runs last one million gradient steps.

The experiments cover the D4RL Gym-MuJoCo locomotion datasets (random, medium, expert, medium-expert, random-expert and mixed for three robots), Adroit hand manipulation from human and cloned data, AntMaze and Franka Kitchen, against BEAR, BRAC-p, BRAC-v, offline SAC and behavior cloning. Gym-MuJoCo baseline scores come from the D4RL paper; no source is given for the others. A QR-DQN version, with QR-DQN's released hyperparameters left untuned, is compared with REM and QR-DQN on Atari games using two slices of the DQN-replay data: the first 20\% of an online DQN agent's samples, and subsets holding 1\% or 10\% of all its samples.

\question{7. What are the results?}

On single-policy Gym-MuJoCo datasets, CQL(H) is usually close to the strongest baseline, losing narrowly on halfcheetah-expert (104.8 against BEAR's 108.2) and winning clearly on walker2d-expert (153.9 against BC's 125.7). The large gains come on most mixed datasets: 92.5, 91.1 and 110.5 on the three random-expert tasks, where the best baselines reach 53.0, 2.7 and 11.1, and 98.7 on walker2d-medium-expert against BC's 11.3, though BC edges it on hopper-medium-expert (111.9 against 111.0).

On AntMaze, CQL(H) reaches 61.2 and 53.7 on the medium mazes, where the best baseline manages 8.0, and on the large mazes only the two CQL variants score above zero (CQL(H) 15.8 and 14.9). On Adroit, CQL(\ensuremath{\rho}) reaches 55.8 on pen-human against BC's 34.4, but BC wins on pen-cloned (56.9 against about 40). On Kitchen, CQL(H) beats BC on all three tasks (43.8 to 51.0). On Atari with the first 20\% of a DQN run, CQL generally matches or beats both baselines and stays stable on three of the four games; with 1\% of the replay data it reaches 14012.0 on Q*bert and 61.1 on Breakout, about 36 and 6 times the best baseline, but with 10\% it trails REM on Seaquest and QR-DQN on Asterix.

CQL(H)'s predicted value lies below the true return on the three hopper datasets tested (\ensuremath{-}43.20, \ensuremath{-}10.93 and \ensuremath{-}7.48), the pointwise variant is far more pessimistic, and ensembles and BEAR overestimate. The sampled gap statistic stays negative for CQL but rises above zero for BEAR on hopper-expert, where BEAR's return later deteriorates. In the ablations, dropping the dataset maximization term lowers hopper-medium return from 1866.1 to 1028.4, CQL(H) beats CQL(\ensuremath{\rho}) on three MuJoCo tasks, and automatic tuning of \ensuremath{\alpha} helps most on AntMaze.

\question{8. What are the impacts and limitations of the contributions?}

CQL's chief legacy is a simple way to make an offline critic pessimistic without modeling the behavior policy, usable with continuous and discrete actions and backed by lower-bound and safe-improvement results. It became a standard comparison point, appearing as a baseline in TD3+BC and IQL.

In their discussion the authors name two open problems: the guarantees cover tabular, linear and NTK-style approximation rather than practical deep networks, and offline RL lacks a validation signal for early stopping. When setting up the problem they also note that the policy may still face unfamiliar states at test time, a shift the action-focused penalty does not address. Critical reading adds more. The bounds require an \ensuremath{\alpha} large enough relative to unknown constants, so in practice \ensuremath{\alpha} is governed by a threshold \ensuremath{\tau} that had to change across domains; \ensuremath{\tau} = 2 drove \ensuremath{\alpha} into the millions, while \ensuremath{\tau} = 10 sometimes let Q-values diverge on the human-demonstration datasets. The better variant (CQL(H) or CQL(\ensuremath{\rho})) and the better backup also depend on the domain, the sampled log-sum-exp adds action evaluations to every update, and the gains on single-policy datasets are usually small.

The theory holds a hint for BCA: the \ensuremath{\alpha} needed for a bound shrinks as a state-action pair's count grows, and Appendix E sketches a count-dependent penalty, yet the implementation applies one global \ensuremath{\alpha} to every transition. Appendix E also attributes the weakness of earlier uncertainty-based methods to poorly calibrated uncertainty, which is the gap that calibrated per-sample weighting of the conservative term is meant to close.

\question{9. How were the results evaluated or validated?}

Validation has three strands. Theoretical claims rest on appendix proofs, mostly for the tabular case. Performance claims are tested on D4RL with its normalized return metric, reporting a smoothed average of undiscounted return over four seeds, and on Atari with the datasets and protocol of Agarwal et al. The lower-bound claim is checked directly by comparing the critic's predicted policy value with the policy's actual discounted return, and the gap-expansion claim by tracking a sampled advantage statistic over training against BEAR. Appendix G ablations isolate the variant, the maximization term and the method of setting \ensuremath{\alpha}.

The validation is broad but statistically thin. Four seeds is a small sample and the tables give only means; the Atari curves carry unexplained shaded bands, and neither the Atari seed count nor the number of evaluation episodes is stated. Gym-MuJoCo baselines are copied from the D4RL paper rather than rerun, and the other baselines have no stated source. Dataset names carry the original -v0 suffix, so scores are not directly comparable with later papers using revised D4RL versions. Some entries look inconsistent: the walker2d-mixed baseline row duplicates walker2d-medium-expert, and hopper-mixed differs between Tables 6 and 7. The authors say \ensuremath{\tau} was chosen from predicted Q-values rather than online returns, but policy learning rates were compared by returns across tasks, implying some environment evaluation.

\question{10. What are the ethical implications? How are they addressed?}

CQL's arXiv v3 (labeled a preprint under review) has no ethics or broader-impact section. NeurIPS 2020 required authors to include such a statement, and a published review of the submission refers to one (remarking that its first paragraph mostly summarized the method), but that statement is absent from the arXiv text, where the closest material is the introduction's point that gathering data by acting in the real world is expensive and can be risky, a discussion naming missing deep-network theory and overfitting as open problems, and an acknowledgment of DARPA Assured Autonomy funding.

Several ethical issues follow. Offline RL appeals most in healthcare, robotics and recommendation, where a policy that looks good offline but fails in use can hurt people. CQL's safe-improvement result is only relative to the behavior policy, probabilistic, and based on assumptions deep networks do not meet, so calling it safe could invite misplaced trust. Because CQL's penalty keeps the policy close to dataset actions, any bias in the logging policy carries over to the learned one. All evidence comes from simulators and Atari, and without an offline stopping rule, performance is read from online returns a real deployment could not obtain beforehand. Compute costs go unreported.

CQL's response to these issues is indirect: conservatism is a technical step toward safer offline learning, but the arXiv text never discusses data bias, deployment safeguards or societal consequences.

\question{11. What themes are suggested by this research? (minimum 3)}

\textbf{Distribution shift and out-of-distribution value overestimation.} CQL is built on the observation that Bellman targets query the critic at actions missing from the data, where errors cannot be corrected offline and a maximizing policy exploits them. The paper documents the problem concretely, showing ensembles and BEAR overestimating policy values, ensembles by many orders of magnitude, and responds by suppressing out-of-distribution values inside the critic rather than relying on the actor to avoid them.

\textbf{Conservatism versus policy improvement trade-off.} The paper treats pessimism as something to be dosed. The pointwise penalty proves needlessly pessimistic, so a dataset maximization term is added to tighten the bound to the policy's expected value, which also raises returns. The coefficient \ensuremath{\alpha} controls the balance, and no single setting suits every task: a fixed \ensuremath{\alpha} of 5 reached only 0.21 on antmaze-medium-diverse, against 0.53 when \ensuremath{\alpha} was tuned automatically against a threshold.

\textbf{Behavior regularization and policy constraints.} Although CQL is presented as an alternative to policy constraints and needs no behavior model, Theorem 3.5 shows it implicitly maximizes empirical return minus a divergence between the learned and empirical behavior policies. This puts critic-side value regularization and actor-side behavior regularization on a common footing, which helps when comparing CQL with TD3+BC, which adds a behavior-cloning term only to the actor, and ReBRAC, which penalizes both actor and critic target with separate coefficients.

\textbf{Theory versus practice.} The guarantees assume tabular or linearized critics, concentration bounds with unknown constants, slowly changing policies and a sufficiently large \ensuremath{\alpha}, whereas the experiments use deep networks, importance-sampled log-sum-exp estimates and a heuristically thresholded \ensuremath{\alpha}. The authors partly bridge this gap with direct empirical checks of underestimation and gap expansion, and they explicitly leave a rigorous analysis for deep networks to future work.

\textbf{Dataset quality and composition.} CQL's advantage depends strongly on how the data were collected. On single-policy datasets it is usually close to the best baseline, with a few clear wins such as walker2d-expert. Its largest margins come on most multi-policy datasets, AntMaze, pen-human and the 1\% Atari subsets, yet BC still edges it on hopper-medium-expert and pen-cloned. The authors argue that heterogeneous data, which also make behavior-policy estimation hard, are the realistic case.

\clearpage

\bookmarked{\subsection}{subsection}{Paper 2: A Minimalist Approach to Offline Reinforcement Learning (TD3+BC)}

\question{1. What is the citation?}

\begin{hangref}Fujimoto, S., \& Gu, S. S. (2021). A minimalist approach to offline reinforcement learning. \emph{Advances in Neural Information Processing Systems}, \emph{34}, 20132--20145. \url{https://proceedings.neurips.cc/paper/2021/hash/a8166da05c5a094f7dc03724b41886e5-Abstract.html}\end{hangref}
\begin{hangref}Version reviewed: Fujimoto, S., \& Gu, S. S. (2021). \emph{A minimalist approach to offline reinforcement learning} (arXiv:2106.06860v2) [Preprint]. arXiv. \url{https://doi.org/10.48550/arXiv.2106.06860}\end{hangref}

\question{2. What is the research area?}

The paper contributes to machine learning through deep reinforcement learning, specifically its offline (batch) branch, in which training draws only on a fixed dataset. Technically, TD3+BC is a model-free continuous-control algorithm built on TD3, an actor-critic method whose deterministic actor follows the deterministic policy gradient. It targets extrapolation error: critics trained on limited data misjudge actions absent from the dataset, and the actor exploits those mistakes. Among existing remedies, the work sits in the policy-constraint or behavior-regularization family (with BCQ and BRAC), as opposed to value regularization (CQL), critic regularization that needs a learned behavior model (Fisher-BRC), or advantage-weighted imitation (AWAC). It borrows its penalty from earlier demonstration learning that added behavior cloning to DDPG or SAC, and joins meta-analyses of how implementation details shape reported results.

\question{3. What is the research question?}

Can an existing online deep RL algorithm reach state-of-the-art offline performance while being changed as little as possible, and how much of the extra machinery in recent offline methods is actually needed?

The authors split this into narrower questions. First, do the code-level adjustments that CQL and Fisher-BRC make to their SAC base matter for their scores, and what do they cost in training time? Second, is a behavior-cloning penalty in TD3's actor update, together with standardized input states, enough to match those methods on the D4RL locomotion tasks? Third, how much does each ingredient (the cloning term, the normalization and the value of \ensuremath{\alpha}) contribute? Finally, the paper raises a side question that it does not resolve: are policies produced by offline training stable across episodes and across late checkpoints?

\question{4. What is the significance of the question?}

The first stake is practical. Offline RL is attractive where gathering new interaction is expensive, slow or dangerous, as with robots or patients, so a method that only works after many trial runs in the environment undercuts its own purpose. Every extra hyperparameter, pre-trained generative model, architectural change or actor pre-training phase is another choice that must be validated, and a truly offline user has no cheap way to validate it. Fewer moving parts mean fewer decisions that silently depend on online testing.

The central technical obstacle is extrapolation error, in which critics overestimate unseen actions and the policy drifts toward them. Many fixes existed, but each arrived bundled with engineering tweaks, so it was unclear whether gains came from the stated idea or from the tweaks. A minimal method isolates the effect of the constraint itself. At the time, the field was producing increasingly elaborate algorithms, and the authors argued that simple options had been neglected. If a few lines of code could match the strongest existing methods, that would raise the bar that new, more elaborate algorithms must clear, and a cheap method matters more as offline datasets grow.

\question{5. What are the contributions reported in this paper?}

The paper offers three kinds of contribution. The first is diagnostic. It catalogs the algorithmic and implementation-level changes that CQL and Fisher-BRC make to SAC (Table 1), such as altered network sizes, actor pre-training, a constant reward bonus and removal of the entropy term from the target, and it measures what happens when those details are stripped away. It also documents a problem the authors consider neglected: policies trained offline can do far worse than their average on single episodes of one evaluation and on single evaluations among the last ten, unlike a policy trained online with TD3.

The second is algorithmic. TD3+BC changes TD3 in two places only: the actor loss gains a squared-error behavior-cloning penalty, and states are standardized with dataset statistics. Scaling the trade-off weight by the average magnitude of Q leaves a single hyperparameter, \ensuremath{\alpha}, held fixed across all tasks. Appendix B adds rules of thumb for setting \ensuremath{\lambda} by hand.

The third is empirical and practical. The authors rerun CQL and Fisher-BRC from the original code under one shared protocol, time CQL, both Fisher-BRC implementations and TD3+BC (and, in the appendix, the Decision Transformer) on the same GPU and CPU, reimplementing Fisher-BRC in PyTorch so frameworks match, ablate each component, and add experiments on mixed random and expert data, AntMaze, state normalization inside other algorithms, and a comparison with the Decision Transformer. The full code is public.

\question{6. What is the research method?}

Methodologically, the paper mixes a critical analysis of existing algorithms, minimal algorithm design and empirical benchmarking with ablations; it has no theoretical results. For the analysis, the authors rerun CQL and Fisher-BRC without their implementation adjustments and report the percent change in score (Figure 1). Separately, for BC, CQL, Fisher-BRC, TD3+BC and online TD3, they measure how far the worst of the ten episodes at the final evaluation, and the worst of the last ten evaluations, falls below average (Figures 2 and 3).

The algorithm starts from TD3, an off-policy actor-critic method whose deterministic actor normally maximizes Q(s, \ensuremath{\pi}(s)). TD3+BC instead maximizes \begin{equation*}\mathbb{E}_{(s,a)\sim\mathcal{D}}\big[\lambda Q(s,\pi(s))-(\pi(s)-a)^{2}\big],\end{equation*} so the actor balances estimated value against squared distance from the logged action. Because that penalty is bounded (at most 4 when actions lie in $[-1,1]$) while Q grows with the reward scale, \ensuremath{\lambda} is set to \ensuremath{\alpha} divided by the mean absolute Q-value, \begin{equation*}\lambda=\frac{\alpha}{\frac{1}{N}\sum_{i}\lvert Q(s_i,a_i)\rvert}.\end{equation*} The mean is computed per mini-batch and not differentiated through, and \ensuremath{\alpha} = 2.5 throughout. The second change standardizes every state feature by its dataset mean and standard deviation, with $\epsilon=10^{-3}$ added to the denominator. All TD3 settings stay at their defaults: Adam with learning rate $3\times10^{-4}$ for actor and critic, batch size 256, discount 0.99, target update rate $5\times10^{-3}$, target policy noise 0.2 clipped to $\pm 0.5$, actor updates every second step, and two hidden layers of 256 ReLU units.

Experiments use the D4RL v0 MuJoCo datasets for HalfCheetah, Hopper and Walker2d at five quality levels (random, medium, medium-replay, medium-expert and expert). Each method trains for one million steps and is evaluated every 5,000 steps. Baselines are behavior cloning (the authors' own code), CQL and Fisher-BRC (rerun from author code) and BRAC-p and AWAC (published scores). Ablations drop the BC term, the normalization or the value term, and sweep \ensuremath{\alpha} over 1, 2, 3 and 4. Run time is measured without evaluation on one GTX 1080 GPU. The appendix adds a 50/50 random-expert mixture, AntMaze (without normalization) and the D4RL v2 datasets for a Decision Transformer comparison.

\question{7. What are the results?}

Summed over the fifteen D4RL v0 datasets, TD3+BC scores 979.3 (\ensuremath{\pm}33.4 across seeds), effectively tied with Fisher-BRC at 974.6 (\ensuremath{\pm}108.3) and well ahead of CQL (764.3), behavior cloning (595.3) and BRAC-p's published total (284.1). Its largest margins over Fisher-BRC are on hopper-medium-expert (112.2 versus 90.6) and walker2d-expert (105.7 versus 79.9), with far smaller seed variance, although CQL comes close on both (111.4 and 103.8). It loses on halfcheetah-random (10.2 against 32.2 for Fisher-BRC and 21.7 for CQL) and walker2d-medium-replay (25.2 against 42.6 for Fisher-BRC). One million training steps took 39 minutes, only seconds more than plain TD3, against 2 h 8 min for the PyTorch Fisher-BRC, 2 h 12 min for the original Fisher-BRC code and 4 h 11 min for CQL.

The ablations show both halves are needed: TD3 without the cloning term performs poorly except on some random datasets, and the full method never falls below plain BC, even on expert data. Normalization has the smallest effect but helps on several tasks. Setting \ensuremath{\alpha} to 2 or 3 gives practically the same results as the default 2.5, whereas \ensuremath{\alpha} = 1 or 4 degrades a few tasks. Removing the implementation tweaks of CQL and Fisher-BRC causes large drops on many datasets, and every offline method, TD3+BC included, shows much larger worst-case drops than online TD3. In the appendix, TD3+BC fails on the Walker2d random-expert mixture (28.8, near Fisher-BRC's 32.0 but far below CQL's 100.3), reaches 78.6 and 71.4 on the two U-maze datasets of AntMaze but at most 10.6 on the medium and large mazes, and roughly matches the Decision Transformer on the nine v2 tasks DT reports (677.4 versus 672.6) while training faster.

\question{8. What are the impacts and limitations of the contributions?}

TD3+BC's lasting effect on the field is a strong, cheap and transparent baseline. Because TD3+BC alters a few lines, keeps TD3's defaults and trains in roughly a sixth of CQL's time, it is easy to adopt with modest compute and makes gains in later methods easier to attribute to their stated ideas; the ReBRAC authors (Tarasov et al., 2023) later describe it as a de facto standard baseline. The paper shows that earlier state-of-the-art scores leaned on unexplained engineering choices, and it urges benchmarks to report the variability of final policies, not just means.

The limitations are substantial. The method is weak wherever cloning the data hurts, as on halfcheetah-random, the Walker2d random-expert mixture and the medium and large AntMaze mazes, which plausibly reflects that a uniform pull toward logged actions cannot separate good transitions from bad ones (for that mixture, the authors instead hypothesize weak cloning on the walker2d-expert data). There is no theory, and a squared-error penalty on a deterministic actor pulls each action toward the average logged action, which can fall between modes when the data mix very different behaviors. The instability the authors diagnose is left open, and all tests are simulated. The detail that matters most for BCA is that \ensuremath{\lambda} rescales Q by a batch-wide mean, so every transition in a batch receives the same constraint strength regardless of how well the data cover it. Calibrated per-sample weighting would let that strength follow coverage instead, and the paper's own advice to lower \ensuremath{\lambda} when value estimates diverge hints at the kind of value-error signal such a scheme could use.

\question{9. How were the results evaluated or validated?}

The claims are validated empirically. Each method the authors ran (BC, CQL, Fisher-BRC and TD3+BC) is trained with five seeds, evaluated every 5,000 steps over 10 episodes, and summarized by the D4RL normalized score (0 for a random policy, 100 for an expert) averaged over the last ten evaluations rather than the best checkpoint, with the standard deviation across seeds. CQL and Fisher-BRC were rerun from author code (CQL with its paper's settings, which beat the repository defaults), and the authors concede this CQL code falls short of its published scores. BRAC-p, AWAC and the three-seed Decision Transformer results were copied from earlier papers. Run times exclude evaluation and come from one GTX 1080 GPU. \ensuremath{\alpha} was swept over five values on hopper-medium and hopper-expert with a seed excluded from reported results; since both datasets are also test tasks, there is no held-out task split, a milder form of the gap the authors criticize in CQL and Fisher-BRC.

The protocol is fair in its reruns and unusually transparent about software versions and commits. Its weaknesses are the absence of proofs, confidence intervals and significance tests, only five seeds despite the high variance the paper documents, ablation bars showing mean percent differences without spread, and a near-tie with Fisher-BRC well within seed spread, so the evidence supports parity rather than superiority.

\question{10. What are the ethical implications? How are they addressed?}

The paper contains a brief Broader Impact statement in Appendix A and no separate ethics or limitations section. The statement says offline RL could open up uses in robotics and healthcare, where gathering data is slow or costly, warns of misuse such as making social media more habit-forming, and notes that learned policies inherit biases from the dataset and the policy that produced it. It then argues that because TD3+BC's performance gain is small, its main effect would be wider and easier access to offline RL rather than new applications. Section 4 also links policy instability to safety-critical use.

Further implications follow from the method itself. Behavior cloning ties the policy to logged behavior, so in medical or social settings it can reproduce inequities encoded in past decisions. The documented episode-to-episode variance means a policy with a good average can still fail badly on some runs, which matters for driving, robotics or patient care. Results come only from MuJoCo simulations and say nothing about real deployment. Tuning \ensuremath{\alpha} with rollouts and dropping normalization for AntMaze after seeing scores are forms of online selection that a true offline user could not repeat, although the appendix discloses them. Lower compute is a modest positive for energy use and access.

Overall, the paper names bias and misuse but proposes no mitigation; its candor about instability is its most useful safety contribution.

\question{11. What themes are suggested by this research? (minimum 3)}

\textbf{Minimalism and simplicity in algorithm design.} TD3+BC is defined by a one-line change to TD3's actor loss plus state standardization, with every other default left alone. The authors argue that such restraint makes performance gains attributable, shrinks the tuning burden and adds no compute, and they end by suggesting that the field has rewarded algorithmic complexity while overlooking simple alternatives.

\textbf{Behavior regularization and policy constraints.} The method counters extrapolation error by pulling the deterministic actor toward dataset actions with a squared-error penalty weighted against Q through \ensuremath{\lambda}. The paper treats the choice of distance as largely arbitrary and shows the constraint is indispensable on most datasets, yet its weak score on halfcheetah-random and its collapse on the Walker2d random-expert mixture (28.8 versus 100.3 for CQL) suggest that a uniform pull is a poor guide when much of the data is low quality.

\textbf{Implementation details and hyperparameter sensitivity.} Table 1 and Figure 1 show that CQL and Fisher-BRC rely on architecture changes, actor pre-training, reward bonuses and similar tweaks whose removal sharply lowers scores on many datasets. TD3+BC narrows the tunable surface to a single \ensuremath{\alpha} by normalizing \ensuremath{\lambda} with the mean absolute Q, and its sweep shows nearly identical results for \ensuremath{\alpha} of 2 and 3 but degradation on some tasks at 1 and 4.

\textbf{Online tuning and offline model selection.} Section 4 argues that every unexplained implementation change needs validation, and in a truly offline setting that validation requires the environment access offline RL is supposed to avoid. The checkpoint-to-checkpoint swings in Figure 3 make the stopping point another hidden choice. TD3+BC itself still fixed \ensuremath{\alpha} with a one-seed sweep on two Hopper datasets and dropped state normalization for AntMaze after finding it hurt performance there.

\textbf{Evaluation rigor, reproducibility and fair comparison.} The authors rerun the strongest baselines under one protocol, publish code with exact commits and software versions, and report averages over late evaluations instead of peaks. They also argue that mean scores conceal large episode-level and checkpoint-level variance in offline policies and call for benchmarks that measure stability, a weakness their own method shares.

\textbf{Computational efficiency.} Training time is treated as a primary criterion rather than an afterthought. Because TD3+BC needs no generative model, no sampled log-sum-exp and no enlarged networks, it trains in 39 minutes on the same GPU where Fisher-BRC needs over two hours and CQL over four, and it is also much faster than the transformer-based Decision Transformer.

\clearpage

\bookmarked{\subsection}{subsection}{Paper 3: Offline Reinforcement Learning with Implicit Q-Learning (IQL)}

\question{1. What is the citation?}

\begin{hangref}Kostrikov, I., Nair, A., \& Levine, S. (2022). Offline reinforcement learning with implicit Q-learning. In \emph{International Conference on Learning Representations (ICLR 2022)}. \url{https://openreview.net/forum?id=68n2s9ZJWF8}\end{hangref}
\begin{hangref}Version reviewed: Kostrikov, I., Nair, A., \& Levine, S. (2021). \emph{Offline reinforcement learning with implicit Q-learning} (arXiv:2110.06169v1) [Preprint]. arXiv. \url{https://doi.org/10.48550/arXiv.2110.06169}\end{hangref}

\question{2. What is the research area?}

Within machine learning, this paper addresses offline (batch) deep reinforcement learning, in which a fixed log of transitions is the only source of experience. IQL is a model-free actor-critic method for continuous control whose state-value and Q-networks are trained by temporal-difference updates that involve no actor, with a Gaussian policy extracted afterward. The problem family it addresses is extrapolation error, which arises when a critic is evaluated at actions the dataset never contains. IQL departs from policy constraints, whether built on a behavior density model (BCQ, BEAR) or a behavior-cloning term in the actor loss (TD3+BC), and from value regularization (CQL, Fisher-BRC). Its value learning is in-sample, like that of Onestep RL, but keeps multi-step dynamic programming, and its actor belongs to the advantage-weighted regression family (AWR, AWAC, CRR), whose implicit constraint acts only at extraction.

\question{3. What is the research question?}

Can an offline RL algorithm perform multi-step dynamic programming toward an optimal policy while training its value functions only on logged state-action pairs, so that the critic never has to evaluate an unseen action?

Three testable claims follow from it. The first is theoretical: that an upper expectile of Q-values taken over dataset actions approaches, as \ensuremath{\tau} goes to one, the best value achievable within the data's support. The second is that this multi-step property matters in practice, so the method should clearly beat one-step methods built on the behavior policy's value on tasks that require stitching sub-optimal trajectories, while staying competitive on locomotion. The third is that such an in-sample method can be cheap to train and can serve as a good initialization for online fine-tuning.

\question{4. What is the significance of the question?}

Offline RL matters because fields such as robotics and industrial planning already keep large records of past decisions, whereas letting an untrained controller experiment on real equipment or real operations can be expensive or unsafe. The technical obstacle is that improving on the data-collecting policy requires estimating returns for actions the dataset does not contain. When a critic is maximized over such actions, its approximation errors are actively selected, and bootstrapping compounds them into overestimation. In 2021 the field handled this with policy constraints or with penalties on out-of-distribution values, and both force a compromise: tighter control reduces error but also caps how far the policy can improve. The alternative of fitting the behavior policy's value and improving on it once sidesteps the problem but, as the authors show, breaks down when near-optimal trajectories are rare or absent and good behavior must be assembled from suboptimal segments. A method that removes out-of-distribution queries from value learning while keeping multi-step backups would ease this central tension, and cost mattered too, since the original CQL implementation needed more than four hours per million updates.

\question{5. What are the contributions reported in this paper?}

At the paper's core is the implicit Q-learning algorithm. Its value-learning stage replaces the mean-squared TD loss of SARSA-style evaluation with expectile regression, applied in a separate state-value network that estimates an upper expectile of target Q-values over dataset actions; the Q-network is then regressed with an ordinary squared loss onto the reward plus the discounted value of the next state. Separating the two networks is itself a design contribution, because it stops the expectile from rewarding transitions that reached good states only by chance under stochastic dynamics. The policy is obtained by advantage-weighted regression, so no stage of training evaluates an out-of-sample action and the critic never depends on the actor. The authors call it the first method to avoid unseen-action queries while keeping multi-step backups.

The paper also contributes a short theoretical analysis: a lemma that expectiles approach the supremum of a bounded random variable as \ensuremath{\tau} goes to one, a monotonicity lemma for the recursively defined value functions, and a theorem that the limiting value equals the optimal value restricted to the behavior policy's support. An appendix relates IQL to BCQ: both target the same batch-constrained objective, but IQL enforces it through expectiles rather than through a generative model that can still propose unseen actions.

Empirically, the authors supply a toy maze illustration contrasting one-step and multi-step evaluation, a D4RL comparison that runs every method on the -v2 locomotion datasets to correct an inconsistency in earlier comparisons, runtime measurements of JAX reimplementations, an expectile ablation, and an offline-to-online fine-tuning study. They also release a JAX/Flax implementation.

\question{6. What is the research method?}

IQL is developed through algorithm design and a limit-based theoretical argument, then tested by empirical benchmarking with an ablation and a fine-tuning study. The technical core is the asymmetric squared loss $L_2^{\tau}(u)=\lvert\tau-\mathbf{1}(u<0)\rvert\,u^{2}$, which for $\tau>0.5$ penalizes underestimates more heavily than overestimates. IQL fits $V_{\psi}(s)$ by minimizing $\mathbb{E}[L_2^{\tau}(Q_{\hat{\theta}}(s,a)-V_{\psi}(s))]$ over dataset pairs, using the minimum of two target Q-networks, and fits $Q_{\theta}$ by minimizing $\mathbb{E}[(r+\gamma V_{\psi}(s')-Q_{\theta}(s,a))^{2}]$, with target networks updated by Polyak averaging at rate 0.005. The policy maximizes $\mathbb{E}[\exp(\beta(Q(s,a)-V(s)))\log\pi(a\mid s)]$, where \ensuremath{\beta} is an inverse temperature and the exponentiated advantage is capped at 100; small \ensuremath{\beta} approaches behavior cloning and large \ensuremath{\beta} approaches greedy improvement.

All networks are two-layer ReLU MLPs with 256 units trained by Adam at $3\cdot10^{-4}$, with a cosine schedule on the actor's learning rate and a Gaussian policy with a state-independent standard deviation. Settings differ by domain: \ensuremath{\tau} = 0.7 and \ensuremath{\beta} = 3.0 for locomotion, \ensuremath{\tau} = 0.9 and \ensuremath{\beta} = 10.0 for AntMaze, and \ensuremath{\tau} = 0.7 and \ensuremath{\beta} = 0.5 with 0.1 policy dropout for Kitchen and Adroit. Locomotion rewards are divided by the spread between the best and worst dataset returns, and AntMaze rewards are reduced by one. Training curves and runtimes cover 1M gradient steps; the batch size is not reported.

The experiments comprise a stochastic u-shaped maze (one optimal and 99 random trajectories, \ensuremath{\gamma} = 0.9, \ensuremath{\tau} = 0.95) comparing IQL with Onestep RL; D4RL locomotion -v2 and AntMaze -v0 datasets against BC, 10\%BC, Decision Transformer, AWAC, Onestep RL, TD3+BC and CQL; a smaller Kitchen and Adroit comparison against BC, BRAC-p, BEAR, CQL and, on cloned Adroit tasks, Onestep RL; a \ensuremath{\tau} ablation over 0.5, 0.7 and 0.9 on the four medium and large AntMaze datasets; and fine-tuning against the original AWAC and CQL code, in which 1M offline gradient steps precede 1M online steps at one update per environment step (\ensuremath{\tau} = 0.8, \ensuremath{\beta} = 3.0, 25,000 offline steps and exploration noise of 0.03 on the hand tasks).

\question{7. What are the results?}

On the toy maze, one-step values faded toward zero a short distance from the goal, yielding a poor policy, whereas IQL's values closely matched the true optimum. Across the nine D4RL locomotion datasets IQL totaled 692.4 normalized points, slightly below CQL's 698.5 and above Onestep RL (684.6) and TD3+BC (677.4). Its advantage lay in AntMaze: 378.0 against 303.6 for CQL and 100.2 to 163.8 for the rest. On the two large mazes it reached 39.6 and 47.5 versus CQL's 15.8 and 14.9, while BC, Onestep RL and Decision Transformer scored near zero beyond umaze. IQL lost on some datasets: antmaze-umaze-diverse (62.2 versus 84.0 for CQL), hopper-medium-expert (91.5 versus 110.9 for 10\%BC) and walker2d-medium-replay (73.9 versus 81.8 for TD3+BC). Its Kitchen and Adroit totals (159.8 and 118.1) beat BC (154.5, 104.5) and CQL (144.6, 93.6), but the Adroit lead over CQL came entirely from pen-human (71.5 versus 37.5), and the hammer, door and relocate tasks stayed below 10 points for every method.

An IQL run took about 20 minutes, a quarter of the authors' JAX CQL (80 minutes) but no faster than TD3+BC, AWAC or Onestep RL. In the \ensuremath{\tau} ablation on the medium and large mazes, 0.5 stayed near zero and 0.7 peaked at roughly 15, while 0.9 reached about 40 to 75. In fine-tuning, IQL raised its AntMaze total from 370.1 to 473.7 (CQL 151.5 to 231.1, AWAC almost unchanged) and its hand total from 38.1 to 124.0. It still finished below CQL on both umaze mazes (96.0 and 84.0 versus 99.4) and below AWAC on pen-binary (60.7 versus 70.3), contrary to the Table 2 caption's claim that pen-binary is the only exception.

\question{8. What are the impacts and limitations of the contributions?}

IQL's main impact is to show that in-sample learning need not give up multi-step dynamic programming: a critic trained only on logged actions can still stitch trajectories. Because the critic never depends on the actor, the extraction step can be replaced or run alongside value learning, which is how IQL keeps training during online fine-tuning. It has since been reimplemented in later codebases, including the Unifloral framework this project builds on.

The paper states few limitations: it notes that \ensuremath{\tau} close to one makes optimization harder, so \ensuremath{\tau} must be tuned, and concedes that Kitchen and Adroit barely separate methods. The analysis and experiments expose others. The theory assumes exact optimization as \ensuremath{\tau} \ensuremath{\to} 1, yet practice uses \ensuremath{\tau} = 0.7 or 0.9, and because a continuous state usually appears with a single logged action, the in-support maximum is reached only through the network's generalization across similar states. The separate value network guards against chance outcomes of the dynamics, but not against noise in the Q estimates themselves, which an upper expectile also tends to favor. Because the actor only reweights logged actions, it can favor only actions resembling those recorded in similar states, and the paper leaves unexplained why IQL trails filtered cloning on hopper-medium-expert.

IQL's \ensuremath{\beta} and its cap of 100 are shared by every transition, so an advantage inflated by critic error receives the same amplification as a reliable one. A calibrated per-sample width that tempers only the portion of the capped weight above one would damp that amplification while leaving the cloning baseline intact.

\question{9. How were the results evaluated or validated?}

Validation rests mainly on simulated D4RL benchmarks reported as averaged normalized scores, without restating the normalization formula. Locomotion results average 10 evaluation trajectories over 10 random seeds; AntMaze results average 100 evaluation trajectories, and no seed count is given for AntMaze, Kitchen, Adroit or fine-tuning. The tables give means only, with no standard deviations, confidence intervals or significance tests, and Table 1 does not say whether final or best checkpoints are reported (Table 2 states final performance).

Baseline numbers come from mixed sources: TD3+BC and Onestep RL scores supplied by their authors (Onestep's scores used a different batch and network size from its timed run), CQL on -v2 from an author-suggested implementation, and CQL AntMaze scores taken from the original paper because the authors could not reproduce them, which is why CQL starts much lower in Table 2. Baseline runtimes come from the authors' JAX reimplementations. Hyperparameters vary by domain with no stated selection procedure, so online tuning cannot be ruled out. The toy-maze evidence that IQL recovers near-optimal values is visual, with no error metric; the proofs are short, one only sketched; and the sole ablation varies \ensuremath{\tau}, leaving \ensuremath{\beta}, the separate value network and the advantage cap untested. The breadth of datasets and the corrected -v2 comparison are strengths.

\question{10. What are the ethical implications? How are they addressed?}

IQL's arXiv v1 has no ethics statement and no broader-impact or limitations section. What comes nearest is the introduction's motivating list of application areas in which learning by trial and error would be risky or expensive, and a footnote linking the released code, which supports reproducibility.

The method and its evaluation still carry ethical implications. Because IQL learns only from logged actions and reweights them by estimated advantage, any systematic bias in how the data were gathered, such as unsafe or inequitable decisions in a clinical or operational log, is inherited and can be amplified whenever those decisions score well under the recorded reward. The upper expectile is deliberately optimistic, and with limited data the exponential weight can magnify critic errors, which is a concern in safety-critical use. IQL is tested only on simulated tasks, and its fine-tuning study involves live exploration with Gaussian noise that would carry physical risk on real hardware. Domain-specific hyperparameters without a stated selection method also raise a research-integrity question, since tuning on online returns would contradict the offline premise. Its low compute cost does reduce energy use.

None of these issues is raised explicitly. Its in-sample design offers some implicit protection against acting far outside the data, but this is presented as a performance property rather than a safety argument.

\question{11. What themes are suggested by this research? (minimum 3)}

\textbf{Distribution shift and out-of-distribution value overestimation.} IQL's design is a direct response to extrapolation error. Rather than penalizing or constraining unseen actions, it removes them from value learning entirely: both the value and Q losses use only dataset actions, and the actor is fit by weighted likelihood of logged actions. The authors argue that even the in-distribution constraints of earlier methods may not fully prevent extrapolation, so avoiding the query altogether is presented as the cleaner remedy.

\textbf{Conservatism versus policy improvement trade-off.} The expectile \ensuremath{\tau} turns this trade-off into a single continuous dial: \ensuremath{\tau} = 0.5 recovers SARSA-style evaluation of the behavior policy, and \ensuremath{\tau} approaching one recovers Q-learning restricted to the data's support. On the medium and large AntMaze datasets, too little optimism (\ensuremath{\tau} = 0.5 or 0.7) cripples performance, although 0.7 is the setting used for locomotion, and the authors note that values near one are harder to optimize, so the dial is set per domain.

\textbf{Behavior regularization and policy constraints.} Although value learning carries no explicit constraint, the policy remains tied to the data through advantage-weighted regression, whose inverse temperature \ensuremath{\beta} moves the actor between plain behavior cloning and greedy improvement. The domain-specific \ensuremath{\beta} of 0.5, 3 or 10 and the cap of 100 on exponentiated advantages therefore act as the real behavior constraint, so IQL moves conservatism from the critic into the actor rather than removing it.

\textbf{Dataset quality and composition.} The paper argues that dataset composition determines which family of algorithms succeeds. Locomotion datasets with a substantial share of near-optimal trajectories let simple one-step methods do well, while AntMaze datasets in which near-optimal paths are scarce or missing demand stitching segments together through repeated Bellman backups. The toy maze, built from one optimal and 99 random trajectories, isolates this effect, and the AntMaze results confirm it at benchmark scale.

\textbf{Minimalism and simplicity in algorithm design.} IQL is presented as a light adjustment to a familiar SARSA-style update: one extra critic trained with an asymmetric squared loss and an actor fit by weighted behavior cloning. The design is cheap as well as simple, since a million updates fit in under 20 minutes on a single GTX 1080 GPU, which makes IQL attractive both as a baseline and as a host for extensions.

\textbf{Offline-to-online fine-tuning.} Because the actor does not influence the critic, IQL can keep learning while fresh experience enters its replay buffer, and the authors attribute its suitability for fine-tuning to a strong offline starting point and weighted-cloning extraction. It gained substantially on AntMaze and the hand tasks, whereas CQL sometimes collapsed (antmaze-medium-play fell from 23.0 to 0.0) and AWAC barely moved on AntMaze, tying the choice of offline constraint to later adaptability.

\textbf{Evaluation rigor, reproducibility and fair comparison.} The paper corrects earlier comparisons that scored CQL on older -v0 locomotion environments while reporting their own methods on -v2, and it evaluates every method on -v2. Its own protocol still has gaps: Appendix C admits that CQL's published AntMaze scores could not be reproduced, baseline numbers come from mixed sources, and the tables report means without any measure of spread or a stated rule for choosing hyperparameters.

\clearpage

\bookmarked{\subsection}{subsection}{Paper 4: Revisiting the Minimalist Approach to Offline Reinforcement Learning (ReBRAC)}

\question{1. What is the citation?}

\begin{hangref}Tarasov, D., Kurenkov, V., Nikulin, A., \& Kolesnikov, S. (2023). Revisiting the minimalist approach to offline reinforcement learning. \emph{Advances in Neural Information Processing Systems}, \emph{36}, 11592--11620. \url{https://proceedings.neurips.cc/paper_files/paper/2023/hash/26cce1e512793f2072fd27c391e04652-Abstract-Conference.html}\end{hangref}
\begin{hangref}Version reviewed: Tarasov, D., Kurenkov, V., Nikulin, A., \& Kolesnikov, S. (2023). \emph{Revisiting the minimalist approach to offline reinforcement learning} (arXiv:2305.09836v2) [Preprint]. arXiv. \url{https://doi.org/10.48550/arXiv.2305.09836}\end{hangref}

\question{2. What is the research area?}

The ReBRAC paper works in offline (batch) reinforcement learning, a machine learning setting in which a control policy is trained on logged transitions without any interaction with the environment. The method is model-free and uses the deterministic actor-critic structure of TD3. Its lineage runs through BRAC and TD3+BC, the behavior-regularization (policy-constraint) family that pulls the learned policy toward dataset actions. The authors compare it with value regularization (CQL), in-sample learning (IQL), more elaborate ensemble-free methods that penalize unfamiliar actions through a learned novelty signal (SPOT and SAC-RND), and ensemble-based uncertainty methods (SAC-N, EDAC, RORL, MSG). Its critic penalty adds a value-regularization component, but the contribution lies in testing implementation choices on an existing method rather than in a new algorithmic principle.

\question{3. What is the research question?}

The central question is: how far can a minimalist, ensemble-free offline RL algorithm such as TD3+BC be improved simply by adopting the small design and implementation choices that newer methods introduced alongside their main ideas, and which of those choices account for the improvement?

The authors test this in stages. They ask whether a TD3+BC variant with deeper networks, critic layer normalization, larger batches, separately weighted actor and critic penalties, and a larger discount factor for sparse-reward tasks can reach state-of-the-art ensemble-free performance on proprioceptive (D4RL) and pixel-based (V-D4RL) benchmarks. They ask whether the same agent fine-tunes well once online interaction is allowed. Finally, they ask which modification carries the gain, whether additional depth helps further, and how sensitive the method is to its two penalty coefficients when only a few policies can be evaluated online.

\question{4. What is the significance of the question?}

Anyone deploying offline RL needs methods that are trustworthy and cheap enough to run routinely, and both qualities depend on knowing what actually drives a method's scores. Every such method must cope with extrapolation error: the critic must estimate values for actions the dataset never contains, and an actor that maximizes those estimates tends to exploit its mistakes. Many algorithms address this problem, but each new method tends to arrive with extra machinery and quiet engineering changes, which makes it hard to tell whether a gain comes from the core idea or from the scaffolding around it.

In 2023, TD3+BC was the standard simple baseline against which most new methods claimed improvements, yet no one had given TD3+BC the same engineering upgrades those methods enjoyed. Ensemble methods had pushed Gym-MuJoCo scores higher, but at substantial computational cost (SAC-N needed as many as 500 critics on hopper, and MSG needed 64 members on AntMaze). Establishing how strong a cheap baseline can become therefore affects how the field credits algorithmic progress, and it gives practitioners a fast default method.

\question{5. What are the contributions reported in this paper?}

ReBRAC's contributions fall into five groups. The first is a retrospective audit of recent methods (TD3+BC, CQL, EDAC, MSG, CNF, LB-SAC and SAC-RND) that catalogs which non-algorithmic choices each one adopted: deeper networks, larger batches, layer normalization, separate actor and critic penalties, and a modified discount factor.

The second is ReBRAC (Revisited BRAC), an algorithm that layers these choices onto TD3+BC. It generalizes BRAC by penalizing both the actor and the critic with a mean-squared-error divergence from dataset actions, each with its own coefficient, while keeping TD3+BC's Q-value normalization and dropping its state normalization.

The third is a broad evaluation on 51 datasets from D4RL (Gym-MuJoCo, AntMaze, Adroit) and V-D4RL, plus an offline-to-online fine-tuning study, with TD3+BC and IQL re-tuned under the same seed protocol so that the comparison is more even-handed.

The fourth is a set of diagnostic studies: a leave-one-out ablation of each modification, tests that add each modification to TD3+BC on its own, a sweep over actor and critic depth, a comparison of LayerNorm with feature normalization, and a sensitivity analysis of the two penalty coefficients using Expected Online Performance.

The fifth is a reproducibility package consisting of an open-source implementation (JAX for D4RL), per-dataset hyperparameter tables for ReBRAC and the re-tuned baselines, and compute-cost tables for the main offline benchmarks, the ablations and the depth study.

\question{6. What is the research method?}

The research method is algorithm design by composition, followed by large-scale empirical benchmarking and ablation; the paper contains no theoretical analysis. ReBRAC starts from TD3+BC, a deterministic actor-critic, and keeps most of its hyperparameters. The actor objective is \begin{equation*}\pi=\arg\max_{\pi}\,\mathbb{E}_{(s,a)\sim\mathcal{D}}\big[Q(s,\pi(s))-\beta_1(\pi(s)-a)^{2}\big],\end{equation*} where, as in TD3+BC, the Q term is divided by its batch-average magnitude so that $\beta_{1}$ is less scale-dependent. The critic regresses toward \begin{equation*}r+\gamma\big(\bar{Q}(s',a')-\beta_2(a'-\hat{a}')^{2}\big),\end{equation*} where $a'$ is the policy's action at the next state, $\hat{a}'$ is the next action stored in the dataset, and $\bar{Q}$ is a target network. Thus $\beta_{1}$ sets actor-side conservatism, while $\beta_{2}$ lowers the bootstrapped value when the policy departs from logged behavior; the original BRAC tied the two weights together.

The other choices are three hidden layers of 256 units in both networks, LayerNorm between critic layers, batch size 1024 with learning rate 1e-3 on Gym-MuJoCo (batch 256 elsewhere), and \ensuremath{\gamma} = 0.999 rather than 0.99 on AntMaze, because $0.99^{1000}\approx4\cdot10^{-5}$ nearly erases a sparse terminal reward. AntMaze rewards are also multiplied by 100. Only $\beta_{1}$ and $\beta_{2}$ are tuned per dataset (plus learning rates on AntMaze), over grids such as $\beta_1\in\{0.001,0.01,0.05,0.1\}$ and $\beta_2\in\{0,0.001,0.01,0.1,0.5\}$ for Gym-MuJoCo and Adroit.

Agents train for one million steps on D4RL Gym-MuJoCo v2 (18 datasets), AntMaze v2 (6) and Adroit v1 (12); 15 V-D4RL image datasets at 84\ensuremath{\times}84 pixels are added. D4RL baselines are TD3+BC, IQL, CQL and SAC-RND (plus behavior cloning on Adroit), with ensemble methods in an appendix; V-D4RL baselines (Offline DV2, DrQ+BC, LOMPO, CQL, behavior cloning) come from Lu et al. (2022). For fine-tuning against TD3+BC, IQL, CQL, SPOT and Cal-QL, ReBRAC sets $\beta_{2}$ to zero and linearly decays $\beta_{1}$ to half its value, without further tuning. Diagnostic runs disable one change at a time, add each change alone to TD3+BC in the same code, vary depth from two to six layers on AntMaze, and compute Expected Online Performance over the tuning grid.

\question{7. What are the results?}

ReBRAC averaged 81.2 on D4RL Gym-MuJoCo, second to SAC-RND (82.6) and ahead of CQL (73.6) and IQL (72.9). It led on AntMaze (76.8 versus 60.9 for SAC-RND) and Adroit (58.6 versus 53.4 for IQL). Ensemble methods stayed ahead on Gym-MuJoCo (EDAC 85.2) and AntMaze (MSG 80.5) but not on Adroit (RORL 41.0). On V-D4RL its average of 35.5 beat DrQ+BC (30.4) and behavior cloning (30.3), yet it was best on only four of fifteen datasets: Offline DV2 won the random and mixed walker and cheetah data, and behavior cloning the expert and most humanoid data.

After online fine-tuning the overall average rose from 55.0 to 91.1, ahead of IQL (65.10) and Cal-QL (58.08), mostly through the cloned Adroit tasks (24.0 to 100.1). On AntMaze, however, Cal-QL (97.33) and SPOT (92.38) finished above ReBRAC (85.2), and antmaze-large-play fell from 47.0 to 39.5. ReBRAC's average cumulative regret, 0.18, was the lowest.

In four-seed ablations, removing the actor penalty dropped the overall score from 72.2 to 11.4, and removing LayerNorm cut it to 38.0, with AntMaze falling to zero. On AntMaze, reverting to \ensuremath{\gamma} = 0.99 or dropping the third hidden layer cut the score from 75.3 to about 20, while the large batch helped only on Gym-MuJoCo. Removing the critic penalty or the decoupling changed the average by less than one percent, and TD3+BC, re-tuned or given any single modification, stayed at or below 62 overall. Four critic and three actor layers worked best on AntMaze. Expected Online Performance showed ReBRAC needing about ten evaluated policies to reach the ensemble-free state of the art; with one policy it beat IQL on AntMaze (67.9 versus 29.8) but trailed it on Gym-MuJoCo and Adroit.

\question{8. What are the impacts and limitations of the contributions?}

ReBRAC's main effect is to reset the baseline. A TD3+BC variant needing roughly 0.3 to 0.6 hours per D4RL run, against about two for IQL (the PyTorch CORL implementation the authors ran, not IQL's faster JAX code) and five for MSG, matches or beats more elaborate ensemble-free methods, although its image-based V-D4RL runs took about ten hours. New algorithms must now show gains over a properly engineered simple method, and the ablations give practical guidance: critic LayerNorm and the presence of an actor penalty are essential, and the discount factor matters for sparse rewards.

The authors concede that they tested only a handful of changes, all on TD3+BC; whether these transfer to IQL, CQL or MSG, and whether untried options such as dropout, learning-rate schedules, wider networks or stochastic actors would add more, is left open. On close reading, performance also depends on per-dataset tuning of $\beta_{1}$ and $\beta_{2}$ chosen by evaluating policies in the environment, and the selected values span two orders of magnitude ($\beta_{1}$ from 0.001 to 0.1 on Gym-MuJoCo). Like other ensemble-free methods, ReBRAC scores below 10 on the human and cloned door, hammer and relocate datasets. No theory explains how the choices interact, AntMaze results vary widely across seeds, and no experiment leaves simulation.

ReBRAC's ablations point BCA to the actor's behavior-cloning multiplier as the term worth calibrating: ReBRAC collapses without it, while the critic penalty moves the average score by under one percent. The best actor coefficient also shifts across datasets, so a global setting must be re-tuned per dataset; the paper never tests whether strength should also vary between samples within a dataset.

\question{9. How were the results evaluated or validated?}

Scores are D4RL-normalized, and V-D4RL returns are divided by ten to the same 0 to 100 scale. Hyperparameters were chosen on four training seeds, and the selected configuration was rerun on ten new seeds (two and five on V-D4RL), reporting the mean and standard deviation of the final checkpoint rather than the best one. Final evaluation used ten episodes on Gym-MuJoCo and Adroit and 100 on AntMaze. TD3+BC and IQL were re-tuned under the same protocol and SAC-RND was rerun, while CQL, the V-D4RL baselines and most fine-tuning baselines were copied from earlier papers. As a control, TD3+BC was run in the authors' code with its original and re-tuned coefficient, showing that the gains stem neither from implementation differences nor from tuning alone. Aggregate comparisons use the performance profiles and probability of improvement of Agarwal et al. (2021); ablations use four seeds; and Expected Online Performance measures sensitivity to the online evaluation budget.

Separating tuning seeds from reporting seeds and reporting last checkpoints makes the protocol stronger than usual. Selection, however, still uses true environment returns, which a real offline user lacks. Copied baselines come from other protocols and seed counts (four seeds against ReBRAC's ten in fine-tuning), no formal significance tests are given, and a few tuned values in the hyperparameter tables fall outside the stated search grids.

\question{10. What are the ethical implications? How are they addressed?}

Societal consequences go undiscussed: the paper has no ethics statement or broader-impact section. Related material is limited to a short limitations paragraph in Section 6, which concerns scientific scope only, together with a compute-cost appendix and a link to the public code.

ReBRAC's intended settings still raise issues the paper leaves aside. Offline RL is aimed at domains such as healthcare, driving and industrial robotics, where logged data are plentiful and exploration is dangerous. A policy constrained toward recorded behavior inherits whatever biases or unsafe habits that behavior contains, yet ReBRAC's evidence comes entirely from simulated locomotion, maze and manipulation tasks. Its hyperparameters were chosen by running candidate policies in the simulator, a step that would itself carry risk in a real system. The fine-tuning recipe also relaxes both penalties while the agent is acting, and on one AntMaze dataset the score fell during that phase; outside simulation, such a regression would mean worse behavior during live operation. On research integrity the paper is comparatively transparent: it describes its tuning protocol, releases code, and reports compute (roughly 18,500 run-hours on V100 and A100 GPUs when the appendix tables are summed). The energy cost is disclosed but not discussed.

In short, the authors document tuning, code and compute carefully but say nothing about deployment safety, data bias or validation outside simulation.

\question{11. What themes are suggested by this research? (minimum 3)}

\textbf{Implementation details and hyperparameter sensitivity.} The paper's premise is that choices often treated as incidental, such as network depth, LayerNorm, batch size and the discount factor, explain much of the progress credited to new algorithms. Its ablations bear this out: removing critic LayerNorm roughly halves average performance, and the default discount cuts AntMaze performance by about 70 percent. The wide spread of the best penalty values across datasets shows that results also hinge on per-dataset tuning.

\textbf{Minimalism and simplicity in algorithm design.} ReBRAC adds no new learned components to TD3+BC: no ensembles, generative models or density estimators, only a second MSE penalty and standard deep-learning practice. It still matches or beats more complex ensemble-free methods and, on D4RL, runs several times faster per run than the IQL, SAC-RND and MSG implementations the authors timed, which argues that simple baselines should be strengthened before complexity is added.

\textbf{Behavior regularization and policy constraints.} ReBRAC is a behavior-regularized actor-critic in which an MSE penalty toward dataset actions is the main safeguard against extrapolation. The ablations separate the two penalties: without the actor penalty the method collapses on nearly every domain, while the critic penalty and decoupling add little on average. The best coefficients still differ by dataset, so constraint strength remains a single scalar re-tuned for each dataset.

\textbf{Evaluation rigor, reproducibility and fair comparison.} The paper's case rests on giving an old baseline the engineering and tuning effort that newer methods had received. Its compute tables show ReBRAC and IQL each receiving 1,440 tuning runs on Gym-MuJoCo, and the authors concede that baselines whose scores they copied might score lower under their protocol. They also tune only the penalty weights, although per-task tuning of the design choices would raise scores further, which keeps the comparison closer to a practitioner's budget.

\textbf{Online tuning and offline model selection.} Every reported score relies on selecting $\beta_{1}$ and $\beta_{2}$ by evaluating policies in the environment. The paper confronts this in part through Expected Online Performance, which shows that ReBRAC needs around ten online evaluations before its expected score reaches the ensemble-free state of the art, and that it is weaker than IQL on two domains when only one policy can be tested. Offline model selection without environment access is left open.

\textbf{Offline-to-online fine-tuning.} ReBRAC was not built for fine-tuning, yet the authors simply continue training it online after switching off the critic penalty and decaying the actor penalty linearly to half, with no new tuning. This recipe gives the lowest average regret and lifts the cloned door and hammer tasks from near zero to above 100, but it trails Cal-QL and SPOT on AntMaze and loses ground on one large maze, so loosening a fixed constraint helps unevenly.

\clearpage

\bookmarked{\subsection}{subsection}{Paper 5: Importance Weighted Actor-Critic for Optimal Conservative Offline Reinforcement Learning (A-Crab)}

\question{1. What is the citation?}

\begin{hangref}Zhu, H., Rashidinejad, P., \& Jiao, J. (2023). Importance weighted actor-critic for optimal conservative offline reinforcement learning. \emph{Advances in Neural Information Processing Systems}, \emph{36}, 49579--49602. \url{https://proceedings.neurips.cc/paper_files/paper/2023/hash/9bb93a3c1a424654aaea6f5b594e94d5-Abstract-Conference.html}\end{hangref}
\begin{hangref}Version reviewed: Zhu, H., Rashidinejad, P., \& Jiao, J. (2023). \emph{Importance weighted actor-critic for optimal conservative offline reinforcement learning} (arXiv:2301.12714v2) [Preprint]. arXiv. \url{https://doi.org/10.48550/arXiv.2301.12714}\end{hangref}

\question{2. What is the research area?}

This paper sits in machine learning, under reinforcement learning, and specifically in offline (batch) reinforcement learning, where the learner sees only a fixed log of transitions. The algorithm is model-free and actor-critic, but the work is mainly theoretical: it studies provably efficient offline RL with general function approximation when the data cover only part of the state-action space. Its problem family is pessimistic, or conservative, policy optimization. The paper joins two earlier lines of work: adversarially trained actor-critic methods (ATAC; Cheng et al., 2022), and marginalized importance sampling, which works with density ratios between a policy's occupancy and the data distribution (Zhan et al., 2022; Rashidinejad et al., 2022). A short experimental section connects the theory to deep offline RL on D4RL.

\question{3. What is the research question?}

A-Crab targets one question: can a single offline RL algorithm be computationally tractable, reach the statistically optimal $1/\sqrt{N}$ suboptimality rate with general function approximators, and still require only a weak form of data coverage?

The paper makes the question concrete through narrower claims. It asks whether replacing the squared Bellman-error penalty in ATAC's critic with an importance-weighted average Bellman error is enough to obtain the optimal rate against any policy the data cover, not only against the optimal one. It asks whether an $\ell_2$ version of single-policy concentrability is weaker than the $\ell_\infty$ and Bellman-consistent coefficients used earlier. It also asks whether the relative-pessimism structure still guarantees improvement over the behavior policy across a wide range of the penalty weight \ensuremath{\beta}. Finally, it asks whether ATAC's slower rate reflects a real defect of the squared penalty, and whether a practical version of the change helps on D4RL.

\question{4. What is the significance of the question?}

Coverage assumptions decide whether a guarantee is usable. Offline RL is most attractive where exploration is expensive or unsafe, and the authors point to precision medicine, robotics and climate as examples. In such settings the logged data almost never cover every candidate policy, so guarantees that assume all-policy coverage are of little use. Even single-policy guarantees stated with an $\ell_\infty$ coefficient can become vacuous when one rarely logged action makes a density ratio very large.

On the theory side, the field faced an uncomfortable trade-off. Methods with mild coverage requirements and the optimal rate (Xie et al., 2021; Uehara and Sun, 2021) could not be computed efficiently. ATAC was tractable but converged at only $1/N^{1/3}$. Efficient methods built on density ratios or linear programming (Rashidinejad et al., 2022; Ozdaglar et al., 2022) reached $1/\sqrt{N}$ only when the optimal policy was covered, and only under $\ell_\infty$ coverage. Bellman-consistent pessimism and ATAC also relied on approximate Bellman completeness, a closure condition, much stronger than realizability, that requires F to contain the Bellman backup of each of its members. A small change to an algorithm that already runs with deep networks, if it removed all of these compromises, would narrow the gap between offline RL theory and practice.

\question{5. What are the contributions reported in this paper?}

The main contribution is A-Crab (Actor-Critic Regularized by Average Bellman error). It keeps ATAC's relative-pessimism actor-critic game but changes the critic's penalty. Instead of a squared Bellman residual, the critic pays the largest absolute importance-weighted average Bellman residual over a class of marginalized importance weights. For any fixed weight function this penalty is linear in the residuals, so its sample average is unbiased, and A-Crab no longer needs the extra subtraction term that ATAC used to offset the bias of the squared residual.

On the theory side, the authors provide a finite-sample guarantee that A-Crab attains the optimal $1/\sqrt{N}$ rate against any policy whose density ratio lies in the weight class. The guarantee requires realizability of every policy's Q-function in F and a uniform $\ell_2$ bound on the weight class, needs no completeness assumption, and degrades gracefully when realizability holds only approximately. They introduce an $\ell_2$ single-policy concentrability coefficient and compare it formally with the $\ell_\infty$ and Bellman-consistent coefficients. They prove a robust policy improvement theorem for every $\beta\ge 0$ and a negative result: a bandit instance on which ATAC cannot reach the optimal rate even under realizability, completeness and its recommended \ensuremath{\beta}. They also argue that the critic's inner problem over F \ensuremath{\times} W is smaller than ATAC's over F \ensuremath{\times} F.

Turning to practice, they adopt the box-shaped weight class of Hong et al. (2023) to obtain a closed-form penalty, mix it with the squared Bellman error, compare the result with ATAC on twelve D4RL datasets and release their code.

\question{6. What is the research method?}

Theory dominates the method: algorithm design is paired with finite-sample analysis and a lower-bound construction, and a small empirical comparison follows. The setting is a discounted MDP with a finite action set and an arbitrary state space. The learner has finite classes of Q-functions F, importance weights W and policies \ensuremath{\Pi}, plus N transitions whose state-action pairs are drawn i.i.d. from the behavior policy's discounted occupancy \ensuremath{\mu}. Q-functions take values in $[0,V_{\max}]$ and weights in $[0,B_w]$, and W always contains the constant function 1, the behavior policy's own density ratio, so the behavior policy is always a valid comparator.

A-Crab is set up as a Stackelberg game. The actor looks for the policy that maximizes the relative value $L_{\mathcal{D}}(\pi,f)=\mathbb{E}_{\mathcal{D}}[f(s,\pi)-f(s,a)]$, which measures how much \ensuremath{\pi} improves on the logged actions under critic f. For the current policy, the critic picks the f that minimizes $L_{\mathcal{D}}(\pi,f)+\beta\cdot\mathcal{E}_{\mathcal{D}}(\pi,f)$, where $\mathcal{E}_{\mathcal{D}}(\pi,f)=\max_{w\in W}\big\lvert\mathbb{E}_{\mathcal{D}}[w(s,a)(f(s,a)-r-\gamma f(s',\pi))]\big\rvert$. The actor is updated by a no-regret oracle such as natural policy gradient ($\pi_{k+1}\propto\pi_k\cdot\exp(\eta f_k)$), and the output is a uniform mixture of the K iterates. The main theorem fixes $\beta=2$.

For the experiments, the inner maximization is made tractable by letting W contain all functions with values in $[0,C_\infty]$. This gives a closed form: $C_\infty$ times the larger of the mean positive part and the mean negative part of the empirical Bellman residuals. Adding the squared Bellman error worked better in practice, so the implemented penalty is $\tfrac{1}{2}\big(2\mathcal{E}_{\mathcal{D}}^{\mathrm{app}}+\beta\cdot\mathbb{E}_{\mathcal{D}}[(f-\mathcal{T}^{\pi}f)^{2}]\big)$. Everything else follows a lightweight ATAC implementation with ATAC's hyperparameters, including \ensuremath{\beta}, which ranges from 0.062 to 64 across datasets. $C_\infty$ is chosen per dataset from $\{1,2,5,10,20,50,100,200\}$. The testbed is twelve D4RL MuJoCo v2 datasets: hopper, walker2d and halfcheetah, each with random, medium, medium-replay and medium-expert data. Both methods are pre-trained for 100 epochs in ATAC's manner, the curves cover the following one million training steps, and ATAC is the only baseline.

\question{7. What are the results?}

Theorem 1 bounds the gap between A-Crab's output and a comparator \ensuremath{\pi} with $w^{\pi}\in W$ by $O(\epsilon_{\mathrm{stat}}+C^{*}_{\ell_2}\cdot\sqrt{\epsilon_F})+\epsilon_{\mathrm{opt}}$. Here $\epsilon_{\mathrm{stat}}$ scales as $V_{\max}\cdot C^{*}_{\ell_2}\cdot\sqrt{\log(\lvert F\rvert\lvert\Pi\rvert\lvert W\rvert/\delta)/N}$ plus a lower-order term $V_{\max}\cdot B_w\cdot\log(\lvert F\rvert\lvert\Pi\rvert\lvert W\rvert/\delta)/N$, where $B_w$ is the weights' sup-norm bound. With exact realizability, enough iterations and the optimal policy as comparator, this gives the $1/\sqrt{N}$ rate, improving on ATAC ($1/N^{1/3}$) and Zhan et al. (2022) ($1/N^{1/6}$). Theorem 2 shows that, for any $\beta\ge 0$, the output trails the behavior policy by at most order $(\beta+1)(\epsilon_{\mathrm{stat}}+C^{*}_{\ell_2}\cdot\sqrt{\epsilon_F})+\epsilon_{\mathrm{opt}}$, which vanishes when $\beta=o(\sqrt{N})$. In Proposition 4.1, a two-armed bandit forces ATAC with $\beta=\Theta(N^{2/3})$ into $\Omega(1/N^{1/3})$ suboptimality with constant probability; the proof covers any \ensuremath{\beta} growing faster than $\sqrt{N}$. The $\ell_2$ coefficient never exceeds the $\ell_\infty$ one, and in a bandit example it is at most $\sqrt{2}$ against $1/\epsilon$.

On the training curves (no table is given), A-Crab clearly beats ATAC on walker2d-random (roughly 5 to 7 versus 2 to 3 late), halfcheetah-medium-replay (about 46 to 47 versus 40 to 42) and walker2d-medium-expert (near 110 versus a noisy 90 to 100). On hopper-medium-expert both finish near 105 to 110, but ATAC falls to about 50 at 200,000 to 300,000 steps, while A-Crab dips only to 83. Smaller edges appear on walker2d-medium-replay (about 85 versus 79), halfcheetah-medium (a smooth 52 against an oscillating 46 to 51) and hopper-random (about 10 versus 7.5 from 430,000 steps, then 19 versus 15 with wide bands). Other medium and medium-replay curves end close. A-Crab loses on halfcheetah-random (flat near 2.2 against a noisy 2.5) and starts slower on halfcheetah-medium-expert. Random-dataset scores stay below 20, so those wins separate weak policies.

\question{8. What are the impacts and limitations of the contributions?}

A-Crab's contribution is mainly conceptual. The paper shows that the form of the Bellman penalty separates a tractable but slow actor-critic method from an optimal one: a linear, importance-weighted penalty succeeds where the squared one can fail. The same change removes the completeness assumption and lets the leading term of the bound depend on averaged $\ell_2$ coverage. Because A-Crab only swaps ATAC's critic penalty and its code is public, trying it costs little.

The analysis assumes finite actions, finite function classes, i.i.d. data and an idealized no-regret actor, unlike the continuous-control, neural-network experiments. Efficiency is claimed only relative to oracles: the critic must solve a min--max problem over F \ensuremath{\times} W, and the authors describe the average Bellman penalty as computationally less efficient than the squared one. The tested algorithm therefore departs from the analyzed one. Its box-shaped weight class contains the comparator's density ratio only if that ratio is bounded by $C_\infty$ everywhere, which brings back an $\ell_\infty$-type condition. The authors name future directions but state no limitations.

A-Crab offers BCA a useful contrast. In principle its weights w(s, a) reweight each transition's residual adversarially. Under the box class, however, the maximizer places $C_\infty$ on residuals of the dominant sign and zero elsewhere, so the penalty's strength rests on one grid-searched scalar, the kind of global setting BCA replaces with calibrated per-sample widths. The double-sampling issue the paper cites also implies that single-sample residuals mix transition noise with estimation error, so a scale model fit to them may inflate BCA's widths in stochastic but well-covered regions.

\question{9. How were the results evaluated or validated?}

The theorems receive explicit appendix proofs built from standard tools: a performance-difference lemma imported from Cheng et al. (2022), Bernstein's inequality for the weighted Bellman error, whose variance the $\ell_2$ norm of w controls, Hoeffding's inequality and union bounds over finite classes. The assumptions are realizability for F, an $\ell_2$ bound on W, the comparator's density ratio in W, i.i.d. data and a no-regret actor. The misspecification step cites another lemma of Cheng et al., and the ATAC lower bound uses an explicit bandit, an adapted binomial anti-concentration result and the simplification that empirical action frequencies equal population ones.

The experiments are learning curves on twelve D4RL MuJoCo v2 datasets. Figure 2 states eight seeds (0 to 7) with mean and standard deviation for its three datasets; the appendix figure does not restate the protocol, and the y-axis is labeled only as return. There is no table of final scores, aggregate metric or statistical test, no ablation of the pure average penalty or of $C_\infty$, and ATAC is the only baseline. \ensuremath{\beta} came from ATAC, and $C_\infty$ was grid-searched per dataset without an offline selection rule, which suggests tuning on the reported returns. The $1/\sqrt{N}$ versus $1/N^{1/3}$ separation is never checked numerically, even on the paper's own bandit, so the experiments support a modified algorithm rather than confirm the theorems.

\question{10. What are the ethical implications? How are they addressed?}

Being mainly theoretical, the A-Crab paper includes no ethics statement, broader impact section or limitations section. Its only societal framing is in the introduction, which presents offline RL as a way to avoid costly and dangerous trial and error and lists precision medicine, robotics and climate as application areas. The acknowledgments report NSF funding and note that the first author did the work while a visiting researcher at Meta. The code is publicly released.

The ethical questions are therefore left to the reader. A $1/\sqrt{N}$ guarantee and a robust improvement result could create false confidence in clinical or robotic deployment. They hold only under realizability, i.i.d. sampling, finite classes and coverage of the comparator policy, and none of these can be verified from logged data. The improvement bound also concerns expected return, not safety in each state, and the practical variant used in the experiments is not covered by these theorems. Because A-Crab competes only with policies the data cover, biases in the logging policy, such as historical treatment patterns, can pass directly into its output. Validation is purely simulated, and tuning $C_\infty$ against environment returns would not be possible in a real high-stakes setting.

The public code, stated seeds and per-dataset hyperparameter table support reproducibility, which is the one point the paper handles well. The other issues are not discussed.

\question{11. What themes are suggested by this research? (minimum 3)}

\textbf{Theoretical guarantees under partial data coverage.} A-Crab guarantees a $1/\sqrt{N}$ rate against any policy whose density ratio lies in the weight class, not against every policy. The leading term depends on $C^{*}_{\ell_2}$, a uniform $\ell_2$ bound on that class, while the comparator's $\ell_\infty$ ratio enters only through $B_w$ in a lower-order $1/N$ term. A policy that seldom visits poorly logged actions can therefore be learned at the optimal rate even when its $\ell_\infty$ coefficient is very large, though not infinite.

\textbf{Theory versus practice.} The analyzed and tested algorithms differ. The theory assumes finite actions, finite function classes, a weight class that contains the comparator's density ratio, and $\beta=2$. The experiments use continuous control, a box weight class reduced to a closed form, an added squared Bellman term of the kind the paper shows can block the optimal rate in ATAC, and ATAC's \ensuremath{\beta} values. The benchmark curves therefore test a heuristic inspired by the theorem rather than the theorem itself.

\textbf{Conservatism versus policy improvement trade-off.} Like ATAC, A-Crab uses relative pessimism: the critic is pessimistic about how much the policy improves on the logged actions, not about its absolute value. This yields robust policy improvement for every $\beta\ge 0$. Because that bound scales with $(\beta+1)$ and vanishes only when $\beta=o(\sqrt{N})$, heavier penalties loosen the guarantee rather than tighten it, and Theorem 1's optimality result covers only $\beta=2$.

\textbf{Adaptive, state-dependent conservatism.} The importance weights w(s, a) let the critic penalty concentrate on the state-action pairs that matter for the comparator policy, which amounts to an adversarial per-sample reweighting of Bellman residuals. In the practical version this flexibility shrinks to placing weight $C_\infty$ on residuals of the dominant sign and zero elsewhere. The authors make this simplification to avoid solving the inner maximization over w, which shows how quickly sample-level adaptivity is traded away for computational convenience.

\textbf{Implementation details and hyperparameter sensitivity.} Empirical performance depends on choices outside the theory: mixing the average and squared Bellman penalties, inheriting ATAC's \ensuremath{\beta} values from 0.062 to 64, and grid-searching $C_\infty$ over eight candidates separately for each dataset, with selected values between 1 and 50. No offline selection rule is stated for these settings and no ablation isolates them, so it is hard to tell how much of A-Crab's stability comes from the regularizer itself.

\clearpage

\bookmarked{\subsection}{subsection}{Paper 6: A Clean Slate for Offline RL (Unifloral)}

\question{1. What is the citation?}

\begin{hangref}Jackson, M. T., Berdica, U., Liesen, J., Whiteson, S., \& Foerster, J. N. (2025). A clean slate for offline reinforcement learning. \emph{Advances in Neural Information Processing Systems}, \emph{38}. \url{https://neurips.cc/virtual/2025/oral/119623}\end{hangref}
\begin{hangref}Version reviewed: Jackson, M. T., Berdica, U., Liesen, J., Whiteson, S., \& Foerster, J. N. (2025). \emph{A clean slate for offline RL} (arXiv:2504.11453v1) [Preprint]. arXiv. \url{https://doi.org/10.48550/arXiv.2504.11453}\end{hangref}

\question{2. What is the research area?}

The paper's field is offline (batch) RL, the branch of reinforcement learning, and so of machine learning, that learns policies from a fixed log of transitions. Unlike papers that propose one new learning rule, it is mainly a methodology, evaluation and software contribution. It spans model-free actor-critic methods (TD3+BC, ReBRAC, IQL, SAC-N, EDAC, CQL), model-based methods that train on synthetic rollouts (MOPO, MOReL, COMBO) and imitation baselines (BC, Decision Transformer). Its problem family is hyperparameter and policy selection under limited online access, together with fair benchmarking. Instead of championing policy constraints, value penalties, in-sample learning, ensembles or dynamics models, it treats each as a configurable component, which places it beside earlier unification and evaluation-audit work such as Rainbow, CORL and studies of online evaluation budgets.

\question{3. What is the research question?}

How should offline RL methods be defined, implemented and compared so that reported performance reflects the online interaction actually spent tuning them, and can one cleanly implemented design space that contains prior algorithms be used to discover better ones?

Several more specific questions sit beneath this. It asks which variant of the offline problem published work implicitly assumes, and how an algorithm's standing changes as the number of allowed online evaluation episodes grows from one to about two hundred. It asks whether any existing method performs well across D4RL tasks once tuning cost is counted, and whether a dozen established methods can be rewritten as compact, consistent JAX files that reproduce earlier scores and train faster. Finally, it tests two hypotheses: that replacing ReBRAC's behavior-cloning term with IQL-style advantage weighting helps (TD3-AWR), and that ReBRAC beats SAC-N or CQL as the policy learner inside a MOPO dynamics model (MoBRAC).

\question{4. What is the significance of the question?}

The practical appeal of offline RL is that useful controllers can be learned from records that already exist, in settings such as driving, robotics or medicine where trial-and-error deployment is expensive or dangerous. That appeal is undermined if published scores quietly depend on many online test runs used to pick hyperparameters, because a method that needs hundreds of evaluation episodes to choose its settings is not offline in any useful sense. The authors argue that most papers work in this tuned regime without stating the budget, so numbers from different papers blend algorithmic quality with tuning effort.

The technical stakes come from the fact that offline algorithms fight value overestimation on poorly covered actions through conservatism hyperparameters (behavior-cloning weights, penalty coefficients, expectiles, ensemble sizes) whose best values vary by dataset, which makes results sensitive to how those values are chosen. By 2025, the field had many competing methods, baselines copied from earlier papers, per-dataset hyperparameter tables (ReBRAC is a prominent case) and code bases whose incidental differences hid what each idea contributed. Without a common protocol and common code, claims of progress could not be checked.

\question{5. What are the contributions reported in this paper?}

The paper contributes on several fronts. Conceptually, it sets out a taxonomy of four offline RL settings, separated by whether policies may be tested online before deployment and whether they may adapt afterward: zero-shot deployment, online policy selection before deployment, policy selection after deployment, and offline-to-online fine-tuning. It also defines an offline RL method as an algorithm plus a fixed hyperparameter sampling range that does not change between datasets.

Methodologically, it introduces an evaluation protocol for the pre-deployment selection setting, in which a UCB bandit chooses among trained policies while observing one noisy episode return per pull, so that performance is reported as a function of the online budget. It supplies a software interface for this protocol and recommends a nine-dataset D4RL subset.

On the software and algorithmic side, it first arranges prior methods in a genealogy in which each algorithm is its parent plus a few added components (BC leads to IQL and to TD3+BC, which in turn leads to ReBRAC; SAC-N leads to LB-SAC, EDAC, CQL and MOPO, and MOPO to MOReL and COMBO). Following this tree, it releases single-file JAX reimplementations of twelve methods (BC, TD3+BC, ReBRAC, IQL, SAC-N, LB-SAC, EDAC, CQL, DT, MOPO, MOReL and COMBO), written so that each child differs minimally from its parent. It then presents Unifloral, a single implementation whose configuration space contains most of these methods, and uses it to build two new configurations, TD3-AWR and MoBRAC, without new code. All code is public.

\question{6. What is the research method?}

The research method is design-oriented and empirical: a conceptual analysis of evaluation practice, construction of a protocol and a software framework, large-scale benchmarking under that protocol, and two hypothesis-driven case studies. There is no new theory or proof.

The protocol has two stages. First, P hyperparameter settings are drawn with replacement, each with a fresh seed, from the method's declared range; the P trained policies are rolled out for many episodes and their returns stored. Second, tuning is simulated from these stored returns: K policies are subsampled, a UCB bandit pulls arms that each reveal one episode's return, and after N pulls the arm with the highest estimate is scored by its true mean, repeated B times.

Unifloral expresses algorithms through one loss family. The critic target is either an expectile-trained value v(s\ensuremath{'}), as in IQL, or the minimum over target critics at a noise-perturbed policy action, as in TD3; behavior-cloning, entropy and EDAC-style gradient-diversity terms are optional. The actor minimizes \begin{equation*}\beta_q L_q+\beta_{\mathrm{BC}}L_{\mathrm{BC}}-\beta_H\,\mathcal{H}(\pi),\end{equation*} where $L_q$ aggregates the ensemble by minimum, mean or first critic and $L_{\mathrm{BC}}$ is squared action error or negative log-likelihood, optionally multiplied by $\min(A_{\max},\exp(\eta(q-V)))$ to give advantage-weighted regression. An optional dynamics ensemble generates rollouts whose rewards are reduced in proportion to ensemble disagreement. The same space also covers architecture and optimizer settings that differ between methods, such as network depth and width, layer and observation normalization, learning rates and their schedule, batch size, critic-ensemble size and whether the actor is deterministic. ReBRAC's actor BC coefficient is sampled from [0.0005, 1.0] with batch size 1024, while IQL's AWR temperature spans [0.5, 10] with advantages capped at 100. TD3-AWR combines ReBRAC's settings with IQL's AWR settings; MoBRAC pairs MOPO's dynamics settings with ReBRAC's elsewhere.

Experiments cover nine D4RL datasets (hopper-medium, halfcheetah-medium-expert, walker2d-medium-replay, pen-human, pen-cloned, pen-expert, antmaze-large-diverse, maze2d-large, kitchen-mixed) and ten baselines, sampling hyperparameters uniformly from each original paper's ranges, with 20 policies per algorithm and dataset, K = 8 and 500 bandit repetitions. Runtime is timed over one million updates on halfcheetah-medium-expert using one L40S GPU.

\question{7. What are the results?}

No prior algorithm was best on every dataset. ReBRAC reached the top score at some budget on five of the nine datasets and IQL on four, and ReBRAC was the only method with a non-trivial score on antmaze-large-diverse (about 20). The three model-based methods never ranked above sixth of ten on any non-locomotion dataset. The authors add that they never beat BC there, but on maze2d-large Figure 3 places MOPO (about 33) and COMBO (about 17) above BC (near zero). Some curves dipped as the budget grew; for ReBRAC on hopper-medium this was traced to distractor policies, and the chance that the bandit favored one rose from roughly 0.3 to above 0.6 while it enumerated arms, with deeper dips when more policies were subsampled.

The reimplementations were compact (59 changed lines from TD3+BC to ReBRAC, versus 943 in CORL) and fast: CQL trained in 9.8 minutes against 203.9 in OfflineRL-Kit and 180.3 in CORL. A five-seed locomotion table is offered as evidence of faithful reproduction.

TD3-AWR's curve strictly dominated ReBRAC on six datasets and was dominated on one (hopper-medium, about 94 versus 101), and it dominated IQL on seven. At full budget it beat both parents on walker2d-medium-replay (about 91 versus 74 for ReBRAC and 81 for IQL) and antmaze-large-diverse (about 26 versus 20, with IQL at zero), and it led ReBRAC at small budgets on pen-expert and halfcheetah-medium-expert although its search space was larger. IQL stayed well ahead on kitchen-mixed (about 52 versus 35). MoBRAC beat the other model-based methods on six of nine datasets by the authors' count (about 51 versus near zero on pen-human), but trailed MOPO on maze2d-large and, like them, scored almost nothing on AntMaze.

\question{8. What are the impacts and limitations of the contributions?}

What the paper chiefly gives the field is a shared yardstick and toolkit. Plotting score against online budget turns hidden tuning into an explicit cost, the fixed-range definition rules out per-dataset retuning, and building TD3-AWR and MoBRAC purely through configuration shows that component-level hypotheses can be tested cheaply.

The authors flag two limits: model-based methods rely on hard-coded termination functions, which they call unrealistic, and K is fixed at 8. Only the pre-deployment setting is implemented, and the new methods were proposed and tested on the nine datasets that ranked the baselines. Rankings depend on the original papers' ranges and on UCB as the selector, and the paper never explains why IQL and ReBRAC, reported by their authors at about 47 (IQL, on older v0 data) and 54 on antmaze-large-diverse, reach only about 0 and 20 here. For ReBRAC, one likely factor is that Unifloral's configuration table (Appendix H) fixes \ensuremath{\gamma} = 0.99 and batch size 1024 on every dataset, whereas ReBRAC used 0.999 and 256 on AntMaze; ReBRAC's own ablations found these two settings cut its AntMaze average from 75.3 to 21.0 (\ensuremath{\gamma} = 0.99) and 41.0 (batch 1024). MoBRAC is compared only with model-based methods, though ReBRAC alone scores far higher on maze2d-large and antmaze-large-diverse (Figures 3 and 12). The stated average speedups (131.5\ensuremath{\times} and 74.8\ensuremath{\times}) cannot be recovered from Figure 6, whose per-algorithm ratios span about 1.3\ensuremath{\times} to 21\ensuremath{\times}.

As BCA's codebase, Unifloral treats each host's conservatism as one global scalar sampled from a fixed range, with the bandit picking among the resulting policies (ReBRAC's BC coefficient, TD3+BC's Q coefficient with BC weight 1, IQL's AWR temperature; the advantage cap stays at 100), so its protocol can test whether calibrated per-sample widths ease that tuning, which would show up as stronger scores at small budgets. BCA's calibration settings would likewise need one range across datasets, and its CQL host lies outside the unified space, which omits CQL by design.

\question{9. How were the results evaluated or validated?}

For each algorithm and dataset, the authors train 20 policies with hyperparameters drawn uniformly from the original papers' ranges, store their episodic returns, and run 500 bootstrapped UCB bandit simulations over eight sampled arms, plotting the mean and 95\% confidence interval of the chosen policy's true score against the number of single-episode pulls. Online tuning is thus explicit, not hidden. Comparisons rest on curve dominance and dataset counts, not formal hypothesis tests. Correctness is checked by five-seed means and standard deviations on fifteen locomotion datasets, speed by timing one million updates, and the distractor effect by 100,000 random policy orderings. There are no proofs and no component ablation beyond the two case studies.

Unifloral's protocol improves on common practice because it treats each evaluation as one random episode return, fixes ranges across datasets and runs every method in one consistent code base. Its weaknesses are that the intervals reflect bandit and subsampling randomness over one pool of 20 policies rather than repeated training; that strict dominance is never defined; that the evaluation episodes per policy, the training length and the hyperparameters behind the reproduction table go unreported, with no reference scores beside it; and that no reported check shows Unifloral configured as ReBRAC or IQL matching the standalone files, so TD3-AWR's gains could partly reflect implementation differences.

\question{10. What are the ethical implications? How are they addressed?}

The arXiv version of Unifloral lacks an ethics statement, a broader-impact section and a dedicated limitations section. Its ethically relevant content is a footnote calling hard-coded termination functions unrealistic, motivating examples such as autonomous search and rescue and autonomous vehicles tested with safety drivers, acknowledgments disclosing EPSRC, Amazon Web Services and Sony funding, and a pledge to publish code and run logs.

Much of the paper is itself a research-integrity intervention: undisclosed online tuning inflates claims of deployability and skews comparisons, and pricing that budget explicitly addresses the problem well. Other issues go unaddressed. In the safety-critical domains the paper cites, every pre-deployment evaluation episode carries risk, yet the protocol treats episodes as interchangeable and selects purely on mean return, ignoring constraint violations and bad tail outcomes. Distractor policies, with erratic returns and occasional high peaks, are precisely what a careless real selection process might deploy. The paper also leaves aside how behavior-regularized learners such as ReBRAC and TD3-AWR reproduce skewed or unsafe patterns in logged data, and every benchmark it uses is simulated. On the positive side, the speedups cut the energy and hardware needed for the many training runs the protocol demands and widen access for smaller labs. On balance, the paper takes research integrity seriously, despite gaps in its own reporting, while deployment safety and data bias are left to the reader.

\question{11. What themes are suggested by this research? (minimum 3)}

\textbf{Online tuning and offline model selection.} The paper's core claim is that most offline RL results silently depend on online evaluation to choose hyperparameters. It names this setting, prices it in evaluation episodes and simulates selection with a UCB bandit that sees one noisy episode return per pull. The distractor-policy finding shows that selection with few noisy evaluations can favor high-variance policies, so the selection rule matters alongside the learning algorithm.

\textbf{Evaluation rigor, reproducibility and fair comparison.} By defining a method as an algorithm plus one hyperparameter range used on every dataset, rerunning all baselines in shared code instead of copying published numbers, reporting confidence intervals over bootstrapped tuning runs, and releasing its code, the paper tries to make comparisons fair. Reporting the whole budget curve also lets readers compare methods at the budget that matches their own deployment constraints rather than at one headline number.

\textbf{Implementation details and hyperparameter sensitivity.} Unifloral expresses many algorithms as settings of one configuration space, ranging from batch size and layer normalization to behavior-cloning weights, advantage caps and ensemble size. TD3-AWR and MoBRAC arise from configuration changes alone, and the small amount of code distinguishing each algorithm from its parent suggests that much of what separates published methods is configuration rather than new principles.

\textbf{Benchmark dependence (D4RL).} Every experiment uses D4RL. The authors show that many earlier methods were tested mostly on MuJoCo locomotion and that the model-based ones, evaluated only there, fall far behind on Adroit, AntMaze, Maze2D and Kitchen. Their nine-dataset subset broadens coverage, yet the new methods were designed and judged on those same nine datasets, so benchmark dependence is reduced rather than removed.

\textbf{Behavior regularization and policy constraints.} The two strongest prior methods under the protocol, ReBRAC and IQL, both keep the policy close to dataset actions, and both new methods change how that pull is applied: TD3-AWR replaces ReBRAC's evenly weighted behavior-cloning penalty with an advantage-weighted one, and MoBRAC swaps SAC-N or CQL for ReBRAC's behavior-regularized learner inside a learned model. In Unifloral, the constraint strength is still one coefficient drawn from a fixed range.

\textbf{Computational efficiency.} Because the whole training loop is compiled in JAX, the reimplementations train far faster than PyTorch libraries; COMBO, for example, needs 22 minutes instead of almost five hours in OfflineRL-Kit on one GPU. Here efficiency is more than convenience: the protocol requires training many policies for every algorithm and dataset, so cheap training is what makes budget-aware evaluation affordable for groups with modest compute.

\clearpage

\bookmarked{\subsection}{subsection}{Paper 7: Quantile Q-Learning: Revisiting Offline Extreme Q-Learning with Quantile Regression (Quantile Q-Learning)}

\question{1. What is the citation?}

\begin{hangref}Gao, X., Li, S., Cai, Y., \& Yu, W. (2026). Quantile Q-learning: Revisiting offline extreme Q-learning with quantile regression. \emph{Transactions on Machine Learning Research}. \url{https://openreview.net/forum?id=tBKznsUimN}\end{hangref}
\begin{hangref}Version reviewed: Gao, X., Li, S., Cai, Y., \& Yu, W. (2026). \emph{Quantile Q-learning: Revisiting offline extreme Q-learning with quantile regression} (arXiv:2511.11973v2) [Preprint]. arXiv. \url{https://doi.org/10.48550/arXiv.2511.11973}\end{hangref}

\question{2. What is the research area?}

Quantile Q-Learning is a machine learning contribution to model-free offline (batch) reinforcement learning. Its algorithm is an actor-critic method that starts from the in-sample family of IQL and Extreme Q-Learning (XQL), where value functions are fitted on dataset actions and the policy is extracted by advantage-weighted regression, and then deliberately lets the value functions see actions proposed by the current policy. The problem family is temperature selection and value regularization for these methods. The work sits directly downstream of XQL and its Maclaurin-stabilized variant MXQL, borrows quantile regression from distributional RL, and adapts the mild-generalization idea of Doubly Mild Generalization (DMG). QQL is itself an offline algorithm, evaluated on D4RL and NeoRL2, so its relevance to offline RL is direct.

\question{3. What is the research question?}

Can an offline algorithm built on Extreme Q-Learning stay stable and competitive across many domains and datasets with one fixed set of hyperparameters, if its temperature is estimated from the data rather than tuned for each dataset?

The question can be unpacked into three sub-questions. The first is theoretical: if XQL's Gumbel error model is allowed a state-dependent scale, can that scale be recovered from two value functions that are themselves fixed quantiles of the Q-values under the data, so that ordinary quantile regression can learn them? The second concerns training dynamics: will a mildly generalizing value update, made more pessimistic for out-of-sample actions, offset the excess caution of purely in-sample learning and steady the critic? The third is empirical: does the resulting method, with fixed settings, match baselines tuned per dataset, and are its distributional assumptions plausible on synthetic and benchmark data?

\question{4. What is the significance of the question?}

The question targets a practical weakness of offline RL. Such algorithms should learn without touching the environment, yet many reported results depend on hyperparameters set separately for each dataset. The paper calls this impractical for XQL and MXQL, and such tuning usually means comparing candidate policies by their returns in the environment, the very interaction offline RL tries to avoid. XQL illustrates this: the authors' reproduction uses temperatures from 0.5 to 10 across D4RL datasets, and Figure 1 shows a large drop on walker2d-medium-replay when its tuned value of 5.0 is replaced by the locomotion-wide value of 2.0. Fixed settings would ease use in costly or risky domains.

XQL's temperature also matters technically, because it enters both the value loss and the policy weights, so a poor choice destabilizes both at once. Recasting it as a learnable, state-dependent quantity links extreme value theory to quantile regression and replaces a global knob with a per-state scale estimated from the data. At the time of publication, MXQL had stabilized the XQL loss but still needed a hand-set temperature, and DMG had argued that purely in-sample methods are too conservative, leaving both gaps open.

\question{5. What are the contributions reported in this paper?}

The first contribution is theoretical. The paper generalizes XQL's Gumbel error model so that the noise scale $\beta(s)$ can vary by state, borrowing the heteroscedastic critic-error assumption of Double Gumbel Q-learning. Under its two assumptions it shows that the gap between an expected optimal value $\hat{V}(s)$ and the soft value $V(s)$ equals $\omega\beta(s)$, where \ensuremath{\omega} is the Euler--Mascheroni constant, and that $V$ and $\hat{V}$ are fixed quantiles of the Q-values over dataset actions. Together these results make the temperature estimable with a pinball loss.

The second contribution is the Quantile Q-Learning (QQL) algorithm, a temperature-free version of XQL with two quantile value networks, a critic whose regression target accounts for the estimated noise offset, and an advantage-weighted policy update whose per-state temperature follows from a KL-constrained objective derived in Appendix A.

The third contribution adapts DMG's mild generalization to the quantile objectives by adding a loss on policy actions at next states with lowered quantile levels, a modification the authors call conservative estimation and justify in Appendix E.

Finally, the paper supplies an empirical study on 17 D4RL datasets and two NeoRL2 tasks under a single hyperparameter setting, a head-to-head comparison with XQL under shared settings, ablations of both regularization components and of \ensuremath{\lambda} and \ensuremath{\zeta}, and two checks of the modeling assumptions: a synthetic Gumbel simulation and a residual analysis on D4RL.

\question{6. What is the research method?}

QQL starts from a short derivation and the algorithm design that follows from it, then turns to benchmarking with ablations and assumption checks. The derivation assumes the gap between $Q^{\star}$ and the learned critic is Gumbel with scale $\beta(s)$ (Assumption 1) and $Q(s,a)-V(s)$ a negated Gumbel with the same scale (Assumption 2). It follows that $\beta(s) = (\hat{V}(s) - V(s))/\omega$, where V is the $\alpha_{1}$-quantile and $\hat{V}$ the $\alpha_{2}$-quantile of $Q(s,a)$, with $\alpha_1 = 1 - e^{-1} \approx 0.63$ and $\alpha_2 = 1 - \exp(-e^{\omega}) \approx 0.83$.

Two value networks, $V_{\psi_1}$ and $\hat{V}_{\psi_2}$, are fitted to $Q_\theta(s, a)$ on dataset pairs with the pinball loss $u(\tau - \mathbf{1}[u < 0])$ at levels $\alpha_{1}$ and $\alpha_{2}$. The critic regresses $Q_\theta(s, a) + (\hat{V}_{\psi_2}(s) - V_{\psi_1}(s))$ onto $r + \gamma \hat{V}_{\psi_2}(s')$, so the added $\omega\beta(s)$ offsets the expected Gumbel bias. The actor maximizes a log-likelihood of dataset actions weighted by $\exp[\omega(Q - \hat{V})/(\zeta(\hat{V} - V)) + \omega(Q - V)/(\hat{V} - V)]$, which is AWR with a learned per-state temperature (with $|\beta(s)|$ floored at 0.1). For mild generalization, each value loss gains a \ensuremath{\lambda}-weighted pinball term on current-policy actions at next states, with targets lowered by $\omega\beta(s')$ for these out-of-sample actions, moving the levels to $\alpha_0 = 1 - \exp(-e^{-\omega}) \approx 0.43$ for V and $\alpha_{1}$ for $\hat{V}$.

The code reuses CORL's PyTorch IQL settings (two 256-unit ReLU layers, learning rate $3 \times 10^{-4}$, batch 256), with \ensuremath{\lambda} = \ensuremath{\zeta} = 1 everywhere and one million training steps. Benchmarks are D4RL locomotion (hopper, halfcheetah, walker2d; medium, medium-replay, medium-expert), Adroit (pen, door, hammer; human and cloned), AntMaze umaze and umaze-diverse, and NeoRL2 RocketRecovery and SafetyHalfCheetah. Baselines keep per-dataset tuning: BC, TD3+BC, CQL, IQL, XQL, and MXQL on D4RL, and BC, CQL, EDAC, BCQ, TD3+BC, and XQL on NeoRL2. Further studies rerun XQL with per-domain temperatures and with one shared 2.0, sweep \ensuremath{\lambda} and \ensuremath{\zeta} from 0.25 to 4 and ablate both regularizers on three locomotion datasets, add scale-1 Gumbel noise to a random network's outputs over Gaussian-sampled actions, and fit normal and logistic laws to 8,192 $\beta(s)$-normalized Bellman residuals on four locomotion datasets.

\question{7. What are the results?}

On D4RL, QQL with fixed settings had the highest average normalized score on 10 of 17 datasets, though several leads are smaller than its own spread (94.3 \ensuremath{\pm} 1.8 against XQL's 94.2 on halfcheetah-medium-expert). The largest margins were on walker2d-medium-replay (90.2 against 83.6 for MXQL) and pen-human (128.3 against 122.1). It trailed MXQL on hopper-medium, hopper-medium-replay, and pen-cloned (77.3, 101.1, and 115.2 against 80.9, 102.7, and 117.4), fell well short on both hammer datasets (8.4 and 8.0 against MXQL's 14.7 and IQL's 11.6), narrowly lost door-cloned to IQL (8.8 against 9.0), and lost antmaze-umaze to TD3+BC (88.5 against 98.5). On NeoRL2, QQL led on RocketRecovery (83.2 against 81.1 for XQL) but scored only 60.3 on SafetyHalfCheetah, below the dataset's own 73.56 and below BC, CQL, TD3+BC, and XQL.

The shared-settings comparison is the central evidence. With one temperature for all domains, XQL's domain averages fell to 73.6 on locomotion, 37.8 on Adroit, and 67.3 on AntMaze, against 83.7, 40.6, and 83.8 when tuned, while QQL reached 85.6, 48.3, and 84.9. Moving \ensuremath{\lambda} or \ensuremath{\zeta} between 0.25 and 4.0 changed scores by at most 2.5 points on three datasets. QQL's learning curves were smoother than MXQL's, although MXQL's curves end well below its own Table 1 scores (near 80 against 110.7 on hopper-medium-expert). Removing value regularization clearly hurt hopper-medium-expert and walker2d-medium-replay, while removing conservative estimation mainly hurt hopper-medium-expert and inflated its Q-values. In the toy study, fitted Gumbel scales stayed between 0.989 and 1.065 with p-values of 0.63 to 0.87, and D4RL residuals fit a logistic law better than a normal one (KS 0.138 against 0.236 on halfcheetah-medium-replay). Training took about 4 hours per million steps, 1.6 times XQL's cost.

\question{8. What are the impacts and limitations of the contributions?}

QQL's principal impact is evidence that XQL's most sensitive hyperparameter can be estimated from data, giving a method competitive on locomotion, Adroit, and AntMaze without per-dataset tuning. Reading $V$ and $\hat{V}$ as fixed quantiles of the critic's values is a reusable idea for other in-sample methods.

The authors mention only the 1.6-fold compute cost and leave online learning for future work. $\beta(s)$ is really the spread between two quantiles of critic values over dataset actions, and it measures Bellman noise only if the two Gumbel scales coincide, which only a toy model checks directly. The paper's residuals (kurtosis 71 to 101, KS near 0.13) fit even the better logistic law only loosely, so the quantile levels are approximate. Without a non-crossing constraint, $\hat{V} - V$ can turn negative; the paper patches this with an absolute value and a 0.1 floor in the actor only, while the signed gap still shifts the critic target. The ablation text overstates Figure 3 (little effect on halfcheetah-medium, none from dropping conservative estimation on walker2d-medium-replay, no curve for dropping both), the weak SafetyHalfCheetah score goes undiscussed, and DMG is not a baseline.

Among non-conformal methods, QQL is a close precedent for BCA: it replaces a global conservatism constant with a per-state scale and acts on the same AWR weight that BCA modifies in IQL. Its scale, however, is learned jointly with the critic under a parametric model and carries no coverage guarantee, whereas BCA fits its scale to detached Bellman residuals and calibrates a frozen width with WBCP; QQL's heavy-tailed residuals favor that distribution-free route.

\question{9. How were the results evaluated or validated?}

The claims rest on benchmarks, ablations, and conditional derivations. Performance is normalized return on each benchmark's own scale, averaged over 5 seeds on D4RL and 3 on NeoRL2, with QQL shown as mean \ensuremath{\pm} an unstated spread. In Tables 1 and 2, baselines are single numbers from CORL, the MXQL code, and an OfflineRL-Lib-based XQL rerun; the NeoRL2 baseline source is not stated, and only Table 3 gives XQL \ensuremath{\pm} spreads. No confidence intervals, significance tests, episode counts, or checkpoint rules are reported. \ensuremath{\lambda} = \ensuremath{\zeta} = 1 was fixed (\ensuremath{\lambda} because it worked well empirically); absent any offline selection rule, the choice apparently rested on benchmark returns. The propositions hold only under the Gumbel assumptions, and Proposition 1's proof treats action-dependent noise as one per-state draw.

The shared-settings comparison tests the main claim directly but depends on XQL's shared temperatures (2.0 everywhere, or 2.0, 5.0, and 0.6 per domain), which are never justified (0.6 matches neither tuned AntMaze value). MXQL is never run with shared settings, and the \ensuremath{\lambda} and \ensuremath{\zeta} sweep covers only locomotion. The reports also conflict: both shared XQL rows use 2.0 on locomotion yet score 75.5 and 73.6, and Figure 1 shows a hopper-medium-expert drop though its tuned temperature is also 2.0. The toy test injects the Gumbel noise it detects.

\question{10. What are the ethical implications? How are they addressed?}

QQL's authors include no broader impact statement, ethics statement, or limitations section. The nearest material is the abstract's point that offline RL suits domains where collecting data online is dangerous or expensive (the introduction lists healthcare among reinforcement learning's applications) and the conclusion's claim that working with one fixed configuration makes the method practical for consequential real-world use.

That claim rests entirely on simulated benchmarks. The one task built around a safety constraint, SafetyHalfCheetah, is where QQL scored below the logged data and most baselines, and only normalized return is reported, so nothing is known about constraint violations. The per-state temperature depends on Gumbel assumptions that the residuals do not fully meet; a misestimated $\beta(s)$ in sparsely covered states could make the policy less cautious than intended where errors are most costly. Logged data also carry the behavior policy's biases, which an in-sample method with mild generalization largely inherits. In its favor, removing per-dataset tuning reduces the need for risky online trial runs, and compute is modest (about 4 hours per run on one RTX 4070). Reproducibility is partial: hyperparameters and the base implementation are given, but no code link is provided.

Only the fixed-hyperparameter design touches these issues, and implicitly at that; none is discussed directly.

\question{11. What themes are suggested by this research? (minimum 3)}

\textbf{Implementation details and hyperparameter sensitivity.} The paper starts from the observation that XQL's temperature must be retuned for each dataset, with reproduced values ranging from 0.5 to 10, and that one shared value costs about 16.5 points on AntMaze. QQL responds by estimating the temperature and fixing \ensuremath{\lambda} and \ensuremath{\zeta} at 1, and a sweep from 0.25 to 4 barely moves the scores. Tuning is relocated rather than abolished, since \ensuremath{\zeta} still scales the learned temperature, but the new knobs appear far less sensitive.

\textbf{Adaptive, state-dependent conservatism.} Instead of one global temperature, QQL learns $\beta(s)$ for every state from the gap between two quantile value functions. States where the critic's values spread more widely across dataset actions get a larger temperature, flatter advantage weights, and so a stronger pull toward the data. Under the model, the same scale also fixes how far value targets for out-of-sample actions are lowered, which the algorithm implements through shifted quantile levels, so one quantity governs caution in both actor and critic.

\textbf{Uncertainty quantification and calibration.} The paper treats $\beta(s)$ as a per-state uncertainty estimate grounded in a heteroscedastic Gumbel error model, although what the networks actually measure is the spread of critic values across dataset actions. The model is checked by a synthetic goodness-of-fit study and by Bellman residuals that a logistic law fits better than a Gaussian, though far from exactly. No calibration or coverage property is attached to the learned scale, which separates it from distribution-free approaches.

\textbf{Conservatism versus policy improvement trade-off.} QQL inherits the caution of in-sample learning and then relaxes it with DMG-style mild generalization on actions sampled from the current policy. To stop this relaxation from importing extrapolation error, targets for those actions are lowered by the estimated noise offset. The balance between generalizing beyond the data and staying pessimistic is thus set through shifted quantile levels and the weight \ensuremath{\lambda}, and the ablations show that the regularization helps on some datasets more than others.

\textbf{Theory versus practice.} QQL's quantile levels of about 0.63, 0.83, and 0.43 follow exactly from Gumbel algebra, but the working algorithm departs from that ideal. The estimated scale has to be forced positive and floored at 0.1 in the actor, and real D4RL residuals are far heavier-tailed than the logistic shape the model implies. The derivation motivates the design, yet the case for the method rests mainly on benchmark returns and ablations rather than on the theory.

\clearpage

\bookmarked{\subsection}{subsection}{Paper 8: UNIQ: Conformal Calibration for Adaptive Conservatism in Offline Reinforcement Learning (UNIQ)}

\question{1. What is the citation?}

\begin{hangref}Upadhyay, A. (2026, July 11). \emph{UNIQ: Conformal calibration for adaptive conservatism in offline reinforcement learning} [Poster presentation]. ICML 2026 Workshop on Decision-Making from Offline Datasets to Online Adaptation: Black-Box Optimization to Reinforcement Learning, Seoul, South Korea. \url{https://icml.cc/virtual/2026/83644}\end{hangref}
\begin{hangref}Version reviewed: Upadhyay, A. (2026). \emph{UNIQ: Conformal calibration for adaptive conservatism in offline reinforcement learning} (arXiv:2606.07592v1) [Preprint]. arXiv. \url{https://doi.org/10.48550/arXiv.2606.07592}\end{hangref}

\question{2. What is the research area?}

UNIQ is aimed at offline (batch) reinforcement learning, a subfield of machine learning. It is a model-free actor-critic method built on Implicit Q-Learning (IQL), which learns values in-sample by expectile regression and extracts a policy by advantage-weighted regression. The problem family is pessimism under distribution shift: deciding how far to distrust value estimates in states and actions that the logged data cover poorly. Earlier work either applies one fixed conservatism coefficient everywhere (CQL, IQL, TD3+BC), obtains adaptive pessimism from large critic ensembles (SAC-N, EDAC), or makes penalties state-dependent through density or dynamics models (COMBO, MOPO, MOReL). UNIQ stays model-free, using a small value ensemble for per-state uncertainty and borrowing split conformal prediction from distribution-free statistics to put that uncertainty on a common scale across tasks.

\question{3. What is the research question?}

The question UNIQ poses is this: can IQL's single, global expectile be replaced by a per-state expectile driven by conformally normalized ensemble uncertainty, so that conservatism tightens where the dataset is sparse and loosens where it is dense, without paying the memory cost of large critic ensembles?

Three narrower questions follow. The first is whether the adaptive version beats IQL across the nine D4RL MuJoCo locomotion datasets, and whether its gains concentrate on datasets with uneven coverage, such as the medium-replay mixtures. The second is whether each component is needed: does dividing ensemble disagreement by a conformal quantile of Bellman residuals matter, does a state-dependent expectile help beyond a fixed one, and how does an extra disagreement penalty \ensuremath{\kappa} in the Q-target interact with it? The third is whether the method offers a good balance between benchmark score and GPU memory.

\question{4. What is the significance of the question?}

Most widely used offline RL algorithms expose a single global setting for conservatism, such as CQL's \ensuremath{\alpha}, TD3+BC's \ensuremath{\alpha} or IQL's \ensuremath{\tau}, and the best value changes from dataset to dataset. Choosing it is difficult in a truly offline setting, because checking a policy normally means running it. A rule that derives the amount of caution from the data would reduce this dependence and, in principle, place caution near the edges of the data rather than spreading it uniformly.

On the technical side, the paper targets a known weakness of ensemble disagreement as an out-of-distribution signal: its raw size depends on the task and its reward scale, so a fixed penalty transfers poorly between environments. Calibrating it on held-out data is one possible remedy. Pessimism theory reviewed in the appendix also implies that the need for pessimism varies by state. Among the baselines the paper compares against, all taken from the 2023 CORL benchmark, the top MuJoCo averages belong to the large-ensemble EDAC (92.9) and the heavily tuned ReBRAC (89.7), and the author describes conformal methods for RL uncertainty as underexplored. A lightweight, calibrated alternative was therefore a timely goal, especially on modest hardware.

\question{5. What are the contributions reported in this paper?}

The paper's principal contribution, the UNIQ algorithm, extends IQL with three linked pieces. The first is a multi-expectile value ensemble: several value networks trained at each of three expectile levels (0.5, 0.7 and 0.9), whose disagreement at the middle level gives an epistemic uncertainty estimate and whose upper-minus-lower gap serves as a diagnostic for the spread of returns. The second is a split conformal step that scores held-out transitions by the absolute Bellman residual of the ensemble mean and uses the resulting quantile to convert raw disagreement into a unitless uncertainty ratio. The third is a sigmoid controller that maps this ratio to a state-specific expectile, paired with a disagreement-penalized value target for the Q-function.

The appendix adds a formal account of the design: properties of expectiles, a link between the penalized target and lower-confidence-bound pessimism, a restatement of the split conformal coverage theorem, a proposition interpreting when the uncertainty ratio exceeds one, a proposition on the monotonicity and limits of the controller, and an argument that IQL is a special case. On the empirical side, the paper supplies a nine-task D4RL comparison against eight published baselines that deliberately keeps the rows where UNIQ trails, an analysis of memory and per-step cost, and an ablation study of calibration, adaptivity, ensemble size and the penalty coefficient. The author presents the work as a new mechanism, potentially transferable to other backbones, rather than as a state-of-the-art claim.

\question{6. What is the research method?}

Algorithm design leads the work, backed by a light formal analysis and by empirical benchmarking with ablations. UNIQ keeps IQL's three learners: a value network fit by asymmetric squared loss on Q(s, a) \ensuremath{-} V(s), a Q-function fit by TD regression, and a policy trained by advantage-weighted regression with weight exp(\ensuremath{\beta}\ensuremath{\pi}(Q \ensuremath{-} V)), \ensuremath{\beta}\ensuremath{\pi} = 3.

Three ensemble members (Nv = 3) are trained at each fixed level 0.5, 0.7 and 0.9, and \ensuremath{\sigma}(s) is the standard deviation of the 0.7-level heads. Every 5,000 steps, each transition in a 20 percent held-out calibration split is scored by $\alpha = |r + \gamma \bar{V}(s') - \bar{V}(s)|$ using the 0.7-level ensemble mean, and $\hat{q}$ is set to the $\lceil (1 - \delta)(n + 1) \rceil$-th smallest score with \ensuremath{\delta} = 0.1. The normalized uncertainty is $u(s) = \sigma(s)/(\hat{q} + \epsilon)$, and the expectile becomes $\tau(s) = \tau_{\min} + (\tau_{\max} - \tau_{\min}) \cdot \operatorname{sigmoid}(-\beta_\tau(u(s) - 1))$, with $\tau_{\min} = 0.5$, $\beta_\tau = 5$, and $\tau_{\max}$ equal to 0.90 or 0.95. Q-targets use the penalized value $\bar{V}(s') - \kappa\sigma(s')$. In Algorithm 1, $\tau(s)$ sets the expectile of a separate primary value network that supplies V in the policy's advantage, updated after the ensemble and before the Q-function and policy; Appendix B instead scales each ensemble level's expectile by $\tau(s)$ and uses the ensemble mean in the advantage.

Config A (\ensuremath{\kappa} = 0, $\tau_{\max}$ = 0.95) is used for halfcheetah-medium-expert and hopper-medium-replay, and Config B (\ensuremath{\kappa} = 0.5, $\tau_{\max}$ = 0.90) for the other seven tasks. Training uses one million Adam steps (learning rate $3 \times 10^{-4}$, batch 256, discount 0.99, target averaging 0.995) and three-layer networks of width 256. Experiments use the nine D4RL MuJoCo v2 datasets (three robots, each with medium, medium-replay and medium-expert data). Scores for eight baselines (BC, TD3+BC, CQL, IQL, EDAC, ReBRAC, SAC-N and Decision Transformer) come from CORL summaries rather than reruns. UNIQ uses seeds 0 to 2 on A100 20 GB MIG partitions. Ablations use seed 0 on four tasks (halfcheetah-medium, hopper-medium, hopper-medium-replay, walker2d-medium), and Appendix D runs the mechanism variants at a fixed \ensuremath{\kappa} = 1.

\question{7. What are the results?}

With three seeds, UNIQ averages 85.2 over the nine tasks against 81.6 for IQL and beats IQL on every task. The gains are small on HalfCheetah (+0.1 to +1.5) and bigger for the other two robots, peaking at +8.1 on hopper-medium and +7.2 on walker2d-medium-replay. By dataset type, the mean gain is about +4.4 on medium, +4.3 on medium-replay and +1.9 on medium-expert, so replay data do not benefit more than medium data. UNIQ tops every listed method on hopper-medium-replay (101.6), hopper-medium-expert (111.8) and walker2d-medium-replay (89.4). It ranks third on average, behind EDAC (92.9) and ReBRAC (89.7) but ahead of SAC-N (83.5) and the rest. It trails badly on halfcheetah-medium (48.9 versus 68.2 for SAC-N) and hopper-medium (75.6 versus 102.3 for ReBRAC). Measured peak memory is 250 MB for UNIQ and 530 MB for IQL; estimates for SAC-N, ReBRAC and EDAC are 700, 1,200 and 2,500 MB.

In the single-seed ablations, Table 3 lists 59.4 for the full method, 49.1 with a fixed expectile and 45.2 without conformal normalization. The full row uses per-task \ensuremath{\kappa}, however, while Appendix D ran these variants at \ensuremath{\kappa} = 1, where the adaptive model averages 48.0. At that matched setting, fixing the expectile barely changes the average (losing on hopper-medium and walker2d-medium, winning on hopper-medium-replay), and removing normalization costs about three points. The penalty coefficient has the largest effect: uniform \ensuremath{\kappa} values from 0 to 2 move the average from 57.7 to 45.3. Five ensemble members score 47.2, far below three under per-task \ensuremath{\kappa} (58.9) but level with three at \ensuremath{\kappa} = 1 (47.5). A single member, which switches off both the adaptive expectile and the penalty, still averages 59.1.

\question{8. What are the impacts and limitations of the contributions?}

UNIQ demonstrates that conformal calibration can be added cheaply to a standard offline learner, giving IQL state-dependent conservatism and a 3.6-point gain at low memory, with candid reporting of losses.

The author concedes that $\hat{q}$ only sets a task-wide scale and certifies no individual state, that coverage holds only at each recalibration, that HalfCheetah gains little, and that \ensuremath{\kappa} must be chosen per task. All evidence also comes from MuJoCo locomotion, omitting AntMaze and Adroit, where coverage gaps bite harder. The ablation's full method uses \ensuremath{\kappa} = 0 on hopper-medium-replay while the mechanism variants use \ensuremath{\kappa} = 1, confounding calibration with \ensuremath{\kappa}, and a one-member ensemble with no uncertainty signal nearly matches the full method (59.1 versus 59.4). Some claims clash with the paper's evidence: Appendix E's memory ratio contradicts Table 2, Appendix B.7 says UNIQ beats ReBRAC, and Figure 2 shows no late recovery on hopper-medium-replay.

UNIQ is BCA's closest prior work: both score held-out Bellman residuals conformally to set per-sample conservatism. UNIQ, however, computes its quantile on residuals but divides ensemble spread by it, so coverage does not carry over to $u(s)$; BCA calibrates residuals relative to its own fitted scale model, so its width and per-sample signal describe the same quantity. UNIQ also alters value learning (V's expectile and the Q-target) and recalibrates during training; BCA freezes a WBCP width and applies it to actor or critic terms in four hosts, in IQL to the actor's capped advantage weight. UNIQ's \ensuremath{\kappa} sensitivity suggests testing whether BCA's width removes per-task tuning.

\question{9. How were the results evaluated or validated?}

UNIQ is judged mainly by D4RL normalized return on nine MuJoCo datasets. UNIQ is averaged over seeds 0 to 2, but every baseline, IQL included, is copied from CORL reports, although Figure 2's IQL curves and measured IQL memory suggest the author ran IQL; the central gain over IQL is therefore not a matched comparison. The number of evaluation episodes is not stated, results are means without intervals or tests, and the seed-level breakdown promised for Appendix D never appears, since that appendix holds only seed-0 ablations. The ablations are inconsistently labeled: the full-method row matches the \ensuremath{\kappa} = 1 run on the three tasks meant to use \ensuremath{\kappa} = 0.5, and Tables 3 and 6 give different \ensuremath{\kappa} settings for the ensemble-size runs. Baseline memory and the 4.6-fold speedup over EDAC (51 versus 11 forward passes) are analytic estimates, not timings. Hyperparameters are said to be chosen offline from returns on the calibration split, yet Appendix C gives only the final assignment and Appendix D cites online-evaluated ablations as its basis. The conformal component rests on a restated split conformal theorem under exchangeability; realized coverage, the distributions of $u(s)$ and $\tau(s)$, and $\hat{q}$ over training are never reported, and Figure 2 lacks the late recovery its caption describes. The 3.6-point gain is plausible but weakly established.

\question{10. What are the ethical implications? How are they addressed?}

Ethics receives no dedicated treatment in UNIQ: there is no ethics statement, broader impact section or dedicated limitations section. Its Discussion and Scope section comes closest, covering calibration lag, sensitivity to \ensuremath{\kappa} and weak HalfCheetah gains, together with the choice to keep underperforming rows in the main table. The introduction motivates offline RL by noting that collecting data online can be costly, risky or ethically constrained, but this point is not developed.

The language of distribution-free, finite-sample guarantees could suggest a safety assurance the method does not provide: the guarantee concerns Bellman residuals on held-out logged transitions, not the uncertainty ratio at states the learned policy will reach, and the controller becomes more optimistic wherever ensemble members agree, which could be harmful in medical or industrial use if they share the same error. Validation is entirely in simulated locomotion. The policy also inherits any biases in the behavior data and is least cautious in the regions those data dominate. No code release is mentioned, which limits reproducibility. The small memory footprint does widen access for researchers with modest hardware.

UNIQ responds to these issues mostly indirectly, through its frank scope statement. Appendix A.4 does concede that states reached by the learned policy may differ from the calibration data and names weighted conformal prediction as a remedy, but the paper never discusses deployment harms or misuse.

\question{11. What themes are suggested by this research? (minimum 3)}

\textbf{Adaptive, state-dependent conservatism.} UNIQ starts from the view that one expectile for every state is too cautious where data are dense and too lenient where they are sparse, and it replaces IQL's global \ensuremath{\tau} with a $\tau(s)$ set from calibrated uncertainty. The paper presents the unevenly covered medium-replay datasets as its clearest win, yet its average gain over IQL there (+4.3) is no larger than on the medium datasets (+4.4), and a per-task global \ensuremath{\kappa} remains, so the adaptivity is only partial.

\textbf{Uncertainty quantification and calibration.} The paper treats disagreement among a small ensemble of value heads as an epistemic signal and the gap between high and low expectile heads as an aleatoric diagnostic. It recognizes that raw disagreement has no common scale across tasks and rescales it against held-out Bellman residuals. At a matched \ensuremath{\kappa} = 1, removing that rescaling lowers the four-task ablation average only from 48.0 to 45.2; the larger headline drop from 59.4 also reflects a change in \ensuremath{\kappa}.

\textbf{Distribution-free (conformal) guarantees.} Split conformal prediction supplies a quantile of absolute Bellman residuals with marginal finite-sample coverage, assuming exchangeable scores. The paper admits that this guarantee does not extend to individual states: $\hat{q}$ acts as a task-level scale, and coverage holds only at each recalibration of a changing ensemble. This gap between what is guaranteed and how it is used recurs whenever conformal methods are placed inside RL training.

\textbf{Implementation details and hyperparameter sensitivity.} Although UNIQ aims to automate conservatism, its results hinge on the pessimism coefficient \ensuremath{\kappa}, which the ablations show matters more than either calibration or the adaptive expectile; hopper-medium-replay collapses once \ensuremath{\kappa} reaches 1. Two hand-assigned configurations, differing in \ensuremath{\kappa} and $\tau_{\max}$, cover the nine tasks, and the author names learning a state-dependent \ensuremath{\kappa}(s) as the main next step. In practice \ensuremath{\kappa} takes over \ensuremath{\tau}'s role as the setting that needs tuning.

\textbf{Computational efficiency.} The paper places its contribution on the trade-off between score and hardware cost rather than on raw score, arguing that nine small value heads and a periodic quantile computation cost far less than scaling up critics as EDAC does. Because only UNIQ's and IQL's memory were measured, and the appendix gives a UNIQ-to-IQL memory ratio that contradicts the main table, the efficiency claim is weaker than the headline tenfold saving suggests.

\textbf{Evaluation rigor, reproducibility and fair comparison.} UNIQ is unusually open about the tasks where it loses, yet its evidence is thin by current standards: three seeds, baselines copied from CORL rather than rerun, single-seed ablations, and several numbers that disagree between tables and text. The paper shows how strongly the credibility of D4RL claims depends on matched protocols, adequate replication and consistent reporting.

\clearpage

\bookmarked{\subsection}{subsection}{Paper 9: Bayesian Flow Networks for Offline Trajectory Planning (BFN Planning)}

\question{1. What is the citation?}

\begin{hangref}Killingberg, L., \& Langseth, H. (2026). \emph{Bayesian flow networks for offline trajectory planning} (arXiv:2608.25163v1) [Preprint]. arXiv. \url{https://doi.org/10.48550/arXiv.2608.25163}\end{hangref}

\question{2. What is the research area?}

BFN-RL comes from the part of offline RL, itself a branch of reinforcement learning within machine learning, that treats decision making as conditional sequence modeling rather than actor-critic or value-based learning. Its problem family is generative trajectory planning, alongside the Trajectory Transformer, Diffuser and Decision Diffuser. The authors' niche is the generator: they replace the Gaussian diffusion model of Diffuser-style planners with a Bayesian Flow Network (BFN) and compare it with categorical diffusion on discrete state spaces. The method is neither cleanly model-free nor model-based. It learns no forward dynamics, reward model or value function, yet it plans by generating future state trajectories and replanning as it acts, placing it closest to broadly model-based planning. Having no bootstrapped critic, it avoids the out-of-distribution value overestimation that troubles value-based offline RL.

\question{3. What is the research question?}

Can Bayesian Flow Networks serve as a single generative backbone for offline trajectory planning that handles categorical and continuous state spaces natively while remaining competitive with diffusion-based planners in each setting?

The question has three parts. First, can a categorical BFN, conditioned on return or goal, with the current state imposed by inpainting and actions supplied by a learned inverse-dynamics model, generate plans that execute in grid worlds without hand-written transition rules? Second, does the BFN generator match or beat categorical diffusion when every other part of the pipeline is held fixed, so that any gap reflects the generator alone? Third, does the continuous version of the same recipe reach D4RL MuJoCo scores comparable to the Decision Diffuser and to value-based methods? A further hypothesis, that inferring actions with inverse dynamics helps BFNs as it helps diffusion, is stated but explicitly deferred to future work.

\question{4. What is the significance of the question?}

Many planning problems have discrete structure, such as grid navigation, key-and-door puzzles and box pushing, and some mix discrete and continuous variables. Diffusion planners such as Diffuser and the Decision Diffuser, which the paper treats as the leading sequence-model planners, are built for continuous vectors, so discrete states call for a separate categorical formulation such as discrete diffusion (autoregressive models like the Trajectory Transformer sidestep this by tokenizing each dimension). A generator that covers both modalities in one probabilistic framework would simplify the design of planners for such domains.

The question also matters for the field's direction. Sequence-model planners drew interest because they avoid the out-of-distribution value overestimation of bootstrapped critics, and offline learning is motivated by domains where exploration is unsafe. BFNs had been reported to work well on discrete data in generative modeling, while the independent Guided-BFNs work on planning considered only continuous state-action trajectories. Testing whether a newer generative family transfers to decision making, against a matched discrete diffusion baseline, shows whether the generator choice in these planners deserves attention. Although the preprint appeared in 2026, its account of the state of the art rests on work published between 2020 and 2024.

\question{5. What are the contributions reported in this paper?}

The paper lists three contributions. The first is BFN-RL itself, a planning framework in which the generator evolves the parameters of a distribution over a state trajectory rather than denoising a noisy trajectory, so the same machinery applies to discrete and continuous spaces. Following the Decision Diffuser, it models only states and feeds the return directly into the network in classifier-free fashion (two of the discrete tasks condition on a goal instead).

The second listed contribution is a discrete planner that pairs a categorical BFN over state sequences with a learned inverse-dynamics model that infers the action linking each pair of successive generated states. Several technical pieces described in the method section, though not listed as separate contributions, make this planner work. The authors adapt diffusion-style inpainting to BFN sampling: for observed variables, the noisy message drawn at each step is overwritten by the clean message for the known value, with a categorical version (Algorithm 2) and a Gaussian one (Algorithm 3). They also reparameterize the endpoint of the quadratic accuracy schedule as C = K\ensuremath{\beta}(1), so that the terminal separation between correct and incorrect categories does not depend on the number of categories, and they add an auxiliary cross-entropy term on the clean category.

The third is a cross-domain empirical study: five discrete tasks with a categorical-diffusion baseline matched to BFN-RL in everything except the generator, plus D4RL MuJoCo locomotion. The paper claims no theoretical results and does not mention a code release.

\question{6. What is the research method?}

BFN-RL is an engineering study: the authors design a planner and benchmark it, without theoretical analysis. In the categorical planner, each state variable over the horizon has a categorical parameter vector \ensuremath{\theta} that starts uniform at 1/K. At step i a temporal U-Net receives \ensuremath{\theta}, the BFN time and the return or goal, and outputs a distribution from which a category k is drawn; a noisy message $y \sim \mathcal{N}(\alpha(K e_k - 1), \alpha K I)$ is sampled, and \ensuremath{\theta} is updated by the Bayesian rule $\theta \propto e^{y} \odot \theta$, with $\alpha = \beta(1)(2i - 1)/n^2$ from the schedule $\beta(t) = \beta(1)t^2$. For observed variables, the message is replaced by the clean value $\alpha(K e_c - 1)$, pinning the current state. Training minimizes the continuous-time BFN loss plus a 0.5-weighted auxiliary cross-entropy that asks the network to name the true category directly. An inverse-dynamics MLP with two 256-unit hidden layers then maps consecutive generated states to actions. Discrete training uses Adam (learning rate $3\times10^{-4}$, batch 64 to 256, linear warmup then cosine decay); the number of training updates is unreported.

The discrete testbed comprises three MiniGrid tasks (Empty-Random-6\ensuremath{\times}6 with 2,000 successful trajectories, DoorKey-8\ensuremath{\times}8 with 5,000 and BlockedUnlockPickup with 3,000, the last two adding episodes where random actions derail the agent before an expert recovers), slippery FrozenLake-8\ensuremath{\times}8 (30,000) and two-box Sokoban-7\ensuremath{\times}7 (10,000). MiniGrid tasks are return-conditioned and the other two goal-conditioned. Horizons are 16, 32, 32, 64 and 32; sampling uses 12 or 16 updates, temperature 0.7 or 0.9 and one candidate. Empty-Random, FrozenLake and Sokoban replan after every action, while the other two execute up to 16 actions and replan early on any mismatch with the plan. The baseline is a categorical diffusion model with uniform corruption and a cosine schedule, whose network estimates the uncorrupted trajectory, using the same update count.

For continuous control, a Gaussian BFN with $\sigma_1 = 0.01$ is trained for $2\times 10^{6}$ steps (Adam, batch 256, learning rate $5\times 10^{-5}$, horizon 20, weight averaging with decay 0.9995) and sampled with a first-order solver in 10 steps, guidance weight 1.2 and temperature 0.5. It is compared on nine D4RL MuJoCo datasets with eight published baselines.

\question{7. What are the results?}

In the matched discrete comparison, neither generator dominates. On Empty-Random both solve nearly every episode, but diffusion reaches the goal faster (mean return 0.941 against 0.895, success 100.0\% against 99.7\%). On DoorKey the BFN is slightly ahead, with 100.0\% success and a return of 0.934 against 99.3\% and 0.916. BlockedUnlockPickup is hard for both: diffusion averages 37.0 \ensuremath{\pm} 13.5\% success and BFN-RL 32.7 \ensuremath{\pm} 8.3\%, a gap smaller than the spread across training seeds. FrozenLake is essentially a tie (BFN-RL 86.7 \ensuremath{\pm} 1.3\% against 85.5 \ensuremath{\pm} 1.0\% for diffusion). Sokoban is the clearest win: BFN-RL reaches 68.7 \ensuremath{\pm} 2.5\% success against 59.0 \ensuremath{\pm} 2.6\%, a 9.7-point gain in mean success, and it leads in all three training runs.

On D4RL MuJoCo, BFN-RL averages 77.0 normalized points, behind the Decision Diffuser (81.8) and the Trajectory Transformer (78.9), level with IQL (77), just below CQL (77.6), and ahead of Diffuser (75.3), DT (74.7), MOReL (72.9) and BC (51.9). It posts the best score in the table on halfcheetah-medium-expert (97.2 \ensuremath{\pm} 0.1, against 90.6 for the Decision Diffuser) and ranks second on hopper-medium-expert (110.3). Its weak spots are walker2d-medium, where 66.9 is the lowest entry in the row and below behavior cloning's 75.3, and walker2d-medium-replay (58.5 against 75 for the Decision Diffuser). Against the Decision Diffuser specifically, BFN-RL wins on only two of the nine datasets, both HalfCheetah (medium-expert, and medium-replay at 42.5 against 39.3). Its scores on hopper-medium and hopper-medium-replay and on all three Walker2d datasets carry standard errors of 4.8 to 7.1 points, compared with 0.1 to 0.5 on HalfCheetah.

\question{8. What are the impacts and limitations of the contributions?}

BFN-RL contributes a proof of concept that BFNs can drive offline planners, including a categorical one, together with a usable inpainting rule for BFN sampling and a matched-baseline protocol that isolates the generator's effect. The Sokoban result hints that BFNs may help on harder combinatorial tasks.

By the authors' own account, BFN-RL does not beat categorical diffusion across the board, the value of inverse dynamics for BFNs remains untested, and mixed discrete-continuous planning is only a future extension. No component is ablated (cross-entropy weight, C endpoint, temperature, step counts), and task-specific settings come without any account of how they were chosen. The discrete datasets hold only successful trajectories, so return conditioning has little poor behavior to separate from good. The discrete study also omits behavior cloning and autoregressive token models such as the Trajectory Transformer, which already handle discretized sequences, and its diffusion baseline uses uniform corruption, whereas the discrete-diffusion papers it cites generally did better with absorbing, mask-based corruption. In continuous control BFN-RL trails the Decision Diffuser by 4.8 points on average, and planning cost is never reported.

BFN-RL has no critic and hence no Bellman residuals, so BCA's calibration machinery cannot attach directly. The closest analogue is its sampling temperature (plus a guidance weight of 1.2 in continuous control), set by hand for each task, which plausibly shapes how strongly plans commit to high-return or high-density regions. The paper does not analyze it this way, but it plays the part of a hand-set conservatism coefficient, the kind of quantity BCA calibrates per sample in its actor-critic hosts.

\question{9. How were the results evaluated or validated?}

The discrete claims are tested with a matched protocol: the BFN and categorical-diffusion generators share training data, temporal U-Net, state representation, conditioning, inpainting, inverse dynamics, horizon and execution rules, and both run with the same environment and sampling seeds. Each method is trained with three seeds and evaluated on 100 episodes per seed (200 for FrozenLake), with the mean and sample standard deviation of success rate and, for MiniGrid, mean return. The continuous claims rest on D4RL normalized scores over three seeds with standard errors, while every baseline figure is copied from the Decision Diffuser paper. Neither set of claims is backed by ablations or theory.

The matched design is the strongest part of the validation, because it attributes differences to the generator rather than to surrounding engineering, and the per-seed check on Sokoban strengthens that claim. Matching has a cost, though: the diffusion baseline inherits the BFN's 12 or 16 update steps and task settings, and the paper does not say whether it was tuned on its own. Three seeds, no significance tests, different dispersion measures in the two tables and a D4RL average without an error bar limit statistical confidence. Hyperparameter selection is never described; a BlockedUnlockPickup width sweep (48 or 96) is mentioned without saying which width produced the reported numbers or how it was chosen.

\question{10. What are the ethical implications? How are they addressed?}

BFN-RL comes without an ethics statement, broader-impact discussion or limitations section. Its only declaration states that the Norwegian University of Science and Technology funded the work internally. The introduction raises safety only as motivation, presenting learning from logged data as attractive for applications like self-driving cars and clinical treatment, where trial-and-error exploration would be dangerous.

That motivation raises issues the paper does not examine. A generative planner can propose state sequences the system cannot realize, and the inverse-dynamics model may then output actions with unpredictable consequences, a risk that grows in stochastic settings like slippery FrozenLake. Conditioning on returns or goals outside the data may push generation into unfamiliar regions without any uncertainty signal or safety guarantee. The experiments never leave simulation, and the discrete datasets are curated successful runs, so behavior under biased or low-quality logs is unknown. Reproducibility is mixed: hyperparameters are listed, but no code release is mentioned, baseline numbers are borrowed, and compute is not reported. Dual-use risk appears low given the toy and locomotion tasks.

Some design choices could still mitigate these concerns, namely replanning after every action (Empty-Random, FrozenLake, Sokoban) or once the agent leaves the plan (DoorKey, BlockedUnlockPickup), and a matched baseline that supports fair comparison, but the authors present these as planning and evaluation details rather than safeguards, and the safety motivation is asserted rather than tested.

\question{11. What themes are suggested by this research? (minimum 3)}

\textbf{Generative and sequence models for decision making.} BFN-RL frames offline RL as sampling state trajectories from a learned generative model conditioned on a desired return or goal, then recovering actions with inverse dynamics. It inherits this framing from Diffuser and the Decision Diffuser and presents it as an alternative to value-based methods that overestimate out-of-distribution values. The paper's novelty lies in the generator, showing that the planning recipe survives a change of generative family.

\textbf{Unified modeling of discrete and continuous spaces.} The central motivation is that Gaussian diffusion fits continuous vectors but not categorical states. Because BFNs update distribution parameters, the same framework covers grid-world states with a categorical model and MuJoCo states with a Gaussian one, and the authors propose mixing both types within one trajectory. The conclusion argues that competitive performance in both settings, with a consistent edge on Sokoban, justifies BFNs as one shared planner for either data type, even though the score gains themselves are mixed.

\textbf{Evaluation rigor, reproducibility and fair comparison.} The discrete study holds every component except the generator fixed, which is a careful way to isolate the generator's effect. The continuous study, by contrast, copies baseline numbers from the Decision Diffuser paper instead of rerunning them, so its comparison mixes codebases and protocols. Both studies rest on three training seeds without significance tests, which limits how firmly either study can rank the methods it compares.

\textbf{Implementation details and hyperparameter sensitivity.} Making the categorical planner work required several engineering choices: rescaling the accuracy schedule endpoint by the number of categories, adding a cross-entropy loss weighted at 0.5, and setting temperatures, step counts, horizons and network widths separately for each task. None of these choices is ablated, and their selection is not explained, so it remains unclear how sensitive the reported results are to them.

\textbf{Dataset quality and composition.} The discrete experiments train only on successful trajectories, and the DoorKey and BlockedUnlockPickup sets add episodes in which random actions push the agent off course before an expert recovers, so they resemble imitation more than learning from mixed-quality logs. On D4RL the largest gaps to the Decision Diffuser appear on walker2d-medium and walker2d-medium-replay, yet BFN-RL leads on halfcheetah-medium-replay, so the paper gives no clear evidence on how data quality affects BFN planners.

\clearpage

\bookmarked{\subsection}{subsection}{Paper 10: Conformal Prediction Beyond the Horizon: Distribution-Free Inference for Policy Evaluation (Conformal Policy Evaluation)}

\question{1. What is the citation?}

\begin{hangref}Gan, F., Lu, Y., Zhang, Y., \& Liu, Y. (2025). Conformal prediction beyond the horizon: Distribution-free inference for policy evaluation. \emph{Advances in Neural Information Processing Systems}, \emph{38}. \url{https://neurips.cc/virtual/2025/poster/118031}\end{hangref}
\begin{hangref}Version reviewed: Gan, F., Lu, Y., Zhang, Y., \& Liu, Y. (2025). \emph{Conformal prediction beyond the horizon: Distribution-free inference for policy evaluation} (arXiv:2510.26026v1) [Preprint]. arXiv. \url{https://doi.org/10.48550/arXiv.2510.26026}\end{hangref}

\question{2. What is the research area?}

Conformal policy evaluation draws on statistics as much as machine learning; its home is distribution-free predictive inference (conformal prediction) applied to sequential decision making. Within reinforcement learning it addresses policy evaluation, both on-policy and off-policy, for infinite-horizon discounted Markov decision processes, and it builds on model-free distributional RL, using quantile temporal-difference (QTD) learning as the base return model. The problem family is uncertainty quantification for returns: producing a prediction interval for the random discounted return from a new state rather than a point estimate of its mean. It extends earlier conformal methods for finite-horizon MDP trajectories, weighted conformal off-policy evaluation, and conformal off-policy prediction for contextual bandits to long horizons. Offline RL enters through off-policy evaluation from a fixed logged dataset when the target policy differs from the behavior policy.

\question{3. What is the research question?}

How can one build prediction intervals for the infinite-horizon discounted return of a target policy, using only finite logged trajectories, that reach a chosen coverage level without depending on a correctly specified return-distribution model, in both on-policy and off-policy settings?

It splits into sub-questions matching the authors' three obstacles. Because rewards beyond the observed horizon are missing, can a usable calibration label be formed from a few observed rewards plus a tail sampled from a learned distribution? Can serially dependent transitions be made close enough to exchangeable for split conformal calibration? Can shifts in the state distribution over time, and in actions between behavior and target policies, be corrected by resampling the calibration set instead of reweighting the conformal quantile? The paper also hypothesizes that the coverage loss is bounded by the error of the estimated weights plus a model-error term that shrinks as the rollout length k grows.

\question{4. What is the significance of the question?}

Policy evaluation usually reports an expected return, but a mean conceals the spread of outcomes. In treatment planning, robotics, or recommendation, a policy can look good on average while producing poor results for some states, and a prediction interval for the return expresses that risk before deployment. Distributional RL already estimates whole return distributions, yet its quantiles inherit any bias from function approximation, and convergence results for quantile TD in finite spaces do not guarantee that the limit equals the true return distribution, so quantile intervals can undercover without warning. Conformal prediction offers model-agnostic validity, but its exchangeability premise fails for Markov data, and in the infinite-horizon case the returns it needs as labels are never fully observed. Earlier conformal RL methods were limited to finite horizons, or to bandits with discrete actions, or relied on trajectory-level importance ratios whose variance grows multiplicatively with the horizon. A method that works from ordinary transition tuples, handles policy shift, and states how model error erodes coverage therefore filled a recognized gap between statistical inference and practical RL.

\question{5. What are the contributions reported in this paper?}

The authors list five contributions. The first is a pseudo-return construction: for each calibration transition, the k observed discounted rewards are added to a discounted tail value sampled from the learned return distribution at the state reached k steps later. This supplies a stand-in label for the unobservable infinite-horizon return, and any distributional RL estimator, such as QTD or C51, can supply the tail. The second is a calibration step that treats the calibration data as a replay buffer and draws random subsamples of transition tuples to weaken temporal dependence. The third is time-aware weighted subsampling: calibration tuples are resampled in proportion to estimated importance weights, namely a density ratio between test (initial) states and calibration states in the on-policy case, and that ratio multiplied by k per-step target-to-behavior action-probability ratios in the off-policy case. The authors resample rather than compute a weighted conformal quantile because that approach needs the test point's full covariates, whereas here only the starting state of a test trajectory is seen and the matching weight would require integrating over unobserved futures. The B resulting intervals are merged by a multi-split voting rule governed by a parameter \ensuremath{\xi}. The fourth contribution is theoretical: asymptotic lower bounds on coverage for both settings, expressed through the 1-Wasserstein error of the return model and the $L_1$ error of the estimated weights. The fifth is an empirical study on a two-state MDP, a two-dimensional continuous-state MDP, Mountain Car, and a 50-dimensional extension, with code released on GitHub.

\question{6. What is the research method?}

The method has three parts: algorithm design, asymptotic analysis, and Monte Carlo simulation. Following split conformal prediction, transition tuples are divided into training and calibration sets. On the training set a distributional model is fit with QTD, which represents the return distribution at a state by m equally weighted atoms aimed at the (2i\ensuremath{-}1)/(2m) quantiles; off-policy, QTD is fit on state-action pairs using target-policy next actions and then averaged over \ensuremath{\pi}. Its mean gives the value estimate $\hat{v}(s)$. For each calibration tuple, the pseudo-return is $\tilde{G}(k)(S_t) = \sum_{h<k} \gamma^h R_{t+h} + \gamma^k \tilde{G}$, with $\tilde{G}$ drawn from the learned distribution at $S_{t+k}$, and the nonconformity score is $|\tilde{G}(k) - \hat{v}(S_t)|$. Importance weights come from a classifier separating initial states (the test distribution) from states pooled over all time steps (logistic regression, or frequency counts when discrete), multiplied off-policy by the product of $\pi/\hat{\pi}_b$ over the k steps, with $\hat{\pi}_b$ estimated from data. In each of B rounds, l tuples are resampled by weight, the $\lceil l(1 - \alpha\xi) \rceil$-th smallest score sets a half-width, and the round's interval is $\hat{v}(S_{\mathrm{test}})$ \ensuremath{\pm} that half-width; the final region keeps values covered by at least a $1 - \xi$ fraction of the B intervals.

The coverage analysis chains the \ensuremath{\gamma}-contraction of the distributional Bellman operator in Wasserstein distance with a Kolmogorov-to-Wasserstein inequality and a total-variation bound on weight error.

Experiments use \ensuremath{\gamma} = 0.8 with N = 400 (two-state MDP) or N = 200 (continuous state) trajectories of length T = 30, QTD with 20 quantiles (tabular) or 30 from a neural network, B = 50, 310 test points per run, k = 1 to 5, a 90\% target, with uncorrected DRL quantile intervals (DRL-QR) as baseline. Policy shift comes from changing the two-state transition probabilities from 0.4 and 0.8 to 0.5 and 0.7, or the continuous example's sigmoid action-mixture weights from 0.5/0.5 to 0.6/0.4. Mountain Car uses \ensuremath{\gamma} = 0.99, Q-learning-derived mixture policies (weight 0.3 for behavior, 0.2 for target), and a kernel density estimate from Monte Carlo rollouts as the wrapped return model, with its raw quantile interval (KDE-QR) as baseline.

\question{7. What are the results?}

Theorems 1 and 2 show that the limiting coverage, as the calibration set grows, is at least $1 - \alpha - \Lambda$, where \ensuremath{\Lambda} adds half the time-averaged $L_1$ error of the estimated weights (state weights on-policy, k-step segment weights off-policy) to $\sqrt{2L\gamma^k \mathbb{E}[\bar{W}_1]}$, with $\bar{W}_1$ the worst-case 1-Wasserstein error of the learned return distribution and L a bound on the true return density.

In the two main simulations, conformal coverage stayed close to the 90\% target for k = 1 to 5, on- and off-policy, with medians of roughly 0.89 to 0.92, while DRL-QR undercovered at roughly 0.84 to 0.87. Interval length grew with k and exceeded DRL-QR in the continuous example (about 9 to 10 versus 8). An appendix study with \ensuremath{\xi} = 0.8 found undercoverage at k = 1 (0.87 in both settings), near-nominal coverage at k = 2 (0.90 on-policy, 0.91 off-policy), and overcoverage rising to 0.94 at k = 8, as mean on-policy length rose from 7.78 to 9.50; the authors recommend k = 2 or 3. Small \ensuremath{\xi} caused overcoverage (0.95 to 0.97 at \ensuremath{\xi} = 0.1), while \ensuremath{\xi} of at least 0.5 with k = 2 or 3 worked well. Against Foffano et al.'s weighted method on a horizon-20 two-state problem, coverage was about 0.91 to 0.93 versus 0.85, though Foffano's intervals were shorter (about 6.5 versus 7.7 to 8.0). In Mountain Car, on-policy medians sat near 0.90 against 0.87 for KDE-QR, but off-policy medians stayed near 0.85 for every k (KDE-QR 0.83), with wider intervals (about 15 versus 10). In the 50-dimensional extension, k = 1 undercovered, k = 2 was near target, larger k overcovered, and DRL-QR covered about 0.81.

\question{8. What are the impacts and limitations of the contributions?}

The paper's central lesson is that a conformal calibration layer can sit on top of an ordinary distributional RL estimator using standard transition data, without full trajectories or horizon-length importance ratios, and the coverage gap splits into a weight-error term and a model-error term, $\sqrt{2L\gamma^k \mathbb{E}[\bar{W}_1]}$, that shrinks geometrically in k. Rollout length thus becomes an interpretable bias-variance dial, and leftover miscoverage traces to two identifiable error sources.

Limitations are substantial. The authors admit that choosing k is unresolved, and their ablation swings from under- to overcoverage across k. The guarantees are asymptotic and marginal over test states, and the absolute-residual score of the main experiments gives every state the same width. The aggregation parameter \ensuremath{\xi} was tuned on long simulated rollouts unavailable offline. Experiments are small and synthetic (the largest is a 50-feature binary extension where the action touches one feature), and Mountain Car, the only standard task, stayed near 0.85 off-policy and has discrete returns that fall outside the proofs' bounded-density condition.

For BCA, Theorem 1 is a caution: with k = 1, calibrating bootstrapped targets certifies coverage only up to a model-error term of order $\sqrt{2L\gamma \mathbb{E}[\bar{W}_1]}$. By analogy, a width calibrated on detached one-step Bellman residuals is valid relative to the critic's own bootstrapped target, whose errors carry into the guarantee. The paper's state and action ratios also show how non-uniform WBCP weights could correct dataset-to-policy shift; BCA's uniform weights avoid density-ratio estimation but calibrate to the dataset distribution. The paper's constant-width intervals contrast with BCA's fitted scale model, which lets the width vary per sample.

\question{9. How were the results evaluated or validated?}

Theorems 1 and 2 prove asymptotic coverage lower bounds in the appendix under Markov, time-homogeneity and stationarity assumptions, a bounded return density, finite weight moments and, off-policy, uniform overlap. The off-policy proof covers only trajectory-wise splitting, and the B-interval aggregation is justified by citing earlier multi-split results, not proved directly. Empirically, a Monte Carlo study uses environments where long rollouts approximate true returns, so coverage is measured directly. Metrics are the share of test returns inside the interval and the mean interval width at a nominal 90\% level, shown as boxplots over 50 to 100 repetitions, plus mean and standard-error tables for the k and \ensuremath{\xi} studies. Baselines are the uncorrected DRL quantile interval, a Monte Carlo KDE interval, and Foffano et al.'s method.

This design isolates coverage cleanly but is narrow. Problems are small and synthetic, \ensuremath{\xi} was chosen with ground-truth returns, the main text and appendix disagree on repetitions, B and quantile counts, and no compute figures appear although the checklist claims them. Some claims outrun the figures: the appendix claims shorter intervals than Foffano et al., yet Figure 3 shows Foffano's (undercovering) intervals are shorter, and the main text calls Mountain Car coverage near-nominal while off-policy medians sit near 0.85. The Foffano comparison also ran a modified method, with Foffano's double-quantile score and Meinshausen et al.'s aggregation.

\question{10. What are the ethical implications? How are they addressed?}

The paper has no ethics, broader-impact, or societal-impact section. Its only treatment is the NeurIPS checklist, where the authors answer No to the broader-impacts and safeguards questions, claim compliance with the NeurIPS Code of Ethics on the grounds that the data are public and no people took part in the study, and point to the conclusion for limitations, which mainly concerns the choice of k. Yet the introduction motivates the method with clinical treatment policies, robotics, and autonomous systems.

That framing raises real concerns. An interval labeled 90\% can create unwarranted confidence when its assumptions (Markov and stationary dynamics, overlap between behavior and target policies, accurate weights) fail, as they often do in observational health records with unmeasured confounding. The guarantee is marginal over test states, so a patient subgroup, the very risk the authors cite, can be undercovered while average coverage looks fine. Coverage is only asymptotic, the evidence is entirely simulated, and one off-policy benchmark fell short of nominal coverage. Some points favor the paper: code is public, the experiments are small simulations (though no compute figures are reported), dual-use risk is low, and exposing uncertainty is itself a safety-promoting goal. Rather than engaging these issues explicitly, the paper relies on the broad argument that reporting intervals makes RL systems easier to trust and inspect.

\question{11. What themes are suggested by this research? (minimum 3)}

\textbf{Distribution-free (conformal) guarantees.} The method places a split conformal calibration layer on top of any return-distribution estimator, so coverage does not hinge on QTD being correctly specified. The guarantee is weaker than classical conformal validity, however: exchangeability is only approximately restored through replay subsampling and weighting, and the bound is asymptotic, with a penalty for weight error and a model-error term that shrinks geometrically as k grows.

\textbf{Uncertainty quantification and calibration.} The paper argues that expected-return estimates are insufficient and that returns should be reported with prediction intervals. Its central empirical message is that raw distributional RL quantiles are miscalibrated, undercovering at a 90\% target, and that a calibration layer on held-out transitions restores near-nominal coverage at the cost of wider intervals, which is the familiar trade between validity and sharpness.

\textbf{Distribution shift and out-of-distribution value overestimation.} Here the shift appears in two forms: the state distribution drifts across time steps, and the target policy chooses actions differently from the behavior policy. The paper corrects both by resampling calibration tuples with estimated density and action-probability ratios, which requires overlap and a good behavior-policy model; weight-estimation error enters the coverage bound directly. The paper concerns the shift half of this theme; overestimation does not arise because no policy is optimized.

\textbf{Online tuning and offline model selection.} Two settings control validity: the rollout length k and the aggregation parameter \ensuremath{\xi}. Coverage moved from 0.87 to 0.94 across k, and \ensuremath{\xi} was picked by cross-fitting against long simulated rollouts whose true returns are known. A practitioner with only logged data cannot do this, so the open k-selection problem the authors acknowledge is an instance of offline model selection.

\textbf{Theory versus practice.} The theorems require a bounded return density, overlap, and a calibration set growing without limit, and they bound marginal rather than conditional coverage. The authors note that Mountain Car's discrete returns fall outside the density condition and argue it is not needed there, yet the proof uses it, and that experiment's off-policy coverage stayed near 0.85. The text also promises finite-sample guarantees that the asymptotic theorems do not deliver.

\clearpage

\bookmarked{\subsection}{subsection}{Paper 11: Bellman Calibration for V-Learning in Offline Reinforcement Learning (Bellman Calibration)}

\question{1. What is the citation?}

\begin{hangref}van der Laan, L., \& Kallus, N. (2026). \emph{Bellman calibration for V-learning in offline reinforcement learning} (arXiv:2512.23694v2) [Preprint]. arXiv. \url{https://doi.org/10.48550/arXiv.2512.23694}\end{hangref}

\question{2. What is the research area?}

The paper lies between machine learning and mathematical statistics, within reinforcement learning, specifically the off-policy evaluation and value-prediction side of offline (batch) RL. It is model-free and value-based, acting on fitted Q evaluation or V-learning outputs for a fixed target policy, not on policy optimization. The problem family is post-hoc recalibration of a trained value function from logged transitions. It extends supervised post-hoc calibration (histogram binning, isotonic regression) and causal calibration of treatment effects to the Bellman setting; its nearest predecessor used isotonic Q-calibration only as an auxiliary correction inside double reinforcement learning estimators of policy value. It complements minimax, conservative and structurally constrained value learners, since it uses only their scalar outputs. Its link to offline RL is the critic, whose reliability underpins both policy evaluation and policy learning.

\question{3. What is the research question?}

Can a value predictor that has already been trained off-policy be checked for systematic errors in its numerical scale, and repaired by learning a single scalar transformation of its outputs from logged transitions, with finite-sample error bounds that need neither Bellman completeness nor a correctly specified value class?

The question is approached in steps. The authors first ask what calibration should mean when the long-run return is never observed, and answer with a criterion built from one-step Bellman backups grouped by predicted value. They then ask how this calibration error relates to ordinary value error, whether it can be estimated from behavior-policy data robustly to errors in either the importance ratio or the outcome model, whether iterated histogram and isotonic calibrators attain rates typical of one-dimensional smoothing without inflating value error, and whether the corrections help in simulations built around specific failure modes and on a neural Hopper critic.

\question{4. What is the significance of the question?}

Value estimates feed real decisions, so their scale has consequences. The authors point to clinicians projecting long-term outcomes of treatment rules, platforms forecasting customer lifetime value, and economists studying counterfactual programs. In each case, a value scale that is systematically too high or too low over some range of predictions can distort comparisons between policies even when the ordering of states is sensible.

Bootstrapped value learning compounds early-stopping, approximation and distribution-shift errors, and the usual finite-sample guarantees for fitted value or Q evaluation rest on strong structural conditions such as Bellman completeness or realizability, which are hard to verify, and seldom hold exactly, for flexible neural or tree-based critics. Minimax alternatives loosen those conditions but bring saddle-point training and demanding coverage requirements. At the time of writing, calibration had a mature theory for classification, regression and treatment-effect prediction, yet there was little model-agnostic theory for diagnosing and repairing the scale of a learned value function after training. A cheap one-dimensional correction with provable rates supplies that missing theory and gives offline RL practitioners a way to check a critic without ground-truth returns.

\question{5. What are the contributions reported in this paper?}

Van der Laan and Kallus claim four contributions. First, they formalize Bellman calibration, which asks that states grouped under a common predicted value have an average one-step target-policy backup (reward plus discounted next-state prediction) equal to that shared prediction. They show that the true value function satisfies this property while the converse fails, interpret it through the reward implied by the predictor, and note that it collapses to ordinary regression calibration when \ensuremath{\gamma} = 0.

Second, they define an $\ell_2$ Bellman calibration error, a single-number summary of systematic miscalibration, and establish a calibration--refinement bound in which value error is at most a concentrability-and-discount factor times the sum of the calibration error and a refinement term, the smallest error any transformation of $\hat{v}$ could achieve. To estimate the error from behavior-policy data, they build a doubly robust Bellman pseudo-outcome whose conditional bias is a product of importance-ratio and outcome-model errors, together with a cross-fitted debiased estimator.

Third, they introduce Iterated Bellman Calibration, a post-hoc procedure in the style of fitted value iteration that works with any model. Its iterates always remain transformations of the initial prediction, and it comes in histogram, isotonic and hybrid isotonic--histogram forms, with the hybrid as the recommended default; the same machinery covers an action-value variant that calibrates $\hat{Q}(S, A)$ on logged pairs. Fourth, they supply finite-sample bounds on calibration error and value error for the histogram and hybrid calibrators and a calibration-error bound for the isotonic one, none assuming Bellman completeness or realizability. Beyond this list, the paper adds evidence from simulations and a Hopper benchmark audit.

\question{6. What is the research method?}

The work is mainly methodological and theoretical, supported by simulations and a benchmark audit. The setting is a discounted MDP with discrete actions; a value predictor $\hat{v}$ is trained on separate data, and a calibration sample of i.i.d. transitions $(S, A, R, S')$ is drawn under a behavior policy $b_{0}$. The population target is $\Gamma_0(v)(s) = \mathbb{E}_\pi[R + \gamma v(S') \mid v(S) = v(s)]$, and the error is $\operatorname{Cal}(v) = \|v - \Gamma_0(v)\|_{2,\rho}$. Because data are off-policy, backups are estimated with the pseudo-outcome $\hat{Y}_v = (\pi \hat{q}_v)(S) + \hat{w}_\pi(A \mid S)\{R + \gamma v(S') - \hat{q}_v(S, A)\}$, where $\hat{w}_\pi$ estimates $\pi/b_0$ and $\hat{q}_v$ the one-step outcome regression; $\hat{q}_v = 0$ recovers importance weighting.

Iterated Bellman Calibration (Algorithm 1) starts at $V_0 = \hat{v}$ and repeatedly regresses the pseudo-outcomes of the current iterate on the fixed inputs $\hat{v}(S_i)$ within a calibrator class, so $V_{k+1} = g_{k+1} \circ \hat{v}$. Histogram calibration averages targets within B equal-mass or equal-width bins of $\hat{v}$; isotonic calibration fits a monotone map by pool-adjacent-violators; the hybrid (Algorithm 2) freezes the flat regions of one isotonic fit as bins and then iterates the histogram update. The experiments also run a weighted affine (linear) calibrator as a simple baseline.

The simulated MDPs have continuous states, three actions, and softmax target and behavior policies with shifted logits. Four regimes appear: an affine reward transform with restricted linear FQE; persistent dynamics with \ensuremath{\gamma} = 0.985 and random-feature FQE stopped after K = 4 iterations; a reward shift with an old-regime batch of 2,000 and a current-regime calibration batch of 100; and a saturating monotone reward with linear FQE. Calibration uses importance-weighted fitted value iteration with ratios truncated at 20 and rescaled to average one, 5-fold cross-calibration with median aggregation (except under the reward shift), and 100 replications. On the Deep OPE benchmark's hopper-medium-v0 data, a frozen neural FQE critic (two 256-unit layers, 1,000 updates, or 100 in a stress test) is calibrated on logged pairs through $\hat{Q}(S, A)$, whose backup $R + \gamma \mathbb{E}_{A' \sim \pi}\, g(\hat{Q}(S', A'))$ takes the target-policy mean over next actions and needs no weights, across eleven target policies and 20 random 60/20/20 trajectory splits.

\question{7. What are the results?}

In theory, histogram calibration error is bounded by a one-dimensional statistical term of order $\sqrt{(B/n) \log(n/B)}$, a product of importance-ratio and outcome-model errors, and a finite-iteration term that shrinks as $\gamma^K$ under the conditions of Theorem 4; since \ensuremath{\gamma} enters only through these last two terms, the leading rate does not grow with the horizon. The isotonic version attains $n^{-1/3}$ without choosing B. With equal-mass bins, $B \asymp n^{1/3}$, and negligible nuisance and iteration error, calibrated value error is at most $\sqrt{C_{B,\gamma}}/(1 - \gamma)$ times the original error plus an $O_p((\log n/n)^{1/3})$ term, and the hybrid inherits both bounds.

In Figure 1's affine panel, all three calibrators cut value MSE to 0.42 of raw and calibration error to 0.19--0.23. With K = 4 early stopping, MSE fell by about a fifth (0.79--0.80). Under monotone misspecification, the linear map barely helped (MSE 0.97, calibration 0.99), whereas histogram and hybrid maps nearly halved calibration error (0.54--0.55) with small MSE gains (0.92--0.94); a second monotone panel gave the hybrid 0.26 (MSE) and 0.04 (calibration). After a reward shift, recalibration on the current-regime batch gave MSE near 0.63 with win rates of 0.97--0.98, although calibration error fell only to 0.72 (linear) and 0.92 (isotonic).

On Hopper with the better-trained critic, the hybrid cut plug-in and debiased calibration error to 0.787 and 0.771 of raw while Q-Bellman MSE (1.002) and OPE error (0.978) stayed flat; histogram gave 0.874/0.860 with MSE 1.064, whereas linear and isotonic maps worsened both diagnostics (about 1.5 and 1.4 times raw for linear and 1.2 for isotonic in Figure 2). With the 100-update critic, every calibrator helped (plug-in calibration ratios 0.39--0.47, MSE 0.32--0.41), yet OPE error improved only to 0.97.

\question{8. What are the impacts and limitations of the contributions?}

Bellman calibration offers a cheap, model-agnostic audit and repair step for any critic: it needs only scalar predictions and held-out transitions, leaves the value model untouched, and avoids the completeness and realizability conditions of fitted value theory. The calibration--refinement bound also shows when recalibration can help (good ranking, wrong scale) and when it cannot (state information lost in the scalar prediction).

The authors name extrapolation, nuisance estimation and coverage as the main practical challenges, concede that fully iterated isotonic calibration lacks a value-error bound and that histogram maps need bin tuning, and stress that calibration is not accuracy. Further, the analysis assumes i.i.d. transitions and centers on doubly robust targets, yet no experiment uses an outcome model: the simulations rely on the importance-weighting special case with weights clipped at 20 and renormalized, so clipping bias enters the bound directly rather than as a product of two small errors, and Hopper calibrates without weights on transitions from dependent trajectories. On the single standard benchmark, scalar OPE error barely moves (0.978 and 0.97), and the better-trained critic improves only in calibration error.

Relative to BCA, Bellman calibration is a complement, not a competitor. BCA reads uncertainty from the spread of detached Bellman residuals, while Bellman calibration targets their conditional mean given the predicted value; systematic within-level bias in a host critic could inflate BCA's scale model, making the $\ell_2$ calibration error a sensible companion diagnostic. Its guarantees concern a frozen critic for a fixed target policy, whereas BCA's hosts update critic and actor together, and they give no coverage statement.

\question{9. How were the results evaluated or validated?}

Validation combines proofs, controlled simulation and a benchmark audit. The theorems are proved in Appendices F through L under boundedness, sample splitting, Lipschitz-type stability of the outcome regression, bounded variation for the isotonic case, and (for value error) finite projected concentrability. The simulations compare against Monte Carlo ground-truth values and report held-out value MSE and a binned calibration error (50,000 independent transitions, 50 quantile bins) as ratios to the raw learner refit on all training data, over 100 replications with seed-matched win rates. The Hopper audit fixes hyperparameters across seeds and target policies, shares splits across calibrators, scores plug-in and debiased calibration error and Q-Bellman MSE on held-out trajectories, keeps official returns for a secondary OPE check only, and reports bootstrap intervals that resample matched seed--policy units over 220 fits.

The protocol guards against leakage and reports cases where calibrators hurt, but the simulated regimes are built around the distortions the method fixes and lack intervals. The state dimension, most sample sizes, bin counts and how settings were chosen go unreported. Affine MSE ratios differ without explanation between Figure 1 (0.42) and Table 2 (0.73--0.75), one Figure 1 panel lacks a described design, and the reward-shift retrain comparator and Hopper win rates are missing. Beyond the raw learner and a linear baseline, no external OPE or recalibration method is compared.

\question{10. What are the ethical implications? How are they addressed?}

Bellman calibration is presented as pure methodology, with no ethics statement, broader impact section, or competing-interest declaration. Related content is limited to the introduction's motivating examples (treatment rules in medicine, retention and customer lifetime value, counterfactual economic programs) and a paragraph in the appendix conclusion (Appendix B) naming extrapolation, nuisance estimation and coverage as the main practical challenges; there is no separate limitations section.

The first concern is fairness. Because calibration conditions only on the predicted value, errors for one subpopulation can cancel against opposite errors elsewhere at the same prediction level, so a low calibration error could hide systematic harm to a group of patients or users, echoing fairness critiques of marginal calibration. A small calibration error might also be over-read as validation when the refinement or coverage terms are large. Off-policy targets depend on overlap, and the bias from clipping weights at 20 is never quantified. Recalibrating from a current-regime batch of only 100 could inject noise into high-stakes choices. The evidence is limited to synthetic MDPs and a locomotion simulator, no code release is mentioned, and no compute figures are given, although a one-dimensional recalibration is cheap by construction. The retention and lifetime-value use cases also carry dual-use risk in engagement optimization.

The paper is open about the technical risks, repeatedly stating that calibration is not accuracy, but it does not discuss subgroup effects, human oversight or deployment safeguards.

\question{11. What themes are suggested by this research? (minimum 3)}

\textbf{Uncertainty quantification and calibration.} The paper transfers to long-horizon value functions the classical requirement that, among cases receiving the same prediction, outcomes should match that prediction on average. Since returns are never observed directly, the reference becomes a one-step Bellman backup, and the paper supplies a scalar calibration-error diagnostic plus histogram, isotonic and hybrid recalibration maps. Calibration here concerns the reliability of the mean prediction, not interval coverage.

\textbf{Theoretical guarantees under partial data coverage.} The calibration-error bounds sidestep Bellman completeness and realizability and carry no concentrability factor. The histogram value-error bound relies on completeness holding automatically once the predictor is binned, yet it, like the calibration--refinement bound, still carries a projected concentrability factor that grows when projected target-policy rollouts drift from the data distribution. The authors list coverage among the main practical challenges and calibrate $\hat{Q}(S, A)$ on logged pairs in Hopper to avoid noisy density ratios.

\textbf{Distribution shift and out-of-distribution value overestimation.} Off-policy fitting, a change in the reward regime and weak overlap are presented as reasons a critic can rank states sensibly yet misjudge their scale. The paper treats this error as two-sided, since early stopping under-scales values, rather than as the out-of-distribution overestimation that conservative methods target. A small recent batch can repair a shifted scale, but off-policy recalibration relies on weights that grow unstable where behavior and target policies diverge most.

\textbf{Theory versus practice.} The guarantees cover calibrators that keep one fixed partition of the predictions, and the authors recommend the isotonic--histogram hybrid because fully iterated isotonic calibration falls outside the value-error analysis. On Hopper, the linear and isotonic maps, which lack value-error guarantees, made the well-trained critic worse, and the authors trace this to their missing closure under the Bellman update, so the boundary between analyzed and merely used procedures shows up in the data.

\textbf{Minimalism and simplicity in algorithm design.} Iterated Bellman Calibration keeps the critic, its features and the target policy fixed and learns only a scalar-to-scalar transformation of the critic's own output, computed by within-bin averaging or a pool-adjacent-violators fit at each step. This deliberately small, post-hoc design is what makes assumption-light guarantees possible, and it shows how much reliability can be recovered without retraining a deep value network.

\clearpage

\bookmarked{\subsection}{subsection}{Paper 12: Weighted Bayesian Conformal Prediction (WBCP)}

\question{1. What is the citation?}

\begin{hangref}Lou, X., \& Luo, P. (2026). \emph{Weighted Bayesian conformal prediction} (arXiv:2604.06464v2) [Preprint]. arXiv. \url{https://doi.org/10.48550/arXiv.2604.06464}\end{hangref}

\question{2. What is the research area?}

The paper is about uncertainty quantification through conformal prediction, a topic in statistical machine learning; it is not a reinforcement learning paper. Within conformal prediction it sits where two lines of work meet: weighted conformal prediction, which reweights calibration scores to cope with covariate shift (Tibshirani et al., 2019) or to localize intervals (Guan, 2023), and Bayesian readings of conformal prediction, in particular Bayesian quadrature conformal prediction (BQ-CP; Snell \& Griffiths, 2025), which places a posterior over the threshold. Its application domain is spatial statistics and geographic machine learning, where it extends the first author's earlier frequentist GeoCP method. The connection to offline RL is indirect but concrete: the paper supplies the calibration procedure that BCA runs, with uniform weights, to convert Bellman-residual scores into a calibrated width.

\question{3. What is the research question?}

The paper asks: can the Bayesian posterior over the conformal threshold that BQ-CP offers for i.i.d. data be carried over to weighted conformal prediction, so that one procedure stays valid under covariate shift or localized weighting and also reports how reliable its own threshold is?

From this come several narrower questions. The authors ask which concentration parameter a weighted Dirichlet model over calibration-score spacings should use, and they propose that Kish's effective sample size is the right one because it equates Bayesian and frequentist variances. They ask whether the posterior width shrinks at a rate set by that effective sample size, and whether BQ-CP's data-conditional guarantee and a conditional coverage gain survive weighting. Empirically, they test whether a spatial version with adaptive kernel bandwidth holds nominal coverage on many geospatial datasets while producing location-specific reliability diagnostics, and whether both the adaptive bandwidth and the Bayesian layer are necessary.

\question{4. What is the significance of the question?}

Conformal prediction has become a standard way to attach intervals with finite-sample coverage to any predictor, yet the threshold it returns is a single number that says nothing about how firmly the calibration data pin it down. The authors make the point with two equal thresholds, one backed by two hundred relevant calibration points and the other by five distant ones; neither standard nor weighted conformal prediction can distinguish them. This matters in practice wherever interval width drives decisions, as in property valuation, environmental monitoring or public-health planning, and wherever calibration data are spread unevenly.

The question was still open because the two available remedies did not combine. BQ-CP, which had just received an outstanding paper award at ICML 2025, delivered a threshold posterior and data-conditional guarantees but relied on exchangeability, while weighted conformal prediction handled shift but stayed a point-estimate method. Bridging them tests whether meta-uncertainty can be measured in exactly the non-i.i.d. settings where it is most needed. For BCA the question is also practical, since a calibrated width that carries an estimate of its own reliability is a natural input for deciding how strongly to regularize an offline RL learner.

\question{5. What are the contributions reported in this paper?}

WBCP itself is the first contribution. It replaces BQ-CP's flat Dirichlet over the gaps between sorted calibration scores with a weighted Dirichlet whose parameters are the normalized importance weights scaled by Kish's effective sample size, $n_{\mathrm{eff}} = (\sum w_i)^2/\sum w_i^2$. The construction draws on Newton and Raftery's weighted Bayesian bootstrap, and Algorithm 1 implements it by Monte Carlo, returning a conservative highest-posterior-density (HPD) threshold, the posterior standard deviation $\sigma_{\mathrm{post}}$, $n_{\mathrm{eff}}$ and the threshold samples. The paper also formalizes the difference between predictive uncertainty and meta-uncertainty.

The second contribution is theoretical. Four results state that, asymptotically, only a concentration equal to $n_{\mathrm{eff}}$ makes the posterior variance agree with the frequentist variance of the weighted empirical distribution; that $\sigma_{\mathrm{post}}$ shrinks as $O(1/\sqrt{n_{\mathrm{eff}}})$; that BQ-CP's stochastic-dominance bound carries over under the assumed weighted Dirichlet model, with the HPD threshold's data-conditional guarantee stated as a corollary for the high-$n_{\mathrm{eff}}$ regime only; and that the HPD threshold lowers expected conditional miscoverage when $n_{\mathrm{eff}}$ is large. Supporting propositions show that uniform weights give $n_{\mathrm{eff}}$ = n and a flat Dirichlet (a calibration-only version of BQ-CP, up to an O(1/n) correction), characterize very wide and very narrow kernels, and prove that an adaptive bandwidth keeps $n_{\mathrm{eff}}$ bounded away from zero.

The third contribution is AdaGeoBCP, a spatial instantiation with Gaussian kernel weights and a per-location bandwidth tied to k-nearest-neighbor distances, which yields maps of $n_{\mathrm{eff}}$ and $\sigma_{\mathrm{post}}$. Finally, the authors provide an empirical study on a synthetic Gaussian random field and 13 geospatial datasets, with a 2\ensuremath{\times}3 ablation, a \ensuremath{\beta} sweep and supplementary bootstrap data.

\question{6. What is the research method?}

WBCP is produced by algorithm design and asymptotic proof sketches and supported by empirical benchmarking with ablations. It starts from split conformal prediction with absolute-residual scores $\rho_i = |Y_i - \hat{f}(X_i)|$ and a weighted $(1 - \alpha)$ quantile of those scores. With the scores sorted and $\tilde{w}(i)$ the normalized weights in that order, each of M draws samples u from $\operatorname{Dir}(n_{\mathrm{eff}} \cdot \tilde{w}(1), \ldots, n_{\mathrm{eff}} \cdot \tilde{w}(n))$ and takes as threshold the smallest score at which the running sum of u reaches $q = 1 - \alpha$. The \ensuremath{\beta}-quantile of the M thresholds is $\lambda_{\mathrm{HPD}}$, the interval is $\hat{f}(x) \pm \lambda_{\mathrm{HPD}}$, and the spread of the draws gives $\sigma_{\mathrm{post}}$. With uniform weights $n_{\mathrm{eff}}$ = n and the procedure matches BQ-CP up to an O(1/n) term (only calibration spacings are modeled), at a cost of O(n log n + nM) per weight profile. The theory assumes the weighted Dirichlet spacing model and, for coverage, correct likelihood-ratio weights plus a Gaussian (Bernstein--von Mises) approximation, under which $\lambda_{\mathrm{HPD}} \approx \hat{q}_w + \Phi^{-1}(\beta)\sigma_{\mathrm{post}}$.

AdaGeoBCP sets $w_i(s) = \exp(-\|s - s_i\|^2/(2h(s)^2))$, where $h(s)$ is $h_{0}$ times the median distance from s to its k nearest calibration points. All runs target \ensuremath{\alpha} = 0.1, and the Bayesian variants use M = 1000 draws and, outside the \ensuremath{\beta} sweep, \ensuremath{\beta} = 0.9. The synthetic benchmark places 3,000 points on a 20 \ensuremath{\times} 20 square with a smooth mean surface and Gaussian-random-field noise of length scale 2.0 and region-dependent amplitude, split 50/50 into calibration and test sets, plus a sweep of length scales from 0.5 to 16. The real-world study covers Seattle house prices, CAMELS hydrology, FIA forestry, US county voting, ERA5 climate and eight CDC PLACES county health indicators. The base model is XGBoost (500 trees, depth 3, learning rate 0.1) with an 80/10/10 split repeated over 10 random seeds (so Seattle's calibration set has only 300 points), and Seattle also receives a 50-split bootstrap. Baselines are standard CP, BQ-CP and GeoCP (fixed bandwidth from 20-nearest-neighbor distances), and two ablations remove either the adaptive bandwidth (GeoBCP) or the Bayesian layer (AdaGeoCP). A \ensuremath{\beta} sweep from 0.5 to 0.99 on Seattle traces the coverage--width trade-off.

\question{7. What are the results?}

Across the 13 real-world datasets, AdaGeoBCP reached a mean coverage of 0.925 \ensuremath{\pm} 0.009, essentially tied with BQ-CP (0.928 \ensuremath{\pm} 0.011) and above standard CP (0.903), while the frequentist spatial baseline GeoCP under-covered at 0.824 \ensuremath{\pm} 0.034. The distinguishing result is heterogeneity: the across-location standard deviation of $n_{\mathrm{eff}}$ and $\sigma_{\mathrm{post}}$ was 13.1 and 0.95 for AdaGeoBCP, against 0.0 and 0.02 for BQ-CP, whose posterior is the same everywhere. On the synthetic field, AdaGeoBCP covered 0.983 with a mean width of 1.87, much narrower than standard CP (3.18) and BQ-CP (3.49), yet the frequentist AdaGeoCP already met the target with narrower intervals (0.959, width 1.66). Across the length-scale sweep every spatial variant over-covered (0.96 to 0.99 once $\ell \geq 2$), while standard CP and BQ-CP slipped below target at $\ell = 16$ (0.88 and 0.89).

The ablation separated the two ingredients. Removing the Bayesian layer (AdaGeoCP, mean coverage 0.895) cut the number of datasets meeting the target from 13 to 3, and removing the adaptive bandwidth (GeoBCP, 0.876) cut it to 1, because the fixed bandwidth left only about seven effective calibration points per location on continental-scale data.

The method also has visible costs. In the 50-split Seattle study, AdaGeoBCP matched BQ-CP's coverage (0.929) but gave the widest intervals (24.91 against 20.41 for BQ-CP and 17.68 for standard CP), and its posterior was far more diffuse (mean $n_{\mathrm{eff}}$ 13.5 and $\sigma_{\mathrm{post}}$ 6.80, against 300 and 1.90). The \ensuremath{\beta} sweep showed that its coverage hinges on the HPD level: 0.873 at \ensuremath{\beta} = 0.5, 0.95 at \ensuremath{\beta} = 0.9 and 0.977 at \ensuremath{\beta} = 0.99 (width 30.32), a much steeper frontier than the other Bayesian variants.

\question{8. What are the impacts and limitations of the contributions?}

The main impact is conceptual: WBCP turns a conformal threshold into a distribution for any weighting scheme, so each interval comes with a measure of how well the calibration data support it. The $n_{\mathrm{eff}}$ and $\sigma_{\mathrm{post}}$ maps cheaply show where intervals rest on thin evidence, with no refitting of the base model.

Three caveats come from the authors themselves: autocorrelated residuals inflate $n_{\mathrm{eff}}$, misspecified weights may void the coverage result, and the theorems are sketches under an assumed Dirichlet model. They claim only that WBCP matches weighted CP's marginal coverage as $n_{\mathrm{eff}}$ grows, not exact finite-sample coverage. Theorem 4 needs large $n_{\mathrm{eff}}$, yet AdaGeoBCP ran with $n_{\mathrm{eff}}$ near 10 to 13 on synthetic and Seattle data, and its Seattle coverage depends mainly on the \ensuremath{\beta} = 0.9 inflation. Random real-world splits test no genuine covariate shift, and $\sigma_{\mathrm{post}}$ is never checked against local miscoverage. Seattle diagnostics conflict as well: Figure 3, labeled AdaGeoBCP, shows $n_{\mathrm{eff}}$ of about 100 to 270 and $\sigma_{\mathrm{post}}$ of 1.6 to 2.5, but Table 7 reports AdaGeoBCP means of 13.5 and 6.80; the figure looks like the fixed-bandwidth GeoBCP output (213 and 2.51).

For BCA, uniform weights make WBCP essentially BQ-CP, so the width inherits an exchangeability assumption that policy-induced shift and temporally correlated transitions can break, the latter also understating $\sigma_{\mathrm{post}}$. With 300 Seattle calibration points, BQ-CP's width exceeded standard CP's by about 15 percent; since $\sigma_{\mathrm{post}}$ shrinks as $1/\sqrt{n}$, that margin fades on large offline datasets, and because uniform weights give one global width, BCA's per-sample adaptivity comes from its scale model, not from WBCP.

\question{9. How were the results evaluated or validated?}

Four empirical designs carry the validation: a single-seed synthetic Gaussian random field with known heteroscedastic noise, a sweep over six spatial autocorrelation length scales, 13 real-world regression datasets each resampled over 10 random 80/10/10 splits, and a 50-split bootstrap on Seattle. Metrics are empirical marginal coverage at \ensuremath{\alpha} = 0.1, mean interval half-width, and the mean and across-location standard deviation of $n_{\mathrm{eff}}$ and $\sigma_{\mathrm{post}}$. Results are reported as means \ensuremath{\pm} standard deviations, with no significance tests. A 2\ensuremath{\times}3 ablation isolates the bandwidth and the Bayesian layer, and a \ensuremath{\beta} sweep recomputes the coverage--width frontier from the same posterior draws. The theory rests on proof sketches under the assumed Dirichlet model and, for coverage, correct weights and a Gaussian approximation.

The breadth of datasets and the clean ablation are strengths. A dataset counts as meeting the target when its bootstrap error bar reaches 0.90, a lenient rule, and the single-split per-dataset figure disagrees with the bootstrap: there GeoBCP reaches 0.90 on 12 of 13 datasets, and AdaGeoBCP misses it on five (Hydro near 0.84). XGBoost settings and \ensuremath{\beta} were fixed with no tuning procedure described, and k and $h_{0}$ are never given. Methods are compared at a fixed \ensuremath{\beta}, not matched width; outside Seattle, real-world widths appear only in a log-scale chart; and local coverage and runtime are not measured.

\question{10. What are the ethical implications? How are they addressed?}

The WBCP preprint omits an ethics statement, a broader impact section and any societal discussion. A heading for acknowledgments and funding disclosure appears with nothing beneath it, and the only reflective material is a short technical appendix on limitations (autocorrelated residuals and weight misspecification). All datasets are public and mostly aggregated to counties or grid cells, which keeps privacy risk low.

The ethical stakes follow from the applications used to motivate the method. Intervals and reliability maps for house prices, disease prevalence or vote shares can shape lending, property taxation, health-resource allocation or political targeting. Because sparsely sampled and peripheral areas receive wider intervals and larger $\sigma_{\mathrm{post}}$, the diagnostic could guide new data collection, or it could be used to justify withholding services from places that are already poorly measured. There is also a risk of false assurance: the guarantees rest on an assumed Dirichlet model, asymptotic approximations and correct weights, yet the outputs look like rigorous certificates. Reproducibility also suffers because no code repository is linked.

The authors address these issues only partly. They soften their claims, call the theory sketches and recommend reporting the full \ensuremath{\beta}, coverage and width triple, but they say nothing about misuse, equity or deployment. The same caution applies when the threshold sets conservatism in offline RL, where a miscalibrated width becomes an unsafe or needlessly timid policy.

\question{11. What themes are suggested by this research? (minimum 3)}

\textbf{Uncertainty quantification and calibration.} WBCP adds a second layer to conformal uncertainty: beyond the predictive interval, it quantifies how uncertain the interval's own width is. The posterior standard deviation $\sigma_{\mathrm{post}}$ and the effective sample size $n_{\mathrm{eff}}$ report whether a threshold is firmly supported by calibration data, and the paper formalizes this as meta-uncertainty. The idea matters for any method, BCA included, that treats a calibrated width as a trustworthy input.

\textbf{Distribution-free (conformal) guarantees.} The paper starts from the distribution-free coverage of conformal prediction and tries to keep it while relaxing exchangeability through weighting. Its guarantees become partly model-based, however: the data-conditional result holds under an assumed weighted Dirichlet spacing model and, for the HPD threshold, only in the high-$n_{\mathrm{eff}}$ regime, while coverage claims need correctly specified likelihood-ratio weights. Extending conformal guarantees beyond i.i.d. data thus trades some of their assumption-free character for flexibility.

\textbf{Adaptive, state-dependent conservatism.} Each test location receives its own threshold posterior, and in the Gaussian approximation the HPD rule sets the threshold roughly $\Phi^{-1}(\beta)\sigma_{\mathrm{post}}$ above the weighted quantile, so intervals widen most where calibration support is thin. The adaptive kernel bandwidth adds a second layer of local adaptation. In a regression setting, this mirrors the idea of setting a safety margin per sample from calibrated uncertainty rather than from one global constant.

\textbf{Theory versus practice.} The four theorems are offered as sketches that hold under an assumed model and, for coverage, in a large-$n_{\mathrm{eff}}$ asymptotic regime. On the synthetic and Seattle data the proposed spatial method runs with $n_{\mathrm{eff}}$ near 10 to 13, a small-$n_{\mathrm{eff}}$ regime in which, by the authors' own account, the Gaussian approximation breaks down and the practical value lies in the diagnostic itself. The paper thus shows a gap between asymptotic justification and actual use.

\textbf{Implementation details and hyperparameter sensitivity.} Outcomes depend heavily on two choices. Moving the HPD level \ensuremath{\beta} from 0.5 to 0.99 shifts AdaGeoBCP's Seattle coverage from below target to well above it, and replacing the adaptive bandwidth with a fixed one drops the number of datasets meeting the target from 13 to 1. The authors therefore recommend reporting \ensuremath{\beta}, coverage and width together rather than a single operating point.

\clearpage

\bookmarked{\section}{section}{Part A: Literature Reviews}

This section answers the eleven literature-review questions for each of the three surveys of offline reinforcement learning.

\bookmarked{\subsection}{subsection}{Literature Review 1: Offline Reinforcement Learning: Tutorial, Review, and Perspectives on Open Problems (Levine et al.)}

\question{1. What is the citation?}

\begin{hangref}Levine, S., Kumar, A., Tucker, G., \& Fu, J. (2020). \emph{Offline reinforcement learning: Tutorial, review, and perspectives on open problems} (arXiv:2005.01643v3) [Preprint]. arXiv. \url{https://doi.org/10.48550/arXiv.2005.01643}\end{hangref}

\question{2. What is the survey scope?}

The survey covers offline reinforcement learning, defined as learning a policy from a fixed dataset gathered by a possibly unknown behavior policy, without any interaction with the environment during training. After reviewing RL basics, it examines three algorithmic families: importance-sampling estimators used both to evaluate policies and to compute off-policy gradients, including marginalized estimators such as DualDICE and GenDICE; approximate dynamic programming with policy constraints, uncertainty penalties, or conservative Q-functions; and model-based methods. Section 6 reviews evaluation practice and applications in robotics, healthcare, autonomous driving, advertising and recommender systems, and dialogue, and Section 7 closes with broader perspectives. Its roughly 190 references run from 1991 to mid-2020, and about three quarters date from 2015 or later. Exploration and a complete account of model-based RL are explicitly set aside, and imitation learning appears only as a point of contrast. The level is a tutorial combined with a conceptual review rather than a formal taxonomy, and it addresses researchers starting work on offline algorithms for deep RL.

\question{3. What is the survey question?}

The guiding question of the survey can be put this way: what prevents current reinforcement learning algorithms from turning a fixed, previously collected dataset into the best policy the data can support, and which ideas have been proposed to overcome those obstacles? The survey works through narrower questions, a few posed outright, such as the title of Section 2.4. First, why is learning from static data harder than ordinary supervised learning, and why does distributional shift sit at the center of the difficulty? Second, how does each major family of methods (importance sampling, dynamic programming, and model-based learning) cope with that shift, and where does each one break down? Third, how should offline methods be evaluated, when can they beat every trajectory in the data, and where have they already been applied? Finally, which open problems must be solved before offline RL can support data-driven decision making at scale?

\question{4. What is the significance of the survey?}

The survey first appeared in May 2020, when offline RL had recently become an active topic through methods such as BCQ, BEAR, and BRAC, yet the literature still used competing names (batch RL, fully off-policy RL) and, as the authors themselves note, lacked settled evaluation protocols. Its contribution was an early, unified, tutorial-level account that places the older statistical work on importance sampling and off-policy evaluation next to recent deep dynamic-programming methods, all under one notation and one diagnosis: the distributional shift created by counterfactual queries.

Newcomers gain a single entry point, and experienced researchers gain open problems specific enough to motivate new algorithms. Its influence is visible in the other surveys in this set. Figueiredo Prudencio et al. single it out as the most prominent earlier review, adapt its illustration of online, off-policy, and offline learning, skip marginalized importance sampling because Levine et al. already cover it well, and present their own open-problems section as an update of its list. He (2023) relies on it for the notion of distributional shift. One open problem the survey raises, conservatism that adjusts itself to the data, is also named in Figueiredo Prudencio et al.'s introduction and is the question BCA takes up.

\question{5. What are the themes?}

The authors arrange the material by the mechanism a method uses to reason about actions that are absent from the data: reweighting returns with importance weights, propagating values through dynamic programming, or rolling out a learned dynamics model. Shared notation from Section 2 links these parts, each of the three method sections (3, 4, and 5) closes with a subsection on challenges and open problems, and distributional shift recurs as the thread connecting all of them.

\textbf{Distributional shift as the core obstacle.} Section 2.4 argues that offline RL fundamentally requires answering counterfactual questions, which breaks the i.i.d. assumption behind supervised learning. A behavior cloning example adapted from Ross et al. shows that, even with perfect expert labels, offline error can grow quadratically in the horizon, whereas interactive collection in the style of DAgger keeps it linear. Section 4.2 then traces the training-time problem to Bellman targets that query actions the data never contained, and Figure 2 shows offline SAC on HalfCheetah collapsing during training regardless of dataset size.

\textbf{Importance sampling and off-policy evaluation.} Section 3 reviews estimators that correct for the gap between the behavior and learned policies by reweighting returns: per-decision and weighted importance sampling, doubly robust estimators, high-confidence bounds for safe policy improvement, and off-policy policy gradients. It gives special attention to marginalized importance sampling, splitting it into forward-Bellman methods such as GenDICE and GradientDICE and backward, duality-based methods such as DualDICE and AlgaeDICE. The authors conclude that variance grows with the horizon, so these tools are dependable only when the learned policy stays close to the data.

\textbf{Policy constraints and penalties.} Section 4.3 unifies methods that keep the learned policy near the behavior policy, distinguishing constraints on the actor update from penalties added to rewards or target values. It compares explicit f-divergence constraints such as the KL penalty used in BRAC, implicit constraints enforced through weighted regression in AWR, AWAC, and ABM, and integral probability metrics such as the MMD used by BEAR. With a one-dimensional lineworld example (Figure 3), it argues that restricting the support of the policy is preferable to matching the full behavior distribution.

\textbf{Uncertainty estimation and conservative critics.} Sections 4.4 and 4.5 present the alternative of making the critic itself cautious. Uncertainty-based methods learn a distribution over Q-functions, usually with ensembles, and subtract an uncertainty term from the mean during policy improvement (Equation 24); the ensemble-combination approach of Agarwal et al. is one example. Conservative Q-learning instead adds a penalty that pushes Q-values down under an adversarial action distribution, and a refined variant also pushes up values on dataset actions, yielding a lower bound in expectation without any explicit uncertainty model.

\textbf{Model exploitation in model-based offline RL.} Section 5 explains that a learned dynamics model suffers the same shift, because a policy optimized inside the model seeks out states where the model is wrongly optimistic. The authors restate the MBPO error bound of Janner et al. to separate model error from policy divergence, survey visual foresight, hybrid reward-prediction methods, and deep imitative models, and then describe MOPO's uncertainty-penalized reward and MOReL's absorbing low-reward state. They caution that ensemble disagreement, the proxy for model error used in prior work, carries no consistency guarantee.

\textbf{Evaluation, benchmarks, and applications.} Section 6 criticizes the habit of training on complete replay buffers from online runs, since such data reveal little about improving on suboptimal behavior, and introduces D4RL as a remedy. It identifies the stitching of partial trajectories (Figure 4, Maze2D) as the main route by which offline RL can surpass its data, notes the roles of off-policy evaluation and human judges, and surveys applications such as QT-Opt grasping, sepsis and ventilation policies on MIMIC-III, driving, ad and news recommendation, and dialogue.

\question{6. What is the survey method?}

The survey is an expert narrative written as a tutorial, not a systematic review. The authors report no search databases, keywords, inclusion criteria, or time window, and the works appear to have been chosen for their value in explaining the problem. The material is organized by algorithmic principle rather than by chronology. Section 2 fixes the notation (the MDP tuple, the behavior policy, state visitation distributions, and the Bellman operator) and gives generic pseudo-code for policy gradients, Q-learning, and off-policy actor-critic methods (Algorithms 1 to 3), noting that the latter two become offline methods once the number of new data-collection steps is set to zero. Sections 3 to 5 then rebuild each family of methods inside that frame.

The analysis is formal but does not introduce new theorems. The authors restate and explain results from the literature, such as the behavior cloning and DAgger error bounds, the normal equations behind LSTD-Q and LSPI, the generic constrained and penalized actor-critic objectives, the derivation of the CQL penalty, and the MBPO performance bound. They run no new experiments and provide no benchmark tables; apart from their own schematic in Figure 1, the figures are taken from earlier work, for example the offline SAC curves from Kumar et al. (2019) and the lineworld from Kumar (2019). Sections 3, 4, and 5 each end with a subsection on challenges and open problems, which is where most of the authors' own critical perspective appears.

\question{7. What are the survey findings?}

Importance-sampling variance grows exponentially with the horizon, and even marginalized estimators degrade once the learned policy drifts far from the behavior policy; apart from off-policy evaluation and high-confidence policy improvement, the policy-improvement variants had seldom been used in a fully offline setting. The attainable gain depends on how suboptimal the behavior policy is, the problem's dimensionality, and the effective horizon. Naive dynamic programming fails because target values query out-of-distribution actions and the resulting errors compound without corrective feedback.

Policy constraints help, but most depend on an estimated behavior policy that is hard to fit when the data are multimodal, and because error can scale with the square of the effective horizon, the constraints must be tight, which caps improvement; a lower bound proved by Kidambi et al. shows that this quadratic dependence cannot be avoided in the worst case. Uncertainty-based methods are principled, but calibrated uncertainty for deep Q-functions is hard to obtain, and caution demands more faithful uncertainty than optimistic exploration, so constraint and conservative methods performed better at the time. Conservative critics avoid behavior modeling but can become overly pessimistic, which the authors suspect is a major barrier on current benchmarks, and how to adjust that pessimism on the fly remains open. Model-based methods look promising, since MBPO works offline unmodified, but they depend on ensemble error estimates, struggle with long horizons and images, and have no proven advantage over model-free methods.

On evaluation, the authors judge benchmarking immature: datasets built from full replay buffers are unrealistic, and how broad or multimodal the data are strongly affects task difficulty. Section 7 casts offline RL as counterfactual inference, expects tools from causal inference, uncertainty estimation, generative modeling, distributional robustness, and invariance to help, and argues that large, diverse datasets may matter as much as new algorithms.

\question{8. What are the limitations of the survey?}

The most obvious limitation is currency. The survey's newest citations are from mid-2020, so it predates later standard baselines, including TD3+BC (2021), IQL (2022), and ReBRAC (2023), sequence-model approaches such as the Decision Transformer, and the pessimism theory for partially covered data, such as VI-LCB, that He (2023) reviews. Figueiredo Prudencio et al. note that several new methods using strategies Levine et al. do not discuss have since appeared. Even within its window the coverage is uneven. Uncertainty-based model-free methods get roughly half a page, which the authors call a brief review, the model-based part is likewise labeled a brief survey, and offline-to-online fine-tuning appears only in passing through the QT-Opt example. The applications section is mostly descriptive, offering only short caveats such as biased clinical records and the contextual-bandit simplification in recommendation.

Selection is a second concern. No search protocol is reported, and more than 30 of the roughly 190 references include at least one of the four authors, among them BEAR, BRAC, CQL, AWR, MOPO, and D4RL, which the text presents as the authors' own benchmark. The survey's preference for support constraints and conservative critics mirrors that body of work. It also offers no quantitative comparison, so claims such as constraint and conservative methods beating uncertainty-based ones rest on citation rather than assembled evidence. Finally, offline hyperparameter selection is never treated directly: the CQL weight is only said to need a suitable value, and off-policy evaluation is mentioned merely as a possible way to choose policies offline.

\question{9. How was the survey assessed relative to other surveys?}

Levine et al. include no section on related surveys. They position the article against the older batch RL literature (Ernst et al., Riedmiller, and the chapter by Lange et al.), explain why they prefer the term offline over batch or fully off-policy, and present the work as a tutorial meant to equip readers rather than an exhaustive catalog. Other overview works appear only as pointers, such as the Gottesman et al. guidelines for RL in healthcare and the Yurtsever et al. survey of autonomous driving.

Compared with the other two surveys in this set, Levine et al. give the most formal and pedagogical treatment. The survey by Figueiredo Prudencio et al. (2024), published in IEEE TNNLS, is broader and more systematic: it proposes a taxonomy that adds one-step, imitation-learning, and trajectory-optimization classes, classifies methods in a table, provides a timeline through 2022, reviews D4RL and RL Unplugged against a list of dataset properties, devotes a section to off-policy evaluation, and summarizes relative performance in a heatmap. With 91 references, about half as many as Levine et al., it trades derivation for coverage of post-2020 methods such as TD3+BC, IQL, and the Decision Transformer. He (2023) is a short preprint with 19 references that summarizes seven model-based methods, including COMBO, RAMBO, and ARMOR, all of which postdate Levine et al.; it is more current on model-based work but offers little analysis. Levine et al. remain the deepest account of why offline RL is hard, while the later surveys are better guides to which algorithms exist.

\question{10. What are the ethical implications? How are they addressed?}

The survey has no section on ethics, broader impact, or societal consequences; its relevant content consists of scattered remarks about safety and data. It rules out trying a half-trained treatment policy on real patients (Section 2.3), describes high-confidence policy improvement as a way to guarantee performance in safety-critical settings (Section 3.1), mentions Lyapunov-based safe regions for model-based control (Section 5.2), and warns that clinical records over-represent severe cases, which can lead a naive agent to conclude that treatment causes death (Section 6.3). Safety in driving and financial risk in recommendation appear mainly as reasons to prefer offline learning.

Several implications of the field as presented go unaddressed: privacy and consent for clinical datasets such as MIMIC-III, unequal treatment of groups that are under-represented in logged data, persuasive dialogue agents (one motivating example is a sales chatbot), the replacement of human operators, and the danger of treating off-policy estimates, which the survey describes as still unsettled, as the final check before deployment. As scholarship, the survey is candid about weaknesses of its authors' own methods, such as the underestimation caused by the basic CQL penalty, yet their prominence still shapes which remedies it presents as most promising. Safety thus serves the tutorial as a technical motivation rather than an ethical question, which leaves its handling of ethics limited.

\question{11. What meta-themes are suggested by this survey? (minimum 3)}

\textbf{Distribution shift and out-of-distribution value overestimation.} Because the survey traces every family's failure to the same mismatch between the data and the actions or states a learner queries, later methods can be compared by where they intercept it: in the actor (TD3+BC), in actor and critic together (ReBRAC), in the critic (CQL), or by never querying unseen actions (IQL). BCA reads its uncertainty signal from the Bellman residual, which is exactly where Section 4.2 locates the damage done by out-of-distribution targets.

\textbf{Conservatism versus policy improvement trade-off.} The survey repeatedly returns to the tension between safety and gain. Tight constraints are needed because error compounds roughly with the square of the horizon, yet they hold the policy near a possibly poor behavior policy; KL constraints force needless randomness; conservative critics overcorrect into underestimation. Since Kidambi et al.'s lower bound makes quadratic error growth unavoidable in the worst case, the survey ends Section 4.6 by asking for constraints that strike a good balance on typical problems, a framing that BCA recasts as a calibration problem.

\textbf{Uncertainty quantification and calibration.} The survey treats calibrated uncertainty as both the most principled remedy and the hardest one to deliver. It argues that acting conservatively demands more faithful uncertainty than exploration does, and that the bootstrap-ensemble disagreement prior model-based work uses as an error oracle has no consistency guarantee under sampling error (Section 5.2). The survey never mentions conformal prediction, but its diagnosis leaves an opening that conformal prediction for policy evaluation and WBCP address by replacing heuristic ensemble spread with intervals calibrated to a chosen error rate.

\textbf{Adaptive, state-dependent conservatism.} Section 4.6 leaves open how a method could tune its own caution on the fly, trading the danger of inflated values against the price of refusing anything the data rarely show. Most constraint and conservative methods the survey describes rely on one global coefficient or threshold, such as the divergence budget or the CQL weight, while the state-dependent penalties of Equation 24 and MOPO depend on uncertainty estimates that the authors consider hard to calibrate. BCA's proposed answer is to set constraint strength per sample from calibrated residual uncertainty.

\textbf{Benchmark dependence (D4RL).} The survey criticizes evaluation on full replay buffers and offers D4RL, released by the same authors at about the same time, as the remedy, tying the tutorial closely to the benchmark that later dominated empirical comparison; Figueiredo Prudencio et al. find that most methods are tested on D4RL alone. Because the survey also observes that cheap online evaluation exists only where a simulator does, results tuned and compared on D4RL, BCA's included, say little about settings where no simulator is available to check them.

\clearpage

\bookmarked{\subsection}{subsection}{Literature Review 2: A Survey on Offline Model-Based Reinforcement Learning (He)}

\question{1. What is the citation?}

\begin{hangref}He, H. (2023). \emph{A survey on offline model-based reinforcement learning} (arXiv:2305.03360v1) [Preprint]. arXiv. \url{https://doi.org/10.48550/arXiv.2305.03360}\end{hangref}

\question{2. What is the survey scope?}

He's survey covers offline model-based reinforcement learning, the setting in which an agent fits a dynamics and reward model to a fixed log of experience and then plans or optimizes a policy inside that model, without further interaction with the environment. After brief primers on offline RL and on model-based RL, the body profiles seven methods that appeared between 2020 and early 2023: MOReL, MOPO, BREMEN, VI-LCB (together with the related PEVI and the Bernstein-penalty analysis of Li et al.), COMBO, RAMBO, and ARMOR. The reference list holds nineteen items, and the main text runs to about five pages. Deep continuous-control algorithms and tabular sample-complexity theory are both included, while model-free offline algorithms appear only as a point of contrast. Benchmarks, quantitative results, off-policy evaluation, and application domains are absent; empirical performance enters only through brief restatements of the original authors' claims. The work is a compact, primer-level review for readers who know basic RL and want a quick entry into model-based offline RL.

\question{3. What is the survey question?}

He's review centers on one question: how have recent offline model-based RL methods coped with the distributional shift that arises when a model learned from logged data is used to optimize a new policy, and what problems remain unsolved?

The survey approaches this through smaller questions. It first asks what offline RL and model-based RL each contribute and why combining them is attractive, given that supervised model fitting can make full use of large historical datasets. It then asks why a learned model is especially vulnerable to shift, pointing to model exploitation and to the absence of exploration that could correct model errors. The method sections ask which part of the problem each algorithm alters (the reward, the transition dynamics, or the value function) to enforce caution, and what guarantees follow. The closing section asks what still stands in the way of real deployment.

\question{4. What is the significance of the survey?}

The survey appeared in May 2023, about three years after MOReL and MOPO had made uncertainty-penalized models a main route to offline RL. At that point the general offline RL surveys gave model-based methods only a few pages: one short section in Levine et al. (2020), written before COMBO appeared, and a single subsection of Figueiredo Prudencio et al.'s algorithmic review that stops at COMBO. He's paper responds with a dedicated, recent, and very short treatment that sets adversarial-model methods and minimax sample-complexity theory beside the earlier reward-penalty approaches.

Its main beneficiaries are students and researchers who need a fast orientation to the subfield and a short reading list of its key papers, as well as researchers in model-free offline RL who want to see how the same caution is expressed through a learned model. Its influence on the field is not clearly established. It is a single-author preprint, marked as under review, with no later revision or peer-reviewed version, so its value lies in serving as a concise map rather than an authoritative reference.

\question{5. What are the themes?}

The survey offers no formal taxonomy. After two background sections, it observes at the start of Section 2 that recent methods counter distributional shift by altering the learned model's transition dynamics, reward function, or value function, and it then gives one subsection to each of seven algorithms in roughly chronological order. The themes below group those profiles by the mechanism they rely on.

\textbf{Model exploitation as the core obstacle.} Sections 1.1 and 2 frame every method as a response to distributional shift. The survey argues that a learned model inherits the shift in both states and actions, and that, because no exploration is possible, a policy optimized against the model can drift into regions where the model is wrong but looks rewarding. Section 1.2 adds the related problem of compounding prediction error and mentions Lipschitz constraints on the model and adversarial distribution matching as remedies proposed in the general model-based RL literature.

\textbf{Uncertainty-penalized MDPs.} MOReL and MOPO (Sections 2.1 and 2.2) both convert ensemble uncertainty into a reward penalty. MOReL thresholds the largest pairwise disagreement among ensemble predictions and diverts the agent to an absorbing HALT state that carries a negative reward, while MOPO uses the largest Frobenius norm among the predicted Gaussian covariances. The survey notes that the resulting penalized reward can be passed directly to classical planners such as the linear-quadratic regulator or Monte Carlo tree search.

\textbf{Behavior regularization within model rollouts.} BREMEN (Section 2.3) leaves the model unpenalized and restrains the policy instead. An ensemble trained with a squared-error loss generates imagined rollouts, the policy starts as a Gaussian centered on a behavior-cloned policy, and each update is bounded by a KL trust region. The survey relays the authors' claims that this keeps divergence from the data in check implicitly, works better than explicit KL penalties, and yields good sample efficiency when data are scarce.

\textbf{Provable pessimism in value iteration.} Section 2.4 turns to theory. VI-LCB subtracts a coverage-based penalty inside the Bellman backup, which the survey describes as reversing the sign of an online exploration bonus, and it is linked to the earlier PEVI algorithm. The section then presents the Bernstein-style penalty of Li et al. and their result that model-based offline RL reaches minimax-optimal sample complexity in tabular MDPs without a burn-in cost, with bounds governed by a single-policy clipped concentrability coefficient.

\textbf{Conservative critics without explicit uncertainty.} COMBO (Section 2.5) is presented as the step that removes the need for an uncertainty estimate. Its critic objective lowers values under a chosen sampling distribution and raises them on dataset pairs, and the survey treats that sampling distribution as a hyperparameter that sets how conservative the method is. The survey highlights COMBO's guarantee of safe improvement over the behavior policy and treats the method as an answer to the unreliability of ensemble-based uncertainty.

\textbf{Adversarial models and relative pessimism.} The last two profiles turn the model itself into an adversary. RAMBO (Section 2.6) trains the dynamics to lower the policy's value while staying close to the maximum-likelihood model, holding the Lagrange multiplier constant as a tuned weight and swapping the total-variation constraint for a likelihood loss. ARMOR (Section 2.7) sets up a two-player game in which the policy maximizes its worst-case performance relative to a reference policy, an idea the survey traces to the relative pessimism of Cheng et al.

\question{6. What is the survey method?}

He selected the material personally and reviews it narratively, naming no databases, search terms, time window, or inclusion criteria and not explaining why these seven algorithms were chosen over others. The organization is linear. Section 1 defines reinforcement learning, offline RL, and model-based RL in words, without formal MDP preliminaries; Section 2 explains how the two fields combine and then gives one subsection per method, ordered roughly by publication date from MOReL (2020) to ARMOR (2023); Section 3 closes with two paragraphs on challenges.

Each method profile follows a similar pattern: a sentence on the core idea, the key objectives or update rules in the notation of whichever source the author followed, and a summary of the original authors' claims. The paper contains twenty numbered equations but no figures, tables, experiments, or benchmark comparisons, and it derives no new results. Because notation changes with the source, there is no unifying framework, and the learned dynamics are written differently from one section to the next. Several passages lean heavily on the survey of Figueiredo Prudencio et al. That survey is cited four times, including for the planner remark and the claim that uncertainty estimates are unreliable, but the MOReL and MOPO equations (Eqs. 1-4), which reproduce its Eqs. 31-34 almost symbol for symbol, carry no citation to it. The method is therefore descriptive summary rather than systematic synthesis.

\question{7. What are the survey findings?}

The survey's main conclusion is that distributional shift, made worse by model exploitation, is the defining problem of offline model-based RL and remains open. Section 3 states that, despite the range of remedies, none of the reviewed methods provides a provable performance guarantee that holds regardless of the data. The second conclusion concerns deployment. Drawing on the ARMOR authors, the survey argues that theoretically grounded methods often go unused because they cannot reliably outperform the policy already in service, and it calls for research on guaranteed improvement over that policy.

At the level of individual methods, the findings are mostly the original papers' own claims. Reward-penalty methods such as MOReL and MOPO are attractive because they reuse existing planners, but they depend on uncertainty estimates that the survey, citing Figueiredo Prudencio et al., regards as error-prone and difficult to compute. COMBO avoids such estimates through a conservative critic while keeping a safe-improvement guarantee, although RAMBO's authors contend that COMBO's pessimistic updates leave persistent local maxima in the Q-function that adversarial model training makes less likely. ARMOR is reported never to fall below its reference policy under a fixed hyperparameter setting and to compete with any policy the data cover when tuned well. On the theoretical side, the only optimality result the survey reports, Li et al.'s proof that a Bernstein-penalized variant of VI-LCB reaches minimax-optimal sample complexity without a burn-in cost, is confined to tabular MDPs, and none of the deep methods is reported to carry a comparable guarantee. The survey also suggests that models may exploit large historical datasets more fully than model-free methods because they are fit by supervised learning.

\question{8. What are the limitations of the survey?}

Coverage is thin. Seven methods and nineteen references leave out much of the subfield as it stood in 2023, including the study by Lu et al. (ICLR 2022) of design choices in ensemble-based uncertainty penalties, trajectory-level sequence models such as the Trajectory Transformer, and the image-based and hybrid predictive models that Levine et al. had already reviewed. With no benchmark results or performance comparison, readers cannot judge how the methods compare in practice. No search or selection procedure is stated, which makes the choice of papers impossible to audit.

Depth is also limited. Most profiles restate equations and pass on the authors' claims with little independent critique, and there is no comparative table showing how the methods relate. Some descriptions are loose. COMBO, for instance, is said to penalize rewards and to rely on adversarial training, whereas the objective the survey prints acts on Q-values, and its degree of caution is credited to the sampling distribution even though that objective carries an explicit weight. The prose contains grammatical slips, and the call for future work on improvement relative to a deployed policy is phrased ambiguously. The paper acknowledges none of these limits. Because it has not been revised since its first version, later model-based work such as MOBILE (ICML 2023), which measures uncertainty through inconsistency among Bellman estimates across the model ensemble, is not reflected.

\question{9. How was the survey assessed relative to other surveys?}

He positions the paper only loosely against earlier surveys. Its references include three reinforcement learning surveys (the tutorial of Levine et al., the taxonomy survey of Figueiredo Prudencio et al., and the model-based RL survey of Luo et al.) and a Monte Carlo tree search survey by Browne et al. cited only for the planner remark. All serve as background authorities for definitions and for claims such as the unreliability of uncertainty estimates; the paper never says what it adds beyond them.

Against the other two surveys in this set, He's review is the narrowest, shortest, and most recent. Levine et al. give offline model-based RL an analytically richer section of about three pages: it explains model exploitation through an explicit error bound from Janner et al., covers image-based and hybrid predictive models and model-based off-policy evaluation, and asks whether models can beat model-free dynamic programming. Figueiredo Prudencio et al. place model-based methods inside a full taxonomy, treat MOReL, MOPO, BREMEN, and COMBO in comparable mathematical detail, and add benchmark and performance reviews that He lacks. He's distinct additions are RAMBO, ARMOR, VI-LCB, and the sample-complexity result of Li et al., none of which appears in either broader survey. In formality, He prints more equations per method than Levine et al.'s model-based section but, unlike both other surveys, has no shared notation. It suits readers wanting a quick list of key model-based papers, while Levine et al. offer conceptual grounding and Figueiredo Prudencio et al. a structured overview with empirical context.

\question{10. What are the ethical implications? How are they addressed?}

Ethics, safety, and broader impact get no section in He's survey, and societal consequences go unmentioned. Only technical material comes near: COMBO's guarantee of safe improvement over the behavior policy, ARMOR's promise not to fall below a reference policy, and the closing observation that methods which cannot reliably beat the current policy are seldom deployed.

The field as He presents it carries ethical weight. Offline learning is framed as improving on a deployed system from that system's own data, and the abstract cites real-world potential without naming any domain. Model exploitation therefore means a policy can look excellent inside a learned simulator yet fail once it replaces that system, while biases in past decisions pass into both the model and the new policy. With off-policy evaluation omitted, the survey never explains how such policies could be validated before release, and several reported methods hinge on a hand-chosen conservatism weight with no advice on setting it offline.

As scholarship, the survey avoids self-citation and draws on several research groups, but it relays the original claims without scrutiny, and its MOReL and MOPO passage follows Figueiredo Prudencio et al. closely while crediting that survey there only for the planner remark. Ethics is touched only implicitly, through the deployment remarks, which is thin treatment for methods meant to replace systems in use.

\question{11. What meta-themes are suggested by this survey? (minimum 3)}

\textbf{Uncertainty quantification and calibration.} The model-based methods in this survey rest on uncertainty heuristics: MOReL thresholds ensemble disagreement and MOPO takes the largest predicted covariance norm. Beyond a passing remark that such estimates can be unreliable, the survey never examines whether either quantity tracks actual model error, and its account of COMBO shows the field sidestepping uncertainty rather than calibrating it. For BCA this points to an opening, since the same penalties could in principle be sized by a calibrated, distribution-free width rather than by an ensemble statistic whose link to error goes unchecked.

\textbf{Adaptive, state-dependent conservatism.} Most surveyed methods vary their caution across state-action pairs but set its overall strength globally. MOReL uses one disagreement threshold, COMBO scales its penalty by a single weight (\ensuremath{\beta} in Eq. 12) that the survey loosely recasts as a choice of sampling distribution, RAMBO hand-sets its Lagrange multiplier, and ARMOR needs the right hyperparameter. VI-LCB comes closest to data-driven sizing, since its penalty shrinks with visit counts, though its Bernstein form keeps a global constant. BCA pursues the same data-driven sizing in deep hosts, letting calibrated residual uncertainty set per-sample strength without visit counts.

\textbf{Theory versus practice.} The survey places Li et al.'s tabular minimax result for a Bernstein-penalized VI-LCB beside deep methods whose guarantees depend on assumptions or well-chosen hyperparameters. Its conclusion that no method guarantees performance independently of the data, and its remark that provable methods often go undeployed, expose the gap. The same tension runs through model-free offline RL, which is why a project like BCA values a finite-sample coverage guarantee whose assumptions, such as exchangeability or correct weights under shift, are stated openly and whose coverage can be checked on benchmarks.

\textbf{Model-based versus model-free offline RL.} Although the survey is about models, its methods repeatedly echo model-free ideas. COMBO's critic has the same push-down and push-up form as CQL, BREMEN uses behavior cloning and trust-region constraints much like policy-constraint methods, and ARMOR imports the relative pessimism of ATAC, a model-free actor-critic method by Cheng et al. The survey's claim that models use large datasets better through supervised learning is asserted rather than tested. For a project built on model-free hosts such as CQL and TD3+BC, this overlap suggests calibrated penalties could transfer across both families.

\clearpage

\bookmarked{\subsection}{subsection}{Literature Review 3: A Survey on Offline Reinforcement Learning: Taxonomy, Review, and Open Problems (Figueiredo Prudencio et al.)}

\question{1. What is the citation?}

\begin{hangref}Figueiredo Prudencio, R., Maximo, M. R. O. A., \& Colombini, E. L. (2024). A survey on offline reinforcement learning: Taxonomy, review, and open problems. \emph{IEEE Transactions on Neural Networks and Learning Systems}, \emph{35}(8), 10237--10257. \url{https://doi.org/10.1109/TNNLS.2023.3250269}\end{hangref}

\question{2. What is the survey scope?}

The subject is deep offline (batch) reinforcement learning, where every policy must come from a fixed log of transitions gathered by a behavior policy, with no further interaction. The survey works in the single-agent MDP setting and states early on that it treats observations as Markov rather than using the POMDP formulation. On the algorithmic side it spans direct and implicit policy constraints, importance sampling, policy and value regularization, uncertainty estimation, model-based methods, one-step methods, imitation learning, and trajectory optimization. Its timeline (Fig. 4) runs from BCQ in late 2018 to early 2022, Table II classifies 26 entries (25 distinct methods, since BRAC appears in two variants), and the bibliography holds 91 references. It also reviews off-policy evaluation, the D4RL, RL Unplugged and DOPE benchmarks, and three short application cases. It leaves most marginalized importance sampling to Levine et al. and does not review theoretical analyses, multiagent methods or recommender applications. The level is taxonomy plus formal review, aimed at researchers who already know deep RL.

\question{3. What is the survey question?}

The question driving the survey is: how can the fast-growing body of offline RL algorithms be placed in one framework that shows which design choices each method makes, how those choices fare on different kinds of data, and where research effort should go next?

The narrower questions map roughly onto the survey's sections. Section II asks why offline learning is harder than off-policy learning and formalizes distributional shift. Section III asks which components any offline algorithm is assembled from. Section IV asks how the main published methods instantiate those components in a common notation. Section V asks how a policy can be judged without deploying it. Section VI asks which properties a good offline dataset should have, how well current benchmarks supply them, and which classes perform best on each property. Section VIII asks which of the open problems listed by Levine et al. remain unsolved.

\question{4. What is the significance of the survey?}

The paper appeared on arXiv in March 2022, when the main broad review of the area was the 2020 tutorial by Levine et al. The authors note that offline RL publications had multiplied rapidly and that many new methods relied on strategies the tutorial never discussed; among those the survey reviews, TD3+BC, one-step RL, IQL, COMBO, the Decision Transformer, the Trajectory Transformer and RvS all appeared after it. The survey answers with an updated review in a single notation and a taxonomy that separates an algorithm's class from the loss terms added to it, so methods can be read as combinations of reusable parts. The eight dataset properties it compiles from the D4RL and RL Unplugged papers, together with the Fig. 5 heatmap, help practitioners match algorithm classes to their data, and its critique of both benchmarks speaks to benchmark designers. Its uptake shows in later surveys: He's 2023 model-based survey cites it several times and presents MOReL and MOPO in nearly identical notation. For the BCA project, it explains three of the four host algorithms (CQL, TD3+BC and IQL) side by side and names dynamic modulation of conservatism as an open problem.

\question{5. What are the themes?}

A two-level taxonomy organizes the survey. Class membership (is-a relations) records what an algorithm learns: a dynamics model, a trajectory distribution, a one-step or multistep actor-critic, or an imitation policy. Optional loss terms (has-a relations), such as policy constraints, regularizers, uncertainty penalties and importance weights, can be attached to any actor-critic. Evaluation, benchmarks and open problems then form a second layer of themes.

\textbf{Distributional shift as the central difficulty.} Section II-C formalizes the problem with the Bellman objective in Eq. 4: the critic is fit on actions from the behavior policy but queried on actions from the learned policy, and small policy errors then move the state distribution away from the data. Most classes in Section III are introduced as ways to prevent, penalize or sidestep this mismatch. The introduction also places importance sampling and trajectory optimization apart, as approaches that do not target the shift directly.

\textbf{Policy constraints and behavior regularization.} The survey separates direct constraints, which estimate the behavior policy (BCQ, BRAC), from implicit ones that rely only on dataset samples (AWR, AWAC, TD3+BC), and it contrasts matching the whole distribution with matching only its support (BEAR's MMD). Value regularization is covered through CQL's lower-bound penalty and policy regularization through SAC's entropy term. It warns that a unimodal fit to multimodal data can break direct methods and that constraints imposed everywhere can be needlessly pessimistic.

\textbf{Uncertainty estimation and model-based pessimism.} Fig. 3 shows the spread of an ensemble of Q-functions serving as an uncertainty signal, and Eq. 11 subtracts a weighted uncertainty term from the policy objective. Section IV-D reviews Agarwal et al.'s disjoint-partition ensembles and REM, noting that such ensembles lack diversity and understate uncertainty. On the model-based side, MOReL penalizes states where ensemble members disagree beyond a threshold, MOPO uses the largest norm of the predicted covariances, BREMEN stresses deployment efficiency, and COMBO transfers CQL into a learned model.

\textbf{Alternatives to iterative actor-critic learning.} Several classes avoid a multistep critic that is repeatedly queried at policy-proposed actions. Brandfonbrener et al.'s one-step methods evaluate the behavior policy once and improve once, succeeding with only weak regularization. IQL trains only on dataset actions, fitting an expectile value that tracks the best in-support actions, then extracts an advantage-weighted policy. Imitation methods turn policy extraction into supervised learning on filtered or outcome-conditioned data (BAIL, ABM+MPO, CRR, RvS), and the Trajectory and Decision Transformers model whole sequences, which the survey credits for strength on sparse rewards.

\textbf{Evaluating policies without deployment.} Section V presents model-based estimation, importance sampling and fitted Q evaluation as the main tools for judging a policy offline and ties them to hyperparameter selection and early stopping. It summarizes Voloshin et al.'s comparison of 33 estimators and Paine et al.'s finding that FQE-based rankings were enough to tune constraint methods such as CRR. Section VI-B then describes the DOPE benchmark and its absolute-error, regret@k and rank-correlation metrics, concluding that FQE and model-based estimators lead but that no estimator wins everywhere.

\textbf{Dataset properties and benchmark design.} Section VI-A1 compiles, from the D4RL and RL Unplugged papers, eight properties a useful offline dataset should exercise, among them narrow data, undirected multitask data, sparse rewards, suboptimal data, non-Markovian behavior and nonstationarity. Table III marks which benchmark environments satisfy each. The survey argues that RL Unplugged relies on online-trained behavior policies, that D4RL lacks nonstationary data, that stochastic dynamics and risky biases (hazards absent from the logs) are underrepresented, and that multiagent datasets are missing. Fig. 5 then scores methods and classes against these properties.

\question{6. What is the survey method?}

Figueiredo Prudencio et al. write as experts surveying the field, not as systematic reviewers: no databases, search strings, inclusion criteria or screening counts are given. The authors cover what they consider the most popular methods and concede that Table II is far from complete. The selection mixes foundational papers with recent preprints.

The organizing device is the taxonomy. Fig. 2 draws an offline algorithm as a pipeline from an optionally filtered dataset to a model, a trajectory distribution or a policy, and Fig. 3 and Table I list the modifications possible relative to a vanilla multistep actor-critic. Section III writes each modification as a loss term in a shared notation, with short derivations such as the KKT route from a KL-constrained objective to advantage-weighted likelihood, and Section IV restates every method in that notation, including the per-decision form of importance sampling. A timeline in Fig. 4 orders methods by preprint date. The benchmark review borrows its dataset properties from the D4RL and RL Unplugged papers, and Table III records which environments exhibit each.

The only quantitative analysis is secondary. Fig. 5 uses D4RL scores collected from the literature, including baselines such as SAC and BC, without naming each score's source. It rescales them per dataset between the lowest and highest score, averages them over the datasets having each property, and gives each class the maximum score of its members. RL Unplugged is omitted because few papers report it, raw numbers go to supplementary material on GitHub, and no new experiments are run.

\question{7. What are the survey findings?}

According to Figueiredo Prudencio et al., most new offline algorithms are different responses to distributional shift, and the main families involve trade-offs. Direct policy constraints depend on a good estimate of the behavior policy and can be overly pessimistic. Implicit constraints such as AWR, AWAC and TD3+BC avoid that estimate and suit online fine-tuning better, with AWR and AWAC beating BEAR and BRAC once training continues online. From BRAC it draws the lesson that hyperparameter search mattered more than the divergence chosen, and from TD3+BC that performance does not track algorithmic complexity. CQL is credited with outperforming constraint methods on AntMaze and Kitchen, tasks that require stitching. One-step methods work because they tolerate weak regularization, yet they trail multistep methods when data are mostly suboptimal or coverage is broad. Uncertainty quantification in model-based methods is described as difficult and often unreliable, and removing the need for it is named as COMBO's main advantage over MOReL and MOPO. Importance sampling, even with per-decision, weighted and doubly robust corrections, stays poorly conditioned because weights multiply across the horizon, which motivates marginalized estimators such as GenDICE.

Fig. 5 indicates that newer methods generally beat older ones, with TT and IQL among the strongest. The best classes are trajectory optimization and one-step methods paired with implicit constraints and value regularization, and planning helps most on sparse-reward and undirected data when guided by a critic trained with dynamic programming. No method was evaluated on nonstationary data.

The open problems are reliable off-policy evaluation and hyperparameter tuning, which still lack a satisfying answer; incremental RL for changing data and fine-tuning; safety-critical RL, including risk-averse CVaR objectives; and the use of large unlabeled datasets, where results by Yu et al. and Yarats et al. suggest that data choices matter as much as algorithms.

\question{8. What are the limitations of the survey?}

With no stated search or inclusion rules, coverage reflects the authors' judgment, and depth is uneven across classes. The uncertainty-estimation section reviews only Agarwal et al.'s ensembles and REM, omitting earlier ensemble methods such as EDAC (An et al., 2021) and PBRL (2022); marginalized importance sampling is deferred to Levine et al.; and nothing is said about the theory of pessimism under partial coverage or about calibrating uncertainty estimates. The introduction promises to discuss a reliable offline workflow and dynamically adjusted conservatism, but Section VIII never returns to conservatism and addresses the workflow only through remarks on validation and early stopping.

The quantitative comparison is fragile. Because Fig. 5 pools scores from many papers, it likely mixes tuning budgets and codebases; it also shows no variance, rescales against whichever algorithms share a dataset, and gives each class its best member's score. Class boundaries blur: IQL counts as one-step although the text says it runs iterative dynamic programming. Small errors appear too: CQL is expanded as constrained rather than conservative Q-learning, and a remark about AWAC is attributed to Jiang and Li instead of Nair et al.

The algorithmic review stops with methods released by early 2022 (a few references added during revision reach late 2022), so 2023 methods such as ReBRAC are absent. Finally, at least 21 of the 91 references are coauthored by Sergey Levine, tilting the review toward one group's methods and benchmarks, though the authors never cite themselves.

\question{9. How was the survey assessed relative to other surveys?}

The survey positions itself against one predecessor, the 2020 tutorial of Levine et al., which it credits with setting out the field's challenges, early methods, applications and open problems. It claims to add a taxonomy, a single notation, the many methods published since, a review of benchmark properties and a performance figure. It adapts Fig. 1 from Levine et al., defers marginalized importance sampling and broad application coverage to Levine et al. and Fu and Di, and refers readers to Arulkumaran et al. and Wang et al. for deep RL background.

Among these three offline RL surveys, Figueiredo Prudencio et al. give the broadest map of model-free methods up to early 2022 and are the only ones to map benchmark environments to dataset properties (Table III) and aggregate published scores (Fig. 5). Levine et al. discuss evaluation and D4RL only briefly (Section 6.1), but their tutorial is longer and more conceptual: it builds the problem up from importance sampling and linear value-function analysis, frames offline RL as counterfactual inference, and covers applications such as recommender systems, though it predates TD3+BC, IQL and the sequence-model methods. He's 2023 preprint is far narrower, with 19 references and seven model-based methods (MOReL, MOPO, BREMEN, VI-LCB, COMBO, RAMBO, ARMOR). It reaches theory-driven and adversarial work that this survey omits but offers no taxonomy or benchmark review, and borrows this survey's notation. BCA is best served by this survey as a single reference on the host algorithms, with Levine et al. as the better conceptual grounding.

\question{10. What are the ethical implications? How are they addressed?}

No section of the survey treats ethics or broader impact, and its safety remarks are scattered. It motivates offline RL by the expense and risk of online interaction, warns that testing a policy outside the data's support could be hazardous in driving or medical diagnosis, and wants off-policy evaluation to supply confidence before deployment. It asks for datasets with risky biases, such as driving logs without a single crash, to test whether algorithms still learn to avoid hazards missing from the data, and Section VIII names safety-critical RL as an open problem. The insulin-dosing case is framed around patient safety, the thermal-power case mainly around cost and time.

Other implications go unmentioned. Its biased-data property (NB) concerns only narrow coverage, yet logged clinical or human data also carry past decision makers' prejudices, which a policy kept near that behavior can reproduce. Privacy, consent, accountability for automated decisions and the misuse of dialogue agents are absent. Fig. 5 rests on published scores rather than rerun experiments, although the raw numbers are released.

As scholarship it discloses funding, is open access and avoids self-citation, though its reliance on one laboratory raises a question of balance. In the end the survey treats safety as the technical task of avoiding unsupported actions and, through CVaR objectives, hazardous in-distribution states, and goes no further into ethics.

\question{11. What meta-themes are suggested by this survey? (minimum 3)}

\textbf{Adaptive, state-dependent conservatism.} The cost of applying one level of caution everywhere is a recurring concern. The survey argues that constraining the policy where every action earns zero reward only burdens the function approximator, and presents uncertainty estimates as a way to relax constraints where the model is trustworthy. Yet each mechanism it describes is scaled by one global number, such as TD3+BC's trade-off coefficient \ensuremath{\alpha}, the penalty weight in Eq. 11 or COMBO's data ratio. It names dynamic modulation of conservatism as an open problem without developing it, the gap BCA targets.

\textbf{Uncertainty quantification and calibration.} Uncertainty appears in the survey mainly as ensemble variance or model disagreement used as a penalty (Eq. 11, MOReL, MOPO). The survey itself says these estimates are unreliable: disjoint-data ensembles lack diversity and understate uncertainty, and COMBO is presented as a way to avoid uncertainty estimates altogether. It never asks whether an uncertainty estimate is calibrated, that is, whether a stated width matches the observed error. BCA's conformal calibration of Bellman-residual scales responds directly to that missing question.

\textbf{Online tuning and offline model selection.} Hyperparameters run through the survey as an unresolved weakness. BRAC's results hinged more on tuning than on the divergence chosen, multistep methods peak at the weakest regularization that still keeps training stable, and most papers stop after a preset budget of updates and score the final policy online. Off-policy evaluation and hyperparameter selection lead its list of open problems. A method that derives constraint strength from data instead of a hand-tuned coefficient reduces, though it does not remove, this dependence on online tuning.

\textbf{Benchmark dependence (D4RL).} The survey's comparative verdicts inherit the shape of one benchmark. Because most methods it reviews report D4RL results and little else, Fig. 5 can rank methods only on the properties D4RL happens to exercise, and Table III marks every D4RL environment as deterministic. Conclusions such as the strength of trajectory and one-step methods may therefore not transfer to stochastic, nonstationary or risk-biased data. Gains from calibrated, per-sample conservatism measured only on D4RL would carry the same caveat.

\clearpage

\bookmarked{\section}{section}{Part B: Themes and Observations}

\bookmarked{\subsection}{subsection}{Collected Themes}

Themes were gathered from the research papers' Q11 themes and the surveys' Q5 themes and Q11 meta-themes, and labels or explanations describing the same idea were merged into sixteen consolidated themes. A work is counted for a theme when it lists that theme or its other answers clearly exhibit it, counts are out of fifteen works, and a meta-theme's prevalence is judged by the total number of theme occurrences it gathers.

{\small
\begin{longtable}{P{3.4cm} P{3.6cm} P{8.2cm}}
\toprule
\textbf{Theme} & \textbf{Works} & \textbf{Significance} \\
\midrule
\endfirsthead
\toprule
\textbf{Theme} & \textbf{Works} & \textbf{Significance} \\
\midrule
\endhead
Distribution shift and out-of-distribution value error & CQL, TD3+BC, IQL, ReBRAC, A-Crab, Quantile Q-Learning, UNIQ, BFN Planning, Conformal Policy Evaluation, Bellman Calibration, WBCP, Levine et al., He, Figueiredo Prudencio et al. (14/15) & Bootstrapped critics queried at actions or states absent from the data make errors that cannot be corrected offline and that a maximizing policy exploits, so nearly every work defines itself by where it intercepts this mismatch: in the critic (CQL), the actor (TD3+BC), actor and critic together (ReBRAC), by never querying unseen actions (IQL), inside a learned model (He), or by resampling or reweighting calibration data under policy or covariate shift (Conformal Policy Evaluation, WBCP). Because the damage shows up in Bellman targets, the Bellman residual is a natural place to read per-sample uncertainty. \\
\addlinespace
Implementation details and hyperparameter sensitivity & CQL, TD3+BC, IQL, ReBRAC, A-Crab, Unifloral, Quantile Q-Learning, UNIQ, BFN Planning, Conformal Policy Evaluation, WBCP, Figueiredo Prudencio et al. (12/15) & Results hinge on choices outside the stated idea, from layer normalization and the discount factor in ReBRAC to per-domain or per-dataset conservatism coefficients such as CQL's threshold, IQL's expectile, XQL's temperature and UNIQ's \ensuremath{\kappa}, and even to calibration settings such as the rollout length in Conformal Policy Evaluation and the HPD level in WBCP. This makes gains hard to attribute and is the main practical argument for deriving conservatism from data, although Quantile Q-Learning and UNIQ show that tuning tends to move to new knobs rather than vanish. \\
\addlinespace
Conservatism versus policy improvement trade-off & CQL, TD3+BC, IQL, ReBRAC, A-Crab, Quantile Q-Learning, UNIQ, Levine et al., He, Figueiredo Prudencio et al. (10/15) & Too much caution ties the policy to possibly poor behavior while too little lets value errors dominate, so each method exposes a dial, such as CQL's \ensuremath{\alpha}, IQL's \ensuremath{\tau}, ReBRAC's $\beta_{1}$ and $\beta_{2}$ or A-Crab's \ensuremath{\beta}, whose best setting changes across datasets or domains, and Levine et al. relay a lower bound showing that the error growth which forces tight constraints cannot be avoided in the worst case. Treating conservatism as a dose to be set per sample rather than per dataset is the natural next step. \\
\addlinespace
Behavior regularization and policy constraints & CQL, TD3+BC, IQL, ReBRAC, Unifloral, Quantile Q-Learning, Levine et al., He, Figueiredo Prudencio et al. (9/15) & Keeping the policy near logged actions, through an actor penalty (TD3+BC, ReBRAC), advantage-weighted extraction (IQL, Quantile Q-Learning), CQL's implicit divergence or a trust region inside model rollouts (BREMEN, in He), is the field's main working defense against extrapolation, and ReBRAC collapses without it. Its strength is almost always one scalar applied identically to every transition, which is precisely the quantity a calibrated per-sample width would modulate. \\
\addlinespace
Theory versus practice & CQL, IQL, A-Crab, Quantile Q-Learning, UNIQ, Conformal Policy Evaluation, Bellman Calibration, WBCP, He (9/15) & Guarantees are proved for tabular, linear, finite-class or asymptotic settings, while experiments use deep networks, heuristic coefficients and modified algorithms, so what gets tested is often a heuristic inspired by a theorem, or a regime outside it, rather than the theorem itself, as with A-Crab's box-weight variant, CQL's thresholded \ensuremath{\alpha} and WBCP's small-$n_{\mathrm{eff}}$ regime. Any calibrated method inherits this gap and should state which assumptions, such as exchangeability, survive inside RL training and then check its coverage empirically. \\
\addlinespace
Adaptive, state-dependent conservatism & CQL, A-Crab, Quantile Q-Learning, UNIQ, WBCP, Levine et al., He, Figueiredo Prudencio et al. (8/15) & Levine et al. and Figueiredo Prudencio et al. name self-adjusting caution as an open problem, the methods He reviews still fix their overall strength globally, and CQL's own theory implies a penalty that shrinks with state-action counts, yet nearly every implementation applies one coefficient to all data. Quantile Q-Learning's learned temperature, UNIQ's conformal expectile, A-Crab's importance weights and WBCP's local thresholds are concrete attempts, but each keeps a global knob or collapses toward one, which is the gap calibrated per-sample conservatism must close. \\
\addlinespace
Uncertainty quantification and calibration & Quantile Q-Learning, UNIQ, Conformal Policy Evaluation, Bellman Calibration, WBCP, Levine et al., He, Figueiredo Prudencio et al. (8/15) & The surveys present uncertainty as a principled basis for pessimism but judge ensemble and model-disagreement estimates hard to calibrate and often unreliable, whereas the conformal and calibration papers test it directly: Conformal Policy Evaluation shows that raw distributional quantiles undercover, Bellman Calibration repairs a critic's scale, and WBCP reports how reliable a threshold itself is. Whether a stated width matches observed error is what makes uncertainty safe to use as a regularization weight. \\
\addlinespace
Online tuning and offline model selection & CQL, TD3+BC, ReBRAC, Unifloral, Quantile Q-Learning, UNIQ, Conformal Policy Evaluation, Figueiredo Prudencio et al. (8/15) & Most reported scores rely on hyperparameters or checkpoints chosen by running policies in the environment, a budget a truly offline user lacks; Unifloral prices that budget in evaluation episodes, and ReBRAC's Expected Online Performance shows that it needs about ten evaluated policies before its expected score reaches the ensemble-free state of the art. Deriving conservatism from data, as Quantile Q-Learning and UNIQ attempt, is one of the few direct responses, and Unifloral's protocol can test whether such a method helps at small budgets. \\
\addlinespace
Dataset composition and benchmark dependence & CQL, TD3+BC, IQL, Unifloral, UNIQ, BFN Planning, Levine et al., Figueiredo Prudencio et al. (8/15) & How the data were collected decides which family succeeds: mixed-quality and stitching-heavy data reward methods that can depart from or recombine logged behavior, while a uniform pull toward the data fails, as TD3+BC's 28.8 against CQL's 100.3 on the Walker2d random-expert mixture shows. Because most offline RL evidence comes from D4RL, whose environments Figueiredo Prudencio et al. mark as deterministic and which lacks nonstationary data, conclusions about which form of conservatism works best remain conditional on that benchmark. \\
\addlinespace
Evaluation rigor, reproducibility and fair comparison & TD3+BC, IQL, ReBRAC, Unifloral, UNIQ, BFN Planning (6/15) & Whether a claimed gain belongs to the algorithm depends on rerunning baselines under one protocol, reporting final rather than best checkpoints, separating tuning from reporting seeds and releasing code, practices that tighten from TD3+BC through ReBRAC to Unifloral. UNIQ, the closest prior work to BCA, shows the cost of neglecting them, since copied baselines, three seeds and single-seed ablations leave its 3.6-point gain over IQL weakly established. \\
\addlinespace
Minimalism and computational efficiency & TD3+BC, IQL, ReBRAC, Unifloral, UNIQ, Bellman Calibration (6/15) & Small, cheap changes make gains attributable, shrink the number of choices that need validation and make large tuning sweeps and budget-aware protocols affordable, which is why TD3+BC and its descendant ReBRAC became the baselines new methods must beat. UNIQ and Bellman Calibration carry the idea to uncertainty, suggesting that a lightweight or post-hoc calibration layer can add reliability without large ensembles or retraining. \\
\addlinespace
Importance weighting and off-policy evaluation & A-Crab, Conformal Policy Evaluation, Bellman Calibration, WBCP, Levine et al., Figueiredo Prudencio et al. (6/15) & Density and action-probability ratios are the standard correction for the mismatch between logged and target distributions, appearing in off-policy estimators, A-Crab's weighted Bellman penalty and weighted conformal calibration, but their variance grows with the horizon and with weak overlap, and weight error enters coverage bounds directly. BCA's uniform WBCP weights sidestep ratio estimation at the cost of calibrating to the dataset distribution rather than the learned policy's. \\
\addlinespace
Model-based and generative alternatives to actor-critic learning & Unifloral, BFN Planning, Levine et al., He, Figueiredo Prudencio et al. (5/15) & Learned dynamics models and trajectory generators relocate the shift problem rather than remove it, since a policy optimized inside a model seeks out regions where the model is wrongly optimistic, and in Unifloral none of the three prior model-based methods ranks above sixth of ten on any non-locomotion dataset. Their penalties echo model-free ones, as COMBO mirrors CQL, so calibrated widths could transfer to them, whereas critic-free planners such as BFN Planning offer no Bellman residual to calibrate. \\
\addlinespace
Distribution-free (conformal) guarantees & UNIQ, Conformal Policy Evaluation, WBCP (3/15) & Conformal calibration offers coverage that does not depend on a correctly specified model, something no ensemble or parametric scale in the set provides, but all three works show the guarantee weakening once exchangeability is relaxed or calibration moves inside RL: it is marginal rather than per-state, it rests on exchangeability (or, in WBCP, on correct weights and an assumed model) that temporal dependence and policy shift strain, and it covers only the quantity actually calibrated. These caveats define what a calibrated per-sample width can and cannot certify. \\
\addlinespace
Finite-sample guarantees under partial data coverage & A-Crab, Bellman Calibration, He (3/15) & A-Crab and the VI-LCB results in He show that strong guarantees require only that the data cover the comparator policy, measured by single-policy or $\ell_2$ concentrability, and A-Crab shows that the form of the Bellman penalty can decide whether the optimal rate is reached. Its importance-weighted residual penalty is in effect an adversarial per-sample reweighting of Bellman residuals, and Bellman Calibration repairs a critic's scale without completeness assumptions, both of direct relevance to residual-based calibration. \\
\addlinespace
Offline-to-online fine-tuning & IQL, ReBRAC, Figueiredo Prudencio et al. (3/15) & Once interaction resumes, the conservatism that protected offline learning can block adaptation: IQL's actor-independent critic keeps improving while CQL sometimes collapsed, and ReBRAC's hand-designed relaxation of its penalties helps unevenly. The right strength of conservatism therefore changes as data accumulate, which a fixed coefficient cannot express. \\
\bottomrule
\end{longtable}}

\bookmarked{\subsection}{subsection}{Meta-Themes}

{\small
\begin{longtable}{P{3.4cm} P{6.6cm} P{5.2cm}}
\toprule
\textbf{Meta-theme} & \textbf{Description and significance} & \textbf{Themes included} \\
\midrule
\endfirsthead
\toprule
\textbf{Meta-theme} & \textbf{Description and significance} & \textbf{Themes included} \\
\midrule
\endhead
Pessimism toward unsupported actions & This algorithmic core covers the extrapolation problem, the constraints and penalties that contain it, how strongly to apply them, how the problem resurfaces inside learned models and how the constraint must loosen once interaction returns. It reaches all fifteen works and gathers by far the most theme occurrences, establishing the strength of conservatism as the field's central design decision. & Distribution shift and out-of-distribution value error; Behavior regularization and policy constraints; Conservatism versus policy improvement trade-off; Model-based and generative alternatives to actor-critic learning; Offline-to-online fine-tuning \\
\addlinespace
Calibrated, data-dependent caution & These themes supply the statistical means to decide how far to trust each value estimate, from ensemble and quantile uncertainty to importance weights and conformal coverage. They are concentrated in the newest papers and in the surveys' critiques and open problems, and they mark the space in which calibrated per-sample conservatism is positioned. & Adaptive, state-dependent conservatism; Uncertainty quantification and calibration; Importance weighting and off-policy evaluation; Distribution-free (conformal) guarantees \\
\addlinespace
Tuning, engineering and the offline premise & Reported progress often rests on engineering choices and on coefficients chosen through online evaluation, which undercuts the promise of learning without interaction, while simple, cheap methods and budget-aware protocols expose and price that dependence. Any method that derives its conservatism from data is answering this meta-theme. & Implementation details and hyperparameter sensitivity; Online tuning and offline model selection; Minimalism and computational efficiency \\
\addlinespace
Credibility and reach of evidence & These themes ask whether proofs and benchmark scores hold where offline RL would actually be used, given idealized assumptions, uneven comparison protocols and heavy reliance on simulated D4RL tasks. They set the standard a new method must meet: stating its assumptions openly and checking them empirically. & Theory versus practice; Dataset composition and benchmark dependence; Evaluation rigor, reproducibility and fair comparison; Finite-sample guarantees under partial data coverage \\
\bottomrule
\end{longtable}}

\bookmarked{\subsection}{subsection}{Summary of Observations}

Sixteen merged themes form four meta-themes. Pessimism toward unsupported actions is the most prevalent: distribution shift appears in fourteen works, the caution-improvement trade-off in ten and behavior constraints in nine, since every offline learner must decide how far to trust values its data cannot verify. Hyperparameter sensitivity, in twelve works, anchors a meta-theme on tuning, because scores tuned through online evaluation overstate what a truly offline user can achieve. Calibrated uncertainty, which the surveys deem principled but unreliable, and the credibility of theory and benchmarks, which limits how far results transfer, complete the set. Three themes found in only three works each carry outsized weight: conformal guarantees give the only uncertainty estimates with stated coverage, partial-coverage theory ties guarantees to how well data cover a policy, and fine-tuning suggests that fixed conservatism impedes adaptation, together motivating the calibrated, per-sample conservatism my project pursues.

{\small\itshape\color{gray}(143 words)}

\bookmarked{\section}{section}{References}

\begin{hangref}Figueiredo Prudencio, R., Maximo, M. R. O. A., \& Colombini, E. L. (2024). A survey on offline reinforcement learning: Taxonomy, review, and open problems. \emph{IEEE Transactions on Neural Networks and Learning Systems}, \emph{35}(8), 10237--10257. \url{https://doi.org/10.1109/TNNLS.2023.3250269}\end{hangref}
\begin{hangref}Fujimoto, S., \& Gu, S. S. (2021). A minimalist approach to offline reinforcement learning. \emph{Advances in Neural Information Processing Systems}, \emph{34}, 20132--20145. \url{https://proceedings.neurips.cc/paper/2021/hash/a8166da05c5a094f7dc03724b41886e5-Abstract.html}\end{hangref}
\begin{hangref}Gan, F., Lu, Y., Zhang, Y., \& Liu, Y. (2025). Conformal prediction beyond the horizon: Distribution-free inference for policy evaluation. \emph{Advances in Neural Information Processing Systems}, \emph{38}. \url{https://neurips.cc/virtual/2025/poster/118031}\end{hangref}
\begin{hangref}Gao, X., Li, S., Cai, Y., \& Yu, W. (2026). Quantile Q-learning: Revisiting offline extreme Q-learning with quantile regression. \emph{Transactions on Machine Learning Research}. \url{https://openreview.net/forum?id=tBKznsUimN}\end{hangref}
\begin{hangref}He, H. (2023). \emph{A survey on offline model-based reinforcement learning} (arXiv:2305.03360v1) [Preprint]. arXiv. \url{https://doi.org/10.48550/arXiv.2305.03360}\end{hangref}
\begin{hangref}Jackson, M. T., Berdica, U., Liesen, J., Whiteson, S., \& Foerster, J. N. (2025). A clean slate for offline reinforcement learning. \emph{Advances in Neural Information Processing Systems}, \emph{38}. \url{https://neurips.cc/virtual/2025/oral/119623}\end{hangref}
\begin{hangref}Killingberg, L., \& Langseth, H. (2026). \emph{Bayesian flow networks for offline trajectory planning} (arXiv:2608.25163v1) [Preprint]. arXiv. \url{https://doi.org/10.48550/arXiv.2608.25163}\end{hangref}
\begin{hangref}Kostrikov, I., Nair, A., \& Levine, S. (2022). Offline reinforcement learning with implicit Q-learning. In \emph{International Conference on Learning Representations (ICLR 2022)}. \url{https://openreview.net/forum?id=68n2s9ZJWF8}\end{hangref}
\begin{hangref}Kumar, A., Zhou, A., Tucker, G., \& Levine, S. (2020). Conservative Q-learning for offline reinforcement learning. \emph{Advances in Neural Information Processing Systems}, \emph{33}, 1179--1191. \url{https://proceedings.neurips.cc/paper/2020/hash/0d2b2061826a5df3221116a5085a6052-Abstract.html}\end{hangref}
\begin{hangref}Levine, S., Kumar, A., Tucker, G., \& Fu, J. (2020). \emph{Offline reinforcement learning: Tutorial, review, and perspectives on open problems} (arXiv:2005.01643v3) [Preprint]. arXiv. \url{https://doi.org/10.48550/arXiv.2005.01643}\end{hangref}
\begin{hangref}Lou, X., \& Luo, P. (2026). \emph{Weighted Bayesian conformal prediction} (arXiv:2604.06464v2) [Preprint]. arXiv. \url{https://doi.org/10.48550/arXiv.2604.06464}\end{hangref}
\begin{hangref}Tarasov, D., Kurenkov, V., Nikulin, A., \& Kolesnikov, S. (2023). Revisiting the minimalist approach to offline reinforcement learning. \emph{Advances in Neural Information Processing Systems}, \emph{36}, 11592--11620. \url{https://proceedings.neurips.cc/paper_files/paper/2023/hash/26cce1e512793f2072fd27c391e04652-Abstract-Conference.html}\end{hangref}
\begin{hangref}Upadhyay, A. (2026, July 11). \emph{UNIQ: Conformal calibration for adaptive conservatism in offline reinforcement learning} [Poster presentation]. ICML 2026 Workshop on Decision-Making from Offline Datasets to Online Adaptation: Black-Box Optimization to Reinforcement Learning, Seoul, South Korea. \url{https://icml.cc/virtual/2026/83644}\end{hangref}
\begin{hangref}van der Laan, L., \& Kallus, N. (2026). \emph{Bellman calibration for V-learning in offline reinforcement learning} (arXiv:2512.23694v2) [Preprint]. arXiv. \url{https://doi.org/10.48550/arXiv.2512.23694}\end{hangref}
\begin{hangref}Zhu, H., Rashidinejad, P., \& Jiao, J. (2023). Importance weighted actor-critic for optimal conservative offline reinforcement learning. \emph{Advances in Neural Information Processing Systems}, \emph{36}, 49579--49602. \url{https://proceedings.neurips.cc/paper_files/paper/2023/hash/9bb93a3c1a424654aaea6f5b594e94d5-Abstract-Conference.html}\end{hangref}

\end{document}
